\documentclass[11pt,letterpaper]{article}
\usepackage[T1]{fontenc}
\usepackage{lmodern,amsmath,amssymb,amsthm,mathtools,booktabs,array,longtable}
\usepackage[margin=0.82in]{geometry}
\usepackage{microtype,fancyhdr,xurl,enumitem}
\usepackage[colorlinks=true,linkcolor=black,urlcolor=black,citecolor=black]{hyperref}
\hypersetup{pdftitle={Subset sums representable by four summands},pdfsubject={Consolidated proof of 71 <= N <= 85}}
\setlength{\parindent}{0pt}\setlength{\parskip}{4pt}
\setlength{\emergencystretch}{3em}
\setlength{\headheight}{14pt}
\setlist{nosep,topsep=4pt}
\pagestyle{fancy}\fancyhf{}\lhead{Subset sums and four summands}\rhead{Consolidated proof}\cfoot{\thepage}
\newcommand{\FS}{\operatorname{FS}}
\newcommand{\Z}{\mathbb Z}\newcommand{\R}{\mathbb R}
\newcommand{\card}[1]{\lvert#1\rvert}
\newcommand{\iv}[2]{[#1,#2]_{\Z}}
\newcommand{\own}{\operatorname{own}}
\newcommand{\calA}{\mathcal A}\newcommand{\calQ}{\mathcal Q}
\newcommand{\LOW}{\mathrm{LOW}}\newcommand{\HIGH}{\mathrm{HIGH}}
\newcommand{\modref}[1]{\texttt{\nolinkurl{#1}}}
\newtheorem{theorem}{Theorem}[section]
\newtheorem{lemma}[theorem]{Lemma}
\newtheorem{proposition}[theorem]{Proposition}
\newtheorem{corollary}[theorem]{Corollary}
\theoremstyle{definition}\newtheorem{definition}[theorem]{Definition}
\newtheorem{remark}[theorem]{Remark}
\begin{document}
\begin{center}
{\LARGE\bfseries Subset sums representable\par by four summands}\\[6pt]
{\large A consolidated proof of the bounds $71\le N\le85$}\\[5pt]
5 October 2026
\end{center}
\begin{abstract}
Let $N$ be the largest cardinality of a finite real set whose every subset sum is a sum of four of its elements, with repetition allowed in the four-term sum. We give an explicit $71$-element construction and organize the exclusion of cardinality $86$ into one coherent argument, obtaining $71\le N\le85$ from the previously established bound $N\le86$. The common covering budget, source estimates and loss-counting principles are stated once. The remaining cases are divided by their actual hole counts and their numbers of empty covering units. Each case is reduced to explicit endpoint, overlap or scalar contradictions. The prior upper bound, structural theorem and stated general and finite-backed additive-combinatorial results are retained inputs; this report does not claim a proof from first principles of those inputs. The companion program checks the explicitly identified finite assertions using exact arithmetic. The main text gives the common argument and its verification boundary; the longer complete stratum proofs are collected in labeled appendices. The exact value of $N$ remains unknown.
\end{abstract}
\setcounter{tocdepth}{2}{\setlength{\parskip}{0pt}\tableofcontents}
\clearpage
\section{The problem and the lower bound}\label{sec:problem}
For a finite set $A\subset\R$, define
\[
\FS(A)=\left\{\sum_{x\in B}x:B\subseteq A\right\},\qquad
4A=\{a_1+a_2+a_3+a_4:a_i\in A\}.
\]
The empty sum is included in $\FS(A)$. A coefficient is used at most once in a subset sum, but repetitions are allowed in $4A$. Call $A$ \emph{feasible} if $\FS(A)\subseteq4A$ and put
\[
N=\max\{\card A:A\subset\R\text{ is finite and feasible}\}.
\]
The previously established theorem $N\le86$ is an input throughout this report. Its proof is not replaced by the exclusion below. In particular, feasibility is not assumed to pass to arbitrary subsets of a feasible set.

\begin{theorem}[Main result]\label{thm:main}
Every feasible finite real set has at most $85$ elements, and there is a feasible set with $71$ elements. Thus $71\le N\le85$.
\end{theorem}
The upper bound will follow by contradiction. A hypothetical counterexample must have exactly $86$ elements by the retained upper bound. The common structural theorem gives an integral $78$-point source and an integral $70$-point core. We exclude every resulting source by splitting according to its core hole count $z$ and the number $j$ of genuinely empty ordinary covering units. The families are $Z_j$ for $z=0$, $P_j$ for $1\le z\le14$, and $H$ for $z\ge15$. Only $0\le j\le5$ occurs in the first two families. In $H$, a stronger defect bound will give $j\le1$.

\subsection{An explicit set with seventy one elements}\label{sec:lower}
Write $[u,v]_{\Z}=[u,v]\cap\Z$ and set
\[
I_0=[-31,31]_{\Z},\qquad
A_0=I_0\cup\{\pm93,\pm155,\pm217,\pm713\}.
\]
This set has $63+8=71$ elements. Its positive mass is
\[
1+\cdots+31+93+155+217+713=496+1178=1674,
\]
so $\FS(A_0)\subseteq[-1674,1674]_{\Z}$. The following table proves that this entire interval is contained in $4A_0$. Every expression in its left column uses four summand positions; zeros in $I_0$ may fill unused positions.
\begin{center}\small
\begin{tabular}{@{}lr@{}}\toprule
Subset of $4A_0$&Interval covered\\\midrule
$4I_0$&$[-124,124]_{\Z}$\\
$217+3I_0$&$[124,310]_{\Z}$\\
$217+93+2I_0$&$[248,372]_{\Z}$\\
$217+155+2I_0$&$[310,434]_{\Z}$\\
$2\cdot217+2I_0$&$[372,496]_{\Z}$\\
$2\cdot217+\{93,155,217\}+I_0$&$[496,682]_{\Z}$\\
$713+3I_0$&$[620,806]_{\Z}$\\
$713+\{93,155,217\}+2I_0$&$[744,992]_{\Z}$\\
$713+217+\{93,155,217\}+I_0$&$[992,1178]_{\Z}$\\
$2\cdot713-\{217,155,93\}+I_0$&$[1178,1364]_{\Z}$\\
$2\cdot713+2I_0$&$[1364,1488]_{\Z}$\\
$2\cdot713+\{93,155,217\}+I_0$&$[1488,1674]_{\Z}$\\\bottomrule
\end{tabular}
\end{center}
Within a row involving a three-element set, the component intervals overlap or are adjacent. Across rows, every integer from $0$ through $1674$ is covered. Since $A_0=-A_0$, reflection covers the negative half as well. Hence $A_0$ is feasible and $N\ge71$.

\subsection{How to read the upper bound}
The source and capacity lemmas in the next sections are reused without changing their deficits or their meanings of activity. A target hit always means membership in the actual subset-sum target, not merely membership in its integer hull. A paired unit has one capacity; choosing an active component does not create a second capacity. Physical signs are never reordered after an owner has been selected. These conventions make the case divisions statements about actual configurations rather than a catalogue of formal descriptors.

The proof has three levels. Elementary source, interval and loss arguments are proved here. The structural theorem and named general additive-combinatorial inputs are stated at their precise boundary. Exact finite lemmas are accompanied by a program and an explicit statement of what it checks. The final section distinguishes these levels so that a successful program run cannot be mistaken for a proof of an imported infinite theorem.

\subsection{A dependency roadmap}
The logical backbone is short even though some endpoint cases require several pages.
\begin{enumerate}
\item \textbf{Start at the retained upper bound.} The established $N\le86$ theorem makes a counterexample to $N\le85$ have exactly $86$ elements. Rationalization preserves this cardinality and a prescribed zero.
\item \textbf{Pass to one common source.} Covering, peeling and the structural theorem give $D,D_1,E$ of sizes $78,74,70$, one original residual deficit $\delta=c-78j$, and the physical side sizes summing to $69$.
\item \textbf{Prove the reusable estimates once.} Actual-resource loss counting, rooted forests, raw gaps, endpoint ranks, source envelopes and filled-interval propagation are the common tools. Later cards state which real resources supply their inequalities.
\item \textbf{Separate the cases by actual holes.} For $z=0$ and $1\le z\le14$, the scalar budget permits only $j=0,\ldots,5$. The zero-hole and positive-hole appendices retain their different target memberships throughout.
\item \textbf{Use dense information only where needed.} At $z\ge15$, the finite-backed defect input forces $c\le154$ and $j\le1$. Common-strip localization then reduces the last one-unit case to the same first-gap and rooted terminal argument as P1.
\item \textbf{Assemble the disjoint alternatives.} Every actual configuration is in an excluded row. The explicit construction supplies the lower bound. The program checks specified finite supporting assertions, while the stated general and infinite implications remain the written proof boundary.
\end{enumerate}
The main reading route ends with the completion and verification sections. The complete longer proofs of Z2--Z5 and P2--P4 then appear in labeled appendices, with live cross-references from the main argument. Read those appendices to check every classification and endpoint card; no case proof has been omitted.


% Consolidated common mathematics. No document preamble.
% Requires amsmath, amssymb, amsthm, booktabs, longtable, array.
% Macros: \FS, \card{...}, \Z, \R, \iv{a}{b} (integer interval).
% Environments: theorem, lemma, proposition, proof.

\section{A common source and its covering budget}\label{sec:common-foundations}

For a finite set $X$, write $\FS(X)$ for all its subset sums, including
the empty sum, and $kX$ for sums of exactly $k$ elements with repetition
allowed. A set $A$ is feasible when $\FS(A)\subseteq4A$. Put
\[
R(X)=\{x+y:x,y\in X,\ x\ne y\},\qquad
\Delta(X)=\card{2X\setminus R(X)},\qquad T_r=\frac{r(r+1)}2.
\]
Here $R(X)$ is a set, whereas a later unadorned $R$ denotes a sum of two
core maxima. A hull-hole count is the number of missing integers between
the minimum and maximum of a finite integer set.

\begin{lemma}[Reduction at exact cardinality]\label{lem:common-rationalization}
If a real feasible set of cardinality $n$ exists, an integer feasible set
of the same cardinality exists. A prescribed zero coordinate can be
preserved. Given the previously established upper bound $N\le86$, a
counterexample to $N\le85$ would therefore give an integer feasible set
of cardinality $86$ containing zero.
\end{lemma}
\begin{proof}
Order the real elements $x_1<\cdots<x_n$. Choose one four-term
representation for each of their finitely many subset sums. Each chosen
identity is a homogeneous linear equation with integer coefficients in
the $x_i$. Add $x_i=0$ when zero is prescribed. The solution space has a
basis of rational vectors, so its rational points are dense. Its strict
ordering region is relatively open and nonempty; it contains a rational
point. Clearing denominators preserves all the chosen identities and
strict inequalities.

The earlier upper bound makes the cardinality of a counterexample
exactly $86$. If zero were absent, adjoining it would preserve
$\FS(A)$ and could only enlarge $4A$, giving a feasible set of size $87$.
This contradicts the earlier theorem. Apply the rationalization just
proved with that zero prescribed. We have used an existence reduction,
not the assertion that every real feasible set itself lies on a rational
lattice. Nor have we assumed that feasibility passes to arbitrary subsets.
\end{proof}

\subsection{The four-extreme cover}
Remove the four largest positive and four smallest negative elements of
$A$, taking every element on a sign with fewer than four. Call their
union $B$, put $D=A\setminus B$, and let $d=\card D$.

\begin{lemma}[Four-extreme covering]\label{lem:common-cover}
For a feasible zero-containing set,
\begin{equation}\label{eq:common-cover}
\card{\FS(D)}\le2\card{2D}+18d+85.
\end{equation}
Before the two signs are paired, each positive target has one translate
of $2D$, ten translates of $D$, and at most $55$ fixed labels.
\end{lemma}
\begin{proof}
On the positive sign write $p_1>p_2>p_3>p_4>0$ and $T=p_1+p_2+p_3+p_4$.
For $x\in\FS(D)$, $x>0$, feasibility represents $T+x$ as
$z_1+z_2+z_3+z_4$ with $z_1\ge\cdots\ge z_4$ in $A$.
The first failure of $z_i\le p_i$ is in position two, three, or four.
The possible forms are
\[
\begin{array}{c|l}
2&2p_1+u+v,\\
3&p_1+2p_2+u,\quad3p_2+u,\\
4&p_1+p_2+2p_3,\quad p_1+3p_3,\quad2p_2+2p_3,
\quad p_2+3p_3,\quad4p_3.
\end{array}
\]
Subtract $T$ and split $A=D\sqcup B$. The first line gives one translate
of $2D$, eight translates of $D$, and $36$ fixed pair labels. The second
gives two more translates of $D$ and $7+7$ fixed labels, because its free
entry is not $p_1$. The last gives five fixed labels. This proves the
single-target $55$ assertion. Reflection gives the other sign.

Suppose both signs have four elements, and write the negative magnitudes
$q_1>q_2>q_3>q_4>0$. Define
\[
\begin{aligned}
\lambda_p&=p_1-p_2-p_3-p_4,&\alpha_p&=p_2-p_3-p_4,
&\beta_p&=2p_2-p_1-p_3-p_4,
\end{aligned}
\]
and similarly for $q$. For $u\le v$,
\[
\card{(u+D)\cap(0,\infty)}+
\card{(v+D)\cap(-\infty,0)}\le d,
\]
since no source point meets both strict half-line inequalities.
Always $\lambda_p-q_1<-\lambda_q+p_1$. If
$\beta_p+\beta_q\le0$, pair $\beta_p$ with $-\beta_q$.
Otherwise pair $\lambda_p-q_2$ with $-\lambda_q+p_2$, using
\[
(\lambda_p-q_2)-(-\lambda_q+p_2)
=-(\beta_p+\beta_q)-2(p_3+p_4+q_3+q_4)<0.
\]
The two distinct pairs replace four capacities $d$ by two, leaving
$18d$. The fixed-label calculation below removes at least $26$ of the
$110$ labels, leaving $84$. Adding the zero subset sum proves
\eqref{eq:common-cover}. If one original sign has fewer than four
elements, $D$ has no subset sums on that sign; its one remaining table
already satisfies the asserted bound.
\end{proof}

Here is the fixed-label calculation in full. Set
\[
B^+=(p_1,p_2,p_3,p_4,-q_1,-q_2,-q_3,-q_4),\quad
B^-=(q_1,q_2,q_3,q_4,-p_1,-p_2,-p_3,-p_4).
\]
For $\varepsilon\in\{+,-\}$, define
\[
C^\varepsilon_{ij}=\lambda_\varepsilon+B_i^\varepsilon+B_j^\varepsilon,
\quad U_i^\varepsilon=\alpha_\varepsilon+B_i^\varepsilon,
\quad V_i^\varepsilon=\beta_\varepsilon+B_i^\varepsilon,
\]
where $1\le i\le j\le8$ for $C$ and $2\le i\le8$ for $U,V$.
For a decreasing quadruple $x$, put
\[
\begin{array}{lll}
F_1=x_3-x_4,&F_2=2x_3-x_2-x_4,&F_3=x_2+x_3-x_1-x_4,\\
F_4=2x_3-x_1-x_4,&F_5=3x_3-x_1-x_2-x_4.
\end{array}
\]
Thus there are $2(36+7+7+5)=110$ labels. Each label is usable only if
its score is positive on the corresponding reflected target. Let
$r=(p_1-p_2,p_2-p_3,p_3-p_4,p_4)$ and define $s$ from $q$.
All eight coordinates are positive. The following table gives the
coefficient vector of each indicated sum in $(r;s)$. A starred row is
also used after interchanging signs.
\begin{center}\small
\begin{longtable}{c l l}
\toprule
&\text{sum of labels}&\text{coefficient vector}\\\midrule
1&$C^+_{25}+C^-_{25}$&$(0,0,-1,-2;0,0,-1,-2)$\\
$2^*$&$C^+_{48}+F^+_3$&$(0,0,0,-1;0,0,0,-1)$\\
3&$C^+_{58}+C^-_{58}$&$(0,-1,-2,-4;0,-1,-2,-4)$\\
$4^*$&$C^+_{77}+F^+_4$&$(0,-1,0,-2;0,0,-2,-2)$\\
5&$U^+_6+U^-_6$&$(0,0,-1,-2;0,0,-1,-2)$\\
6&$V^+_6+V^-_6$&$(-1,0,-1,-2;-1,0,-1,-2)$\\
$7^*$&$C^+_{15}+F^+_1+C^-_{55}$&$(0,-1,-1,-3;0,-1,-2,-3)$\\
$8^*$&$C^+_{26}+C^-_{57}+V^-_4$&$(0,0,-2,-3;0,0,-2,-3)$\\
$9^*$&$C^+_{35}+V^+_3+C^-_{46}$&$(0,0,0,-2;0,-1,-2,-2)$\\
$10^*$&$C^+_{36}+U^+_4+U^-_5$&$(0,0,-1,-2;0,0,-1,-2)$\\
$11^*$&$C^+_{44}+U^+_8+F^+_5$&$(0,-1,0,-1;0,0,0,-1)$\\
$12^*$&$C^+_{45}+V^+_8+C^-_{67}$&$(0,0,-3,-4;0,-1,-2,-4)$\\
$13^*$&$C^+_{47}+V^+_7+F^+_2$&$(0,0,0,-2;0,0,-2,-2)$\\
$14^*$&$C^+_{56}+V^+_2+C^-_{66}$&$(0,0,-2,-4;0,-2,-3,-4)$\\
$15^*$&$C^+_{68}+U^+_3+V^-_5$&$(0,0,-1,-3;-1,0,-1,-3)$\\
\bottomrule
\end{longtable}\end{center}
These $4+2\cdot11=26$ groups are disjoint; their sums are negative.
Each therefore contains a nonpositive label. This proves the fixed
allowance $84$ by direct symbolic addition, independent of any case search.

\begin{lemma}[Additional fixed credit]\label{lem:common-fixed-credit}
Suppose disjoint collections of fixed labels $\mathcal H,\mathcal L$
are known respectively to lie above their target maxima and at or below
zero. Suppose no baseline group in the displayed table is wholly
contained in $\mathcal H$. Let $g(\mathcal L)$ be the number of baseline
groups meeting $\mathcal L$. Then the fixed allowance can be reduced
from $84$ by at least
\[
e=\card{\mathcal H}+\card{\mathcal L}-g(\mathcal L).
\]
This credit is part of the original deficit, so remaining losses must
share $\delta-e$.
\end{lemma}
\begin{proof}
Count all labels in $\mathcal H\cup\mathcal L$ as unusable. A baseline
group avoiding $\mathcal L$ contains a nonpositive label; such a label
cannot be HIGH, whose score is above a positive target maximum.
It is an additional unusable label outside the two collections. The
disjoint groups supply at least $26-g(\mathcal L)$ such labels. The
total is at least $26+e$, as required. This proof counts labels, so
coincidences of their values cannot invalidate an upper bound.
\end{proof}

\subsection{Peeling and exact slack}
\begin{lemma}[Four-element peeling]\label{lem:common-peeling}
If $F$ contains zero and at least four nonzero elements, delete its two
largest positive and two smallest negative elements, filling from the
other sign if necessary, to obtain $F'$. Then
\[
\card{\FS(F)}\ge\card{\FS(F')}+\card{R(F)}-1.
\]
\end{lemma}
\begin{proof}
Let $\alpha$ and $\beta$ be the sums of the negative and positive
deleted elements, respectively, and $t=\beta-\alpha>0$.
The extremal choice gives $R(F)\subseteq[\alpha,\beta]$.
Both $\alpha+\FS(F')$ and $\beta+\FS(F')$ lie in $\FS(F)$.
In each occupied class modulo $t$, their union has at least one more
point than $\FS(F')$. The interval $[\alpha,\beta]$ has at most one
representative per class, except for its two endpoints. Thus $R(F)$
occupies at least $\card{R(F)}-1$ classes. Each such class either gains
a point from the two translates or supplies a new point of $R(F)$.
Summing proves the inequality.
\end{proof}

Apply it twice to $D\supset D_1\supset E$, with sizes $d,d-4,d-8$.
Deleting four elements affects at most $4d-10$ distinct-element pair
labels. Consequently
\[
\card{\FS(D)}\ge\card{\FS(E)}+2\card{R(D)}-4d+8,
\quad \card{\FS(E)}\le22d+2\Delta(D)+77\le24d+77.
\]
For a positive $r$-element set $X$, define
\[
\mu_X=\card{\FS(X)}-1-T_r,\quad
\delta_0(X)=\card{\{x\in X:2x\notin R(X\cup\{0\})\}},\quad
\rho_X=\mu_X+2-2\delta_0(X).
\]
The increasing-chain bound gives $\mu_X\ge0$; the largest element is
always defective, so $\rho_X\le\mu_X$.

Suppose $E=P\cup\{0\}\cup(-Q)$, with sizes $a,b$ and an opposite
pair. Put $\Lambda=1+T_a+T_b$, $s=\card{\FS(E)}$,
$\mu=\mu_P+\mu_Q$, $\rho=\rho_P+\rho_Q$, and define
\[
\begin{aligned}
\kappa&=s-\Lambda-\mu,&
\theta&=\delta_0(P)+\delta_0(Q)+8-\Delta(D),\\
c&=2\card{2D}+18d+85-\card{\FS(D)},\\
\ell_1&=\card{\FS(D)}-\card{\FS(D_1)}-\card{R(D)}+1,\\
\ell_2&=\card{\FS(D_1)}-s-\card{R(D_1)}+1,\\
g&=4d-10-\bigl(\card{R(D)}-\card{R(D_1)}\bigr).
\end{aligned}
\]
All these quantities are nonnegative. For $\kappa$ use the elementary
sumset inequality; for $\theta$ use the opposite pair to remove zero
as a diagonal exception and charge the eight peeled elements at most
once each. The other signs follow from covering, peeling and pair counts.

\begin{lemma}[Exact slack identity]\label{lem:common-budget}
With these definitions,
\begin{equation}\label{eq:common-budget}
\kappa+\rho+2\theta+c+\ell_1+\ell_2+g=22d+97-\Lambda.
\end{equation}
\end{lemma}
\begin{proof}
Direct cancellation gives
$c+\ell_1+\ell_2+g=22d+77+2\Delta(D)-s$ and
$\rho+2\theta=\mu+20-2\Delta(D)$.
Add $\kappa=s-\Lambda-\mu$. No inequality or second expenditure of
slack is involved.
\end{proof}

\subsection{The retained structural theorem}
The following is the precise structural input used in the case proofs.
Its proof uses the covering and peeling lemmas above, the standard
additive results listed at the end of this section, and the explicitly
specified structural finite lemmas. It is retained as a theorem; the
local exclusions below do not claim to reprove its inverse-theoretic
and finite-exception arguments.

\begin{theorem}[Common integral core at cardinality 86]\label{thm:structure86}
A feasible zero-containing integral set of size $86$, after the outer
four-on-each-sign deletion and two further four-element peels, admits a
positive normalization in which
\[
\card D=78,\quad\card{D_1}=74,\quad
E=P\cup\{0\}\cup(-Q),\quad\card E=70,
\quad a+b=69,\quad15\le a,b\le54.
\]
The source $D$ is integral and $P,Q$ are individually primitive integer
sets. Let $h_P,h_Q,z,w$ be the hull-hole counts of $\FS(P),\FS(Q),
\FS(E),\FS(D)$. Set
\[
I=T_a+T_b-1225=k(k+1),\qquad k=34-\min(a,b)\in\{0,\ldots,19\}.
\]
Then
\[
h_P,h_Q\le\kappa,\quad
z=h_P+h_Q-\kappa,\quad0\le z\le\min(h_P,h_Q)\le\kappa,
\quad0\le w\le4z.
\]
The two extreme subset-sum translates in each peel actually intersect.
Writing $\mathcal T=\kappa+\ell_1+\ell_2+g$, one has
\begin{equation}\label{eq:common-structural}
c+\rho+\mathcal T\le587-I,\qquad
\mathcal T\ge\kappa+146-3z,\qquad
\mathcal T\ge\max(49,z,146-2z).
\end{equation}
Moreover $\rho_X\ge0$ for either core side, and if its subset-sum
support is not an arithmetic progression then
\[
\mu_X\ge |X|-3,\qquad
\rho_X\ge n(|X|),\qquad
n(r)=\begin{cases}r-14,&15\le r\le20,\\r-5,&r\ge21.\end{cases}
\]
In addition, if $R_v$ is the maximum of its positive core side, then
$R_v\le r_v+\rho_v+h(P_v)+2$, and hence
$R_s+R_t\le73+\rho+\kappa+z$.
No assertion about integrality of the original outer eight elements or
activity of a leading covering translate is part of this theorem.
\end{theorem}

For orientation, the constant $146$ arises from the sharp four-extreme
restricted-sum inequality
\[
\eta_D+\eta_{D_1}+g\ge146,
\]
where $\eta$ counts restricted-sum hull holes. Actual overlap bounds
$\eta_{D_1}\le z+\ell_2$ and $\eta_D\le2z+\ell_1$ imply
$\mathcal T\ge\kappa+146-3z$. Since $\mathcal T\ge\kappa\ge z$,
$3\mathcal T\ge146$, giving $\mathcal T\ge49$ by integrality.
The primitive conclusion is important: if $z>0$, neither side support
can be an arithmetic progression. A primitive such support would be
a full integer interval. Therefore in the positive-hole cases
\begin{equation}\label{eq:common-positive-floors}
\mu\ge63,\quad \tau:=\sum P+\sum Q-T_a-T_b
=\mu+h_P+h_Q\ge63+2z,
\quad \rho\ge\begin{cases}59,&\min(a,b)\ge21,\\50,&\min(a,b)\le20.
\end{cases}
\end{equation}
Sorting the two sizes in the scalar formula for $I$ does not rename a
physical sign in any later ownership argument.

\subsection{Actual resources and the residual deficit}
Use physical signs $s,t$, with $P_s=P$, $P_t=Q$, $D_s=D$, $D_t=-D$. Write $h(P_v)$ for
the hull-hole count of a core-side support, reserving $h_v=v_2-v_3$
for an outer gap. Write $r_s=a,r_t=b$,
$R_v=\max P_v$, $M_v=\max D_v$, and $S_v$ for the positive mass in
$D_v$. Thus $\min D_v=-M_o$ for the opposite sign $o$.
Put $R=R_s+R_t$, $M=M_s+M_t$, $K=\card{D+D}$, and
\[
F_v=\FS(D_v)\cap(0,\infty),\qquad
\eta_v=S_v-\card{F_v},\qquad\eta_s+\eta_t=w.
\]
The original outer magnitudes are $v_1>v_2>v_3>v_4>M_v$.
Define $\lambda_v,\alpha_v,\beta_v$ as in the cover and the ten centers
\[
\begin{array}{llll}
O_v=\lambda_v+v_1,&H_v=\lambda_v+v_2,&
\mathrm{own2}_v=\lambda_v+v_3,&\mathrm{own3}_v=\lambda_v+v_4,\\
C4_v=\lambda_v-o_1,&C5_v=\lambda_v-o_2,&
C6_v=\lambda_v-o_3,&C7_v=\lambda_v-o_4,
\end{array}
\]
followed by $\alpha_v,\beta_v$. Their ordinary source is $D_v$.
The leading resource is $L_v=\lambda_v+(D_v+D_v)$, of capacity $K$.

Branch A is $\beta_s+\beta_t\le0$, including equality; branch B is
$\beta_s+\beta_t>0$. In both branches pair1 is $\{C4_s,C4_t\}$.
Pair2 is $\{\beta_s,\beta_t\}$ in A and $\{C5_s,C5_t\}$ in B.
There are two paired units and sixteen singleton units, each whole unit
having capacity $78$. A unit is missing only when all its target
contributions are empty. A retained pair supplies a genuinely active
component, but does not give that selected component a second capacity.
Let $j$ count missing whole units. The unchanged residual deficit is
\begin{equation}\label{eq:common-delta}
\delta=c-78j
=2K+78(18-j)+84-\card{F_s}-\card{F_t}\ge0.
\end{equation}
Equivalently, with $S=S_s+S_t$,
\begin{equation}\label{eq:common-mass}
S=2K+1488-78j+w-\delta.
\end{equation}

\begin{lemma}[Grouping actual losses]\label{lem:common-loss}
Replace any disjoint groups of retained full resources by their actual
target unions, with each ordinary whole unit and each leading resource
used at most once. The sum of the resulting capacity deficiencies is
at most $\delta$. Compatible extra fixed credit $e$ reduces the
available amount to $\delta-e$.
\end{lemma}
\begin{proof}
The target union is covered by the selected groups and the remaining
original resources. Bound the latter by their original capacities and
the fixed allowance. Subtract its cardinality from the full right side
of \eqref{eq:common-delta}. This proves the inequality. A paired unit
uses the sum of its two target contributions against its single
capacity $78$. If only one component contributes, its native loss can
be written $78$ minus that sole contribution. If both contribute,
discarding one contribution would create a false larger loss.
\end{proof}

\begin{lemma}[Rooted forests]\label{lem:common-forest}
On each target choose a forest whose vertices are distinct retained
leading or singleton full resources and whose edges join intersecting
resources. A sole paired component can be included only under a
separately established whole-capacity license. Designate
at most one vertex in each tree component as a root. The sum of all
edge intersection cardinalities, plus all root cardinalities outside
their target, is at most the group capacity deficiency. Across disjoint
groups it is therefore at most $\delta$ (or $\delta-e$ with compatible
fixed credit). Additional losses of unused retained singleton resources
may be included within the same bound.
\end{lemma}
\begin{proof}
At a point $x$, let $n(x)$ count containing vertices, $b(x)$ containing
edges, and $a(x)$ containing roots. Inside the target, the induced
subgraph is a forest, so $b(x)\le n(x)-\mathbf1_{n(x)>0}$.
Outside the target, different designated roots lie in different
original components, hence different induced components, and
$b(x)+a(x)\le n(x)$. Sum these pointwise inequalities and then apply
Lemma~\ref{lem:common-loss}. Restricted copies may lower-bound edge
intersections but never replace a full vertex's original capacity.
\end{proof}

\begin{lemma}[Activity and integrality]\label{lem:common-integrality}
An active ordinary center is integral. If all four own centers on one
sign are active, its $\lambda_v$ and four outer magnitudes are integral.
The same conclusion holds if the leading center is integral and at
least three own centers are active. A leading resource is active
whenever $K>\delta$.
\end{lemma}
\begin{proof}
Activity gives $Z+d=f$ with integral $d\in D_v$, $f\in F_v$, so
$Z\in\Z$. For four own centers use
$\lambda_v=-O_v+H_v+\mathrm{own2}_v+\mathrm{own3}_v$ and
$v_i=(\lambda_v+v_i)-\lambda_v$. If $\lambda_v$ and three $v_i$
are known, its defining equation supplies the fourth. An empty leading
resource has native loss $K$, contrary to Lemma~\ref{lem:common-loss}.
Leading activity itself implies integral $\lambda_v$.
\end{proof}

\begin{lemma}[Single and double raw gaps]\label{lem:common-raw-gaps}
Put $d_v=v_1-v_2$. If $O_v$ is active and $H_v+M_v>0$, then
\begin{equation}\label{eq:common-single-gap}
d_v-M\le56+\eta_v.
\end{equation}
The same bound holds when $\beta_v,H_v$ are active and $O_v$ is empty.
If $O_v,\beta_v$ are both active, then
\begin{equation}\label{eq:common-double-gap}
2(d_v-M)\le57+\eta_v.
\end{equation}
The latter assertion needs no activity of $H_v$. For the first assertion, activity of $H_v$, of $\mathrm{own2}_v$, or of the leading resource is sufficient for its lower guard. The bounds hold even before whole-outer integrality.
\end{lemma}
\begin{proof}
Every ordinary center other than $O_v$ is at most $H_v$; the leading
maximum $\lambda_v+2M_v$ is below $H_v+M_v$. Thus no nonfixed
resource meets $(H_v+M_v,O_v-M_o)$. A positive gap has length $d_v-M$ and contains at least $d_v-M-1$ integers. The upper guard is $O_v-M_o\le S_v$. The stated lower guard holds under $H_v$ activity, under $\mathrm{own2}_v$ activity since $H_v>\mathrm{own2}_v$, or under leading activity through $H_v+M_v>\lambda_v+2M_v\ge1$. Thus the gap lies inside $[1,S_v]$, and the same $55$
raw fixed labels and at most $\eta_v$ target holes cover it.

Write $h_v=v_2-v_3$, $X_v=v_4-M_v$. Since
$\beta_v=h_v-d_v-v_4$, beta activity gives
$h_v-d_v\ge X_v+1>0$. Hence $\alpha_v>\mathrm{own2}_v$.
Every ordinary center other than $O_v,H_v$ is at most $\alpha_v$,
and the leading maximum is below $\alpha_v+M_v$.
This gives the second interval $(\alpha_v+M_v,H_v-M_o)$.
When $O_v$ is empty and $H_v$ active, this one interval proves the
single-gap assertion. When $O_v$ and beta are active, both intervals
lie inside $[1,S_v]$, regardless of $H_v$ activity, and are disjoint.
They must share the same $55$ labels and the same $\eta_v$ holes.
Thus $2(d_v-M-1)\le55+\eta_v$. Nonpositive interval lengths make
all the asserted bounds automatic.
\end{proof}

\subsection{Two elementary source envelopes}
\begin{lemma}[Four positive offsets]\label{lem:common-offsets}
For integers $1\le a<b<c<e$, let $B$ be the number of distinct sums
$x+y>e$ with $x,y\in\{a,b,c,e\}$, repetitions allowed. Then
\[
2B\le a+b+8,\qquad 2B-(a+b+c)\le3.
\]
\end{lemma}
\begin{proof}
There are ten unordered pairs, so the first assertion is immediate
when $a+b\ge12$. For $a+b\le11$ the following elementary cases suffice.

If $a=1$, the pairs $2a,a+b,a+c$ are at most $e$, giving $B\le7$
and settling $b\ge5$. For $b=2$, among pairs avoiding $e$ only $2c$
can add a new value beyond the four $e$-containing sums: $c+2>e$
only at $e=c+1$, when it repeats $e+1$. Thus $B\le5$.
For $b=3,e\ge6$, the four pairs $2,4,6,c+1$ are at most $e$,
giving $B\le6$; $(1,3,4,5)$ has $B=5$.
For $b=4,e\ge8$, the four pairs $2,5,8,c+1$ are at most $e$;
the smaller $(c,e)=(5,6),(5,7),(6,7)$ have $B=6$.

If $a=2$, the pairs $4,b+2$ give $B\le8$, settling $b\ge6$.
For $b=3,e\ge6$, the pairs $4,5,6$ are below the cutoff; besides
$2c$, at most one of $c+2,c+3$ supplies a new value, since either
$e=c+1$ and $c+3=e+2$, or $e\ge c+2$ and $c+2\le e$.
Thus $B\le6$; $(c,e)=(4,5)$ gives $B=5$.
For $b=4,e\ge8$ we have $B\le7$; the smaller $(c,e)$ values
$(5,6),(5,7),(6,7)$ give $6,6,7$.
For $b=5,e\ge10$ again $B\le7$; at
$(6,7),(6,8),(6,9),(7,8),(7,9),(8,9)$ the values are $7,6,6,7,6,7$.

If $a=3$, then $2a\le e$. If $e\ge b+3$, also $a+b\le e$,
so $B\le8$. Otherwise $c=b+1,e=b+2$ and $b+e=2c$ merges
two pair labels, again giving $B\le8$. This settles $b\ge5$.
For $b=4,e\ge8$, three pairs are below the cutoff, giving $B\le7$;
the smaller $(c,e)=(5,6),(5,7),(6,7)$ give $6,6,7$.

If $a=4$, only $b=5,6,7$ remain. For $b=6,7$ the pair $8\le e$
gives $B\le9$. For $b=5,e\ge9$, the pairs $8,9$ give $B\le8$;
the smaller $(c,e)=(6,7),(6,8),(7,8)$ give $7,7,8$.
The final pair is $(a,b)=(5,6)$: for $e\ge10$ we have $B\le9$;
at $(c,e)=(7,8),(7,9),(8,9)$ the values are $7,8,9$.
Each bound is at most $\lfloor(a+b+8)/2\rfloor$.

For the second assertion the first gives $2B-(a+b+c)\le8-c$,
settling $c\ge5$. If $c=3$, necessarily $(a,b)=(1,2)$;
for all $e\ge4$, $B=4$, giving value $2$.
If $c=4$, for $e\ge8$ only the four $e$-containing pairs remain.
For $e=5,6,7$, the $B$ values for $(a,b)=(1,2),(1,3),(2,3)$
are respectively $(5,4,4),(5,5,4),(5,5,5)$. Each gives at most $3$.
These are direct additions of ten pairs, not samples of an unbounded
enumeration.
\end{proof}

On each sign write the four peeled magnitudes of $D$ as
$R_v+a_v<R_v+b_v<R_v+c_v<R_v+e_v=M_v$, with positive integer
offsets. To avoid confusion, $c_v$ here is an offset, whereas $c$ is
the original covering slack. Set
\[
\begin{gathered}
U=a_s+b_s+a_t+b_t\ge6,\qquad C=c_s+c_t,\qquad
J=(e_s-c_s)+(e_t-c_t),\\
L=C+J=M-R,\qquad\mathcal A=I+\tau+\delta,\qquad m_0=78j-263-w.
\end{gathered}
\]
\begin{lemma}[Source-only envelopes]\label{lem:common-envelope}
Without any outer-activity assumption,
\begin{align}
2K&=m_0+\mathcal A+4R+U+2C+J,\label{eq:common-envelope-mass}\\
K&\le2R+C+168,\qquad J+U\le336-m_0-\mathcal A,\label{eq:common-two-max}\\
\mathcal A&\le437-78j+w,\label{eq:common-envelope-A}\\
K-\delta&\ge78j-123-w+I+\tau,\label{eq:common-envelope-Kdelta}\\
K-M&\ge R+78j-425-w+I+\tau+\delta.\label{eq:common-envelope-KM}
\end{align}
\end{lemma}
\begin{proof}
The actual mass is $S=1225+I+\tau+4R+U+2C+J$; combine this
with \eqref{eq:common-mass} to obtain the first identity.
Delete the two source maxima, leaving $76$ points. All their sums
except same-sign pairs of the three remaining peeled points lie in
an integer interval of $2R+C+1$ positions. There are at most $12$
exceptional pair labels. Reinserting the maxima adds at most
$2\cdot76+3=155$ labels. This proves $K\le2R+C+168$ and the
bound on $J+U$. Consequently
\[
K-M=R+\frac{m_0+\mathcal A+U-J}{2}
\ge R+m_0+\mathcal A+U-168
\ge R+m_0+\mathcal A-162,
\]
which is \eqref{eq:common-envelope-KM}.

For the full source the analogous central interval has $2R+L+1$
positions, and the two fringes have at most $B_s+B_t$ values. Thus
$K\le2R+L+1+B_s+B_t$. The first four-offset inequality yields
$J\ge m_0+\mathcal A-18$. Together with the bound on $J+U$ and
$U\ge6$, this gives $\mathcal A\le174-m_0$.
The second four-offset inequality gives $L\ge m_0+\mathcal A-8$.
Since $R\ge69$ and $U+C\ge12$, the mass identity, written as
$2K=m_0+\mathcal A+4R+L+(U+C)$, gives
$2K\ge2m_0+2\mathcal A+280$.
Subtract $\delta$ to obtain \eqref{eq:common-envelope-Kdelta}.
\end{proof}

\subsection{Quantitative layer inequalities}
Let $U_0$ be the sum of the first two actual peeled magnitudes on each
sign and $V_0$ the sum of the last two. Put $C_0=1225+I+\tau$,
$C_1=C_0+U_0$, $S=C_1+V_0$, and $M_0=U_0-C_0$.
The symbol $M_0$ is a layer separation, unrelated to $M=M_s+M_t$.
If $\mu\ge\underline\mu$ and $\rho\ge\underline\rho$, the exact
slack identities and structural overlap give
\begin{align}
V_0-C_1+\delta+2z
&\le-336-2I-78j-\underline\mu-\underline\rho,
\label{eq:common-outer-layer}\\
M_0+\delta-2z&\le-208-2I-78j,\qquad M_0\le0.
\label{eq:common-inner-layer}
\end{align}
Here is the calculation. Write $F_0=\card{\FS(E)}$,
$F_1=\card{\FS(D_1)}$, $F_2=\card{\FS(D)}$.
Substitution in the definitions gives
\[
F_2-2F_1=-337-2I-c-\rho-\mu-2\kappa-2\theta-2\ell_2-2g<0.
\]
The two extreme $\FS(D_1)$ translates intersect and their common
hull has $S-2V_0+1$ integer positions. Thus
\[
2V_0-S+c+\rho+\mu+2\kappa+2\theta+2\ell_2+2g\le-336-2I.
\]
Use $\kappa\ge z$ and $c=78j+\delta$ for the outer row.
For the inner row, $\card{R(E)}\le2R-1$ and adding four points
creates at most $4\cdot70+6$ pair labels. The structural maximum
removal bound $R\le73+\rho+\kappa+z$ then gives
\[
F_1-2F_0\le-209-2I+\rho-\mu+2z-c-2\theta-\ell_1-g.
\]
Actual inner overlap gives $M_0\le1+F_1-2F_0$ and $M_0\le0$;
discard the indicated nonnegative quantities and use $\rho\le\mu$.
This proof uses actual intersections supplied by the structural
theorem, not just intersecting hulls.

\subsection{Common-strip localization and endpoint ranks}
\begin{lemma}[One common strip]\label{lem:common-strip}
Let integer supports $X\subseteq[0,A]$, $Y\subseteq[0,B]$ contain
both endpoints, and have a total of $H$ hull holes. If $A,B\ge H$,
then $[H,A+B-H]\cap\Z\subseteq X+Y$.
If a subset-sum support of total mass $L$ contains $[H,L-H]$,
adjoining a positive integer $x\le L-2H+1$ preserves a central
interval $[H,L+x-H]$.

In particular, if these hypotheses propagate through all absolute
source entries while each signed source mass exceeds $H$, the core
and source holes lie in the same two reflected boundary strips. If
their omitted lower-strip sets are $\Gamma_E,\Gamma_D$, then
\[
\Gamma_D\subseteq\Gamma_E,\qquad
z=2\card{\Gamma_E},\quad w=2\card{\Gamma_D},\quad
0\le w\le z,\quad\eta_s=\eta_t=w/2.
\]
\end{lemma}
\begin{proof}
For $t\in[H,A+B-H]$, the complete hulls have at least $H+1$
representations $t=i+j$. At most $H$ are spoiled by a missing
coordinate. This proves the first assertion. The intervals
$[H,L-H]$ and $[x+H,L+x-H]$ overlap or abut under the stated bound
on $x$, proving the propagation assertion.

For a signed set, replace each negative element by its magnitude and
remember its label. Complementing the selected negative entries
translates this absolute-multiset support to the signed subset-sum
support; equal opposite magnitudes remain separately selectable.
Complementing all entries reflects the support about half its mass.
The central interval therefore leaves exactly two reflected sets
of holes in the common strip $\{1,\ldots,H-1\}$. Selecting no tails
gives containment of the old absolute support in the new one, hence
$\Gamma_D\subseteq\Gamma_E$. The two signed masses exceeding $H$
put the reflected strips on opposite sides of zero, giving the target
hole counts.
\end{proof}

The preceding lemma is conditional on the displayed widths and tail
bounds. A case proof must establish them before using $w\le z$,
parity, or equal target holes. The structural estimate $w\le4z$
alone supplies none of those stronger assertions.

\begin{lemma}[Uniform first-gap-one endpoint ranks]\label{lem:common-endpoint}
Suppose the P4 source has first gap $1$, so each core magnitude on
sign $v$ is at least its rank plus $1$. Write its mass as
$T_{r_v}+r_v+\sigma_v$, with $\sigma_v\ge0$, and assume
$\sigma_s+\sigma_t+I+\delta\le58$ and $\delta\le62-I$.
For an integral individually charged retained ordinary translate
$Z+D_v$, let $\ell_Z,h_Z$ count source points clipped below $1$
and above $S_v$, respectively. If
$\ell_Z,h_Z\le\epsilon=\delta-e$, $e\ge0$, then
\begin{equation}\label{eq:common-END1}
Z\ge r_o+4-\ell_Z,\qquad
Z\le S_v-r_v-3+h_Z.
\end{equation}
This applies to a sole contributing paired component only where its
whole-unit native loss has actually been licensed.
\end{lemma}
\begin{proof}
For $Z\ge1$, at most $Z-1$ negative source values avoid lower
clipping, yielding the stronger $Z\ge r_o+5-\ell_Z$.
For $Z\le S_v$, the positive values give the stronger
$Z\le S_v-r_v-4+h_Z$. It remains to treat $Z\le0$ and $Z>S_v$.

If $Z\le0$, all $r_o+4$ negative entries and zero are lower-clipped.
If $j$ positive entries are clipped, $\ell_Z=74-r_v+j$.
Since $\ell_Z\le\delta\le62<74$, one has $j<r_v$, and these are
the first $j$ core entries. The next core value $x_{j+1}$ satisfies
$Z\ge1-x_{j+1}$. If $x_{j+1}\ge j+3$, monotone deviations from
the baseline $2,3,\ldots,r_v+1$ give
\[
\sigma_v\ge r_v-j=74-\ell_Z\ge74-\delta+e,
\]
contrary to $\sigma_v\le58-I-\delta$. Hence $x_{j+1}=j+2$,
and $Z\ge-j-1=r_o+4-\ell_Z$.
For $Z>S_v$, count upper-clipped negative magnitudes after all
$r_v+4$ positive entries and zero. The identical argument gives
$h_Z=74-r_o+j$, the next magnitude $j+2$, and
$Z\le S_v+j+2=S_v-r_v-3+h_Z$.
\end{proof}

\begin{lemma}[Endpoint support propagation]\label{lem:common-endpoint-bridge}
Let $B$ consist of $m$ distinct positive integers, $L=\sum B$, and let $H\subset[1,L/2)\cap\Z$ satisfy
\[
\FS(B)=[0,L]\cap\Z\setminus\bigl(H\cup(L-H)\bigr),
\quad q=\card H,\quad h=\max H,\quad2h<L,
\]
with $h=q=0$ when $H$ is empty. Let $x>\max B$ be integral and
\[
\card{\FS(B)\cup(x+\FS(B))}-\card{\FS(B)}=m+1+e_x,
\qquad e_x\le E.
\]
If
\[
m\ge\max\left(h,1+\left\lfloor
\frac{1+\sqrt{9+8(E+2q+2h)}}2\right\rfloor\right),
\]
then adjoining $x$ preserves the lower holes $H$ and their reflected
upper holes, and $x=m+1+e_x$.
\end{lemma}
\begin{proof}
We have $L\ge T_m$, $x>\max B\ge m\ge h$, and the strict threshold
gives $T_m-2h-2q>m+1+E$. If $x\ge L-2h$, the shifted support has
at least $\min(x,L+1)-2q\ge L-2h-2q$ points above $L$, contradicting
the permitted increment. Thus $x\le L-2h-1$.

The filled central intervals are
$[h+1,L-h-1]$ and $[x+h+1,x+L-h-1]$. They overlap or abut as
integer intervals and fill $[h+1,L+x-h-1]$. Since $x>h$, the lower
holes remain exactly $H$. Subset complementation gives precisely
$L+x-H$ as the upper holes; the endpoints are subset sums. This also
covers $H=\varnothing$. The cardinality increases from $L+1-2q$
to $L+x+1-2q$, proving the assertion. In particular the right endpoint
of the old filled central interval is $L-h-1$, not $L-h$.
\end{proof}

\subsection{The exhaustive coarse partition}
\begin{proposition}[Thirteen actual strata]\label{prop:common-partition}
Every hypothetical normalized size-86 configuration lies in exactly
one of
\[
Z_0,\ldots,Z_5,\quad P_0,\ldots,P_5,\quad H,
\]
where $Z_j$ means $z=0$ and exactly $j$ empty whole units,
$P_j$ means $1\le z\le14$ and exactly $j$ empty whole units,
and $H$ means $z\ge15$.
\end{proposition}
\begin{proof}
If $z=0$, \eqref{eq:common-structural} gives
$0\le\delta\le441-I-78j$, which is negative for $j\ge6$.
For $1\le z\le14$, use $\mathcal T\ge118$ and the non-AP floors:
\[
0\le\delta\le469-I-78j-n(r_s)-n(r_t).
\]
The sum of the floors is $59$ when both sizes are at least $21$, and
$50$ otherwise. For $j\ge6$ the last bound is at most
$469-468-50<0$. All equality and zero-residual cases are retained.
The three ranges for $z$ are disjoint and exhaustive. Within the first
two ranges the actual integer $j$ is unique. No enumeration or
uniqueness of an auxiliary descriptor is involved.
\end{proof}

\subsection{What remains an external or finite input}
The preceding chapter proves the elementary covering, loss and source
lemmas used repeatedly below. Its retained inputs are the preceding
global theorem $N\le86$ and Theorem~\ref{thm:structure86} with the
following mathematical dependencies:
\begin{enumerate}
\item Carpenter--Defant--Kravitz subset-sum stability and the
zero-increment equality theorem; the Balandraud--Girard--Griffiths--
Hamidoune finite-group estimates together with Kneser's theorem;
Stanchescu's simultaneous inverse and short-diameter theorems; and
Lev's central-interval theorem.
\item The general defect and small-increment deductions, including
$\rho_X\ge |X|-5$ for non-progression supports with $|X|\ge21$.
The latter uses the complete $2{,}340$-case deleted-interval base
and its written signature/depth reduction to arbitrary sizes.
\item The structural scalar inequalities and four fixed nonnegative
integer row combinations for the $152$ exceptional systems, as well
as the sharp four-extreme restricted-sum theorem. The structural
proof establishes lattice and overlap before applying these integer
hull-hole estimates.
\item Only in the high-hole chapter: the additional bounded-density,
profile and direct-defect finite lemmas, at their stated domains.
\end{enumerate}
These are mathematical inputs, not claims that the accompanying
program formalizes all reasoning from first principles. The program's
documented checks distinguish finite arithmetic from the analytical
reductions and the published theorems. The source map identifies full
proofs of the retained intermediate results.


\section{No empty unit and one empty unit}\label{sec:easy}
Throughout this section $v,o$ denote opposite physical signs. We use the common centers $O_v,H_v,\own2_v,\own3_v,C4_v,\ldots,C7_v,\alpha_v,\beta_v$. Put
\[
t_v=d_v-M,\qquad G=g_s+g_t,\quad X=X_s+X_t.
\]
The activity identities needed repeatedly are
\begin{align}
C7_v\text{ active}&\ \Longrightarrow\ t_v\ge2X_v+X_o+g_v+1,\label{eq:e-c7}\\
C5_v\text{ active}&\ \Longrightarrow\ t_v\ge2X_v+X_o+G+h_o+1,\label{eq:e-c5}\\
C4_v\text{ active}&\ \Longrightarrow\ t_v\ge2X_v+X_o+G+h_o+d_o+1,\label{eq:e-c4}\\
\beta_v\text{ active}&\ \Longrightarrow\ h_v\ge d_v+X_v+1.\label{eq:e-beta}
\end{align}
They follow by expanding the literal centers and using $u+M_v\ge1$ for an integral active center $u$. The first implication also holds when $C6_v$ is active, because $C7_v>C6_v$.

\subsection{One proof for Z0 and P0}\label{sec:z0p0}
\begin{proposition}[Z0 and P0]\label{prop:z0p0}
The original covering grammar with $j=0$ is impossible whenever $w\le56$.
\end{proposition}
\begin{proof}
All four own centers on each side are active and integral. The identity
\[
\lambda_v=-O_v+H_v+\own2_v+\own3_v
\]
therefore makes $\lambda_v$ and all eight outer magnitudes integral. Thus $d_v,h_v,g_v,X_v$ are positive integers. Since $R\ge69$, the four source peels give $M\ge77$. Equation~\eqref{eq:e-c7} gives $t_v\ge5$ and hence $d_v\ge82$.

An active beta exists: in A choose an actual contributor of the retained beta pair, while in B both beta singleton units are active. Name its side $s$. The raw-gap inequalities give
\[
t_s\le\frac{57+\eta_s}{2}\le\frac{113}{2},\qquad
 t_t\le56+\eta_t\le112.
\]
If $C4_s$ were active, \eqref{eq:e-c4} would give $t_s>d_t\ge82$. If $C4_t$ were active, \eqref{eq:e-c4} and \eqref{eq:e-beta} would give $t_t>h_s+d_s>2d_s\ge164$. Both alternatives contradict the displayed caps. Both components of pair1 are empty, although $j=0$ requires that whole unit to be retained. This proves the proposition.
\end{proof}
For Z0, $w=0$. For P0, the original $w\le4z$ and $z\le14$ give $w\le56$. No localized-hole theorem or leading-resource activity is needed in this argument.

\subsection{The complete Z1 grammar}\label{sec:z1}
\begin{proposition}[Z1]\label{prop:z1}
There is no original configuration with $z=0$ and $j=1$.
\end{proposition}
\begin{proof}
Here $w=0$ and both targets are complete intervals. We first prove outer integrality without assuming leading activity. At least one side $o$ has all four own centers active, so the preceding identity gives integral $\lambda_o,o_1,\ldots,o_4$. If the other side has a missing own center, that is the unique missing unit. Hence $C6_o,C7_o$ are active, giving integral $v_3,v_4$. One of $\own2_v,\own3_v$ is active, giving integral $\lambda_v$. Now $d_v=\lambda_v+v_3+v_4$ is integral; an active $O_v$ or $H_v$ determines one of $v_1,v_2$, and their difference determines the other.

At least one of $C6_v,C7_v$ is active on each side. Therefore \eqref{eq:e-c7} gives $d_v\ge82$. If $O_v$ is active, the raw gap gives $t_v\le56$: if $H_v$ is the missing unit, the retained $\own2_v<H_v$ still provides the lower endpoint guard. If $\beta_v,H_v$ are active, the second raw gap also gives $t_v\le56$. If all three $\beta_v,O_v,H_v$ are active, the two disjoint gaps together give $2t_v\le57$.

Suppose pair1 is missing. In A choose an active beta contributor on side $v$ and the opposite retained $C5_o$; in B choose an active $C5_o$ contributor and the opposite retained beta singleton. Equations~\eqref{eq:e-c5} and \eqref{eq:e-beta} give $t_o>h_v>d_v\ge82$, while retained $O_o$ gives $t_o\le56$. Thus pair1 is retained.

Choose an active $C4_v$. Equation~\eqref{eq:e-c4} gives $t_v>d_o\ge82$, so $O_v$ must be the unique missing unit. Since $H_v$ remains active, $\beta_v$ cannot be active. Branch B is therefore impossible. In A, the retained beta pair has its actual contributor on the opposite side $u=o$. Consequently
\[
O_v\text{ is the sole missing unit},\quad \beta_v\text{ is empty},\quad
\beta_u,O_u,H_u\text{ are active},\quad 2t_u\le57.\tag{Z1.1}
\]
Every own2, own3, $C7$, alpha and $C5$ singleton used below is retained.

Write $\tau_v=\sum P_v-T_{r_v}$ and $\tau=\tau_s+\tau_t$. The common two-maximum source estimate at $j=1,w=0$ is
\[
K-M\ge R-347+I+\tau+\delta.\tag{Z1.2}
\]
This is equation~\eqref{eq:common-envelope-KM} with $j=1,w=0$; it uses the unchanged original deficit.

Delete the largest six core entries on each sign, retaining zero, to obtain a $58$-point set $Q$. For sorted core entries $b_i=i+e_i$, the deviations are nonnegative and nondecreasing. The six deviations at ranks $r_v-6,\ldots,r_v-1$ are each at least $e_{r_v-6}$, and their sum is at most $\tau_v-(R_v-r_v)$. Thus
\[
\operatorname{span}(Q)\le57+\frac{\tau-R+69}{6}.
\]
Two translates of $Q$ at nonnegative integral displacement $q$ therefore intersect in at least $f-q$ points, where
\[
6f=279+R-\tau.\tag{Z1.3}
\]
On each target use the full retained-resource tree
\[
L_v\,--\,\own3_v\,--\,\own2_v,\qquad L_v\,--\,C7_v.
\]
Copies of $Q$ inside these resources certify gaps $X_v,g_v,X_o$. The six edge lower bounds sum to $6f-2X-G$. Designate $L_u$ as the only root and put $\ell_u=K-\card{L_u\cap F_u}$. The rooted-forest lemma gives
\[
6f-2X-G+\ell_u\le\delta.\tag{Z1.4}
\]
The active $C7_u$ guarantees a positive leading maximum $d_u-2X_u-g_u$. Since the $K-\ell_u$ target points of $L_u$ are distinct positive integers below this maximum,
\[
d_u\ge K+2X_u+g_u-\ell_u.
\]
Also $\alpha_v$ and $C5_u$ are active, and their exact sum is
$\alpha_v+C5_u=d_u-2(v_4+u_4)-G$. Hence $t_u\ge2X+G+2$.
Combining these facts with (Z1.4), and then (Z1.2)--(Z1.3), yields
\[
2t_u\ge K-M-\delta+6f+2
\ge2R-66+I\ge72+I.
\]
This contradicts $2t_u\le57$ from (Z1.1), exhausting Z1.
\end{proof}


\section{Shared one unit loss arguments}\label{sec:one-tools}
This section proves the reductions used in the positive-hole one-unit cases. The hypotheses are recorded explicitly so that no conclusion from one hole range is imported into another.

\subsection{A star with one empty leaf and the two cap obstruction}
\begin{lemma}[One empty leaf]\label{lem:one-empty-leaf}
In the two leading stars with extreme leaves $\own3_s,\own3_t,C7_s,C7_t$, suppose at most one leaf is empty and every other leaf is a retained singleton. Then, once the outer coordinates are integral,
\begin{equation}\label{eq:star4}
2X+4R+\delta\ge556.
\end{equation}
\end{lemma}
\begin{proof}
A translate of the $70$-point core has overlap at least $139-R-q$ with a second translate at gap $q$, by counting their union in a hull. The four leading-to-extreme overlaps have gaps $X_s,X_t,X_s,X_t$, so their lower bounds sum to $556-4R-2X$.

Only the leading and retained leaves carry capacity. At a point of the leading, let $k$ retained leaves and $e\in\{0,1\}$ empty leaves also contain it. Inside the actual target, $e=0$ and the union deficiency is $k$, the edge count. Outside it, the deficiency is $1+k\ge k+e$. Thus all four overlaps, including the possible empty-leaf overlap, are bounded by the same retained-resource deficiency. Summing over the separately tagged targets gives \eqref{eq:star4}. This proof licenses one empty leaf, not arbitrarily many.
\end{proof}

\begin{lemma}[Four top tail]\label{lem:four-top-tail}
For the original cover with $j$ empty whole units,
\begin{equation}\label{eq:tail}
d_s+d_t+G+\delta-2K\ge806-78j.
\end{equation}
\end{lemma}
\begin{proof}
For $u\in F_v$, add the four largest outer magnitudes $v_1+v_2+v_3+v_4$. This is a subset sum of the original feasible set. If a representing four-term sum uses at most two occurrences from $\{v_1,v_2\}$, it is at most $2v_1+2v_3$, so $u\le d_v+g_v$. Above this threshold, three occurrences from the top two give four multisets and at most $84$ choices for the remaining summand, or $336$ values; four top occurrences give at most five values. Consequently $\card{F_v}\le d_v+g_v+341$. Sum over signs and use $\card{F_s}+\card{F_t}=2K+1488-78j-\delta$.
\end{proof}

\begin{lemma}[Two caps]\label{lem:two-caps}
Suppose both leading resources are active, the outer coordinates are integral, the star estimate \eqref{eq:star4} holds, and
\[
d_v-M\le56+\eta_v\quad(v=s,t),\qquad \eta_s+\eta_t=w.
\]
Then
\begin{equation}\label{eq:two-caps}
4I+4\tau+\delta+234j\le562+6w.
\end{equation}
\end{lemma}
\begin{proof}
The leading maximum on side $v$ is $d_v-2X_v-g_v$. Its actual loss is $\ell_v$, and $\ell_s+\ell_t\le\delta$. Thus
\[
d_s+d_t\ge2K+2X+G-\delta.
\]
Put $Y=2M-2K+112+w+\delta$. The two caps, the last inequality and \eqref{eq:tail} imply
\[
G+2X\le Y,\qquad G\ge806-78j-Y.
\]
With \eqref{eq:star4}, this gives $2Y\ge1362-78j-4R-\delta$. The common source estimate
$K-M\ge R-425+78j-w+I+\tau+\delta$ gives
\[
Y\le962-156j+3w-2R-2I-2\tau-\delta.
\]
Comparing the two bounds proves \eqref{eq:two-caps}.
\end{proof}

\subsection{A first gap baseline without a profile catalogue}
\begin{lemma}[Small first gap]\label{lem:first-gap-seven}
Suppose the labeled absolute-value core has a missing positive sum, both side sizes are at least $17$, and
\[
I+\tau\le375.
\]
Then its first missing positive integer $m$ satisfies $1\le m\le7$. If $k_v$ is the number of entries on side $v$ below $m$, then $k=k_s+k_t\le4$ and $k_v\le3$. Keep those actual low entries and replace every higher rank on side $v$ by the consecutive baseline $m+1,\ldots,m+r_v-k_v$. If $B_v=r_v+m-k_v$, $B=B_s+B_t$, and $b$ is this baseline's total excess over $T_{r_s}+T_{r_t}$, then
\begin{equation}\label{eq:baseline-286}
R\ge B,\qquad \tau=b+\sigma\quad(\sigma\ge0),\qquad I+b-5B\ge-286.
\end{equation}
\end{lemma}
\begin{proof}
The greedy first-failure rule says that the actual entries below $m$ sum to $m-1$ and generate $[0,m-1]_{\Z}$; there is no entry equal to $m$. Their two side masses are at least $T_{k_s},T_{k_t}$. Summing the coordinatewise high baselines and completing a square gives
\begin{equation}\label{eq:first-gap-square}
I+\tau\ge70m-\frac54-(m+35)k+\frac{k^2}{4}
 +\left(r_s-\frac{69}{2}-\frac{k_s-k_t}{2}\right)^2.
\end{equation}
The same expression is exactly $I+b$ for the actual low entries and high baseline. This identity uses their total low mass $m-1$, without enumerating low-entry patterns.

Since $m$ precedes a remaining coefficient, $m<\max R_v\le52+\tau\le427$. Also $k(k+2)/4\le m-1$, so $k\le40$. If $m\ge8$ and $k\le4$, discarding the nonnegative quadratic terms in \eqref{eq:first-gap-square} gives
\[
I+\tau\ge66m-565/4\ge1547/4>375.
\]
If $5\le k\le40$, substitute $m\ge k(k+2)/4+1$ and use $70-k>0$ to obtain
\[
I+\tau\ge q(k)/4,\qquad q(k)=-k^3+69k^2-4k+275.
\]
For $5\le k\le39$, $q(k+1)-q(k)=-3k^2+135k+64>0$, while $q(5)=1855$. This is again greater than $4\cdot375$. Thus $m\le7$, and the triangular low-mass bounds give $k\le4,k_v\le3$.

Finally $B=69+2m-k$. Subtract $5B$ from the equality for $I+b$. If $m=1$, then $k=0,B=71,b=69$, giving $I+b-5B=I-286$. For $m=2$, $k\le1$ gives at least $-1033/4$ after dropping the nonnegative quadratic terms. For $m=3$, $k\le2$ gives at least $-929/4$. For $4\le m\le7$, $k\le4$ gives at least $56m-1865/4\ge-969/4$. Each bound is at least $-286$. Since $m\le7<r_v$, each side has a high rank, so its baseline maximum satisfies $B_v\le R_v$; the baseline construction is therefore valid on both sides.
\end{proof}

\begin{lemma}[The final missing O form]\label{lem:one-hole-terminal}
Suppose $j=1$, the first-gap hypotheses of Lemma~\ref{lem:first-gap-seven} hold, both leading resources are active, all outer coordinates are integral, and $\eta_s=\eta_t=w/2$. Suppose the unique missing whole unit is $O_v$, the branch is A, $\beta_v$ is empty and the retained beta pair contributes on $u=o$. Then
\begin{equation}\label{eq:one-hole-terminal}
131-w\le2(d_u-M)\le57+w/2,
\qquad\text{hence}\qquad74\le3w/2.
\end{equation}
\end{lemma}
\begin{proof}
All other singleton units are retained. Delete the largest six core entries per sign, obtaining an actual $58$-point set $Q$. Its retained boundary ranks are above all low entries because $r_v\ge17$ and $k_v\le3$. On the high ranks, deviations from the consecutive baseline are nonnegative and nondecreasing. The six ranks $r_v-6,\ldots,r_v-1$ each have deviation at least that of rank $r_v-6$, while their total is at most $\sigma_v-(R_v-B_v)$. Therefore
\[
\operatorname{span}(Q)\le B-12+\frac{\sigma-R+B}{6},\qquad
6f=762-7B+R-\sigma
\]
and translates at gap $q$ overlap in at least $f-q$ points.

On each target take the full retained-resource tree
$L_v-\own3_v-\own2_v$ and $L_v-C7_v$. Restricted copies of $Q$ certify gaps $X_v,g_v,X_o$. With only $L_u$ designated as root, the common forest lemma gives
\[
6f-2X-G+\ell_u\le\delta.
\]
The leading maximum bounds its actual target points by
$d_u\ge K+2X_u+g_u-\ell_u$. Activity of $\alpha_v,C5_u$ and the identity
$\alpha_v+C5_u=d_u-2(v_4+u_4)-G$ give $d_u-M\ge2X+G+2$.
Eliminating $\ell_u$ and then the opposite $2X_v+g_v$ gives
\[
2(d_u-M)\ge K-M-\delta+6f+2.
\]
The common $j=1$ source envelope yields
\[
\begin{aligned}
2(d_u-M)&\ge417-7B+2R+I+b-w\\
&\ge417-5B+I+b-w\ge131-w.
\end{aligned}
\]
The double raw gap on side $u$ gives the upper bound in \eqref{eq:one-hole-terminal}. The two bounds use one and the same original deficit; every occurrence cancels algebraically.
\end{proof}


\section{The positive hole case with one empty unit}\label{sec:p1}
\begin{proposition}[P1]\label{prop:p1}
No original configuration has $1\le z\le14$ and $j=1$.
\end{proposition}
We prove the statement in four steps: a source estimate, common-strip localization, two-cap exclusion, and the shared terminal contradiction.

\subsection{The original scalar range and source estimates}
Positive core holes and individual primitivity exclude arithmetic-progression subset-sum supports. Thus the general non-AP input gives
\[
\mu\ge63,\qquad \tau\ge63+2z,\qquad
\rho\ge n(r_s)+n(r_t).
\]
The original budget gives $\delta\le391-I-n(r_s)-n(r_t)$. If the smaller side were at most $16$, then $I\ge342$ and the defect sum would be at least $50$, making this upper bound negative. Hence
\begin{equation}\label{eq:p1-range}
17\le r_s,r_t\le52.
\end{equation}
Put $\mathcal A=I+\tau+\delta$. The common source envelope at $j=1$ gives
\begin{equation}\label{eq:p1-envelope}
\mathcal A\le359+w,\qquad K-M\ge R-347-w+\mathcal A.
\end{equation}
The same lemma gives
\begin{equation}\label{eq:p1-J}
K-\delta\ge-45-w+I+\tau,\qquad
J+U\le521+w-\mathcal A.
\end{equation}
Here $U=\sum_v(a_v+b_v)$ and $J=\sum_v(e_v-c_v)$ use the common peel offsets, not the outer covering centers.

\subsection{Localization using all subsets of each peel}
Let $\mathcal H=h(P_s)+h(P_t)=\kappa+z$. Then $\tau=\mu+\mathcal H$, and \eqref{eq:p1-envelope} gives
\begin{equation}\label{eq:p1-hole-strip}
\mathcal H\le296+w-I-\delta\le352-I-\delta.
\end{equation}
For each side size $r$, put $t=r-23$. The identity
\[
T_r+I=408+\frac{t(3t+1)}2\ge408
\]
shows that each core mass $C_v=\sum P_v$ exceeds $\mathcal H$. The common strip lemma therefore gives
\[
[\mathcal H,C_0-\mathcal H]_{\Z}\subseteq B_E,
\qquad C_0=C_s+C_t=1225+I+\mu+\mathcal H,
\]
where $B_E$ is the subset-sum support of the labeled absolute-value core. Its central width is
\begin{equation}\label{eq:p1-width}
L_0=C_0-2\mathcal H\ge936+2I+\delta\ge936.
\end{equation}

Let $U_0$ be the total of the first two peeled source magnitudes on each side and $V_0$ the total of the final two. Put $C_1=C_0+U_0$ and $S=C_1+V_0$. Each inner magnitude is at most $R_v+b_v$. The other three terms in $U$ sum to at least $4$, and $J\ge2$, so \eqref{eq:p1-J} implies
\[
R_v+b_v\le r_v+\tau_v+515+w-\mathcal A
\le r_v+515+w-I-\delta\le623.
\]
For a list of positive entries, adjoining one entry $x$ to a subset-sum set joins the old set and its translate; the largest consecutive subset-sum gap is at most the largest entry in the list. More precisely, in any fixed order it is at most
\[
\max\bigl(t_1,t_2-t_1,\ldots,t_n-\textstyle\sum_{i<n}t_i\bigr).
\]
This follows by induction: disjoint hulls create only their between-hull gap, and overlapping hulls create no larger gap than the old within-hull gaps. Repeated entries cause no change. Thus the actual interval in \eqref{eq:p1-width}, shifted by every subset of the four inner entries, joins to give
\[
[\mathcal H,C_1-\mathcal H]_{\Z}\subseteq B_{D_1}.
\]
Its width $L_1=C_1-2\mathcal H$ is at least $1080$, since $U_0\ge2R+6\ge144$.

For the last four entries, sorted $t_1\le\cdots\le t_4$, the first three gap terms are at most $t_3\le V_0/2$. The last gap term is at most the difference between $t_4$ and its same-sign companion, hence at most $J\le571$. The quantitative outer-layer row is
\[
C_1-V_0\ge Q_0:=414+\mu+\rho+2I+\delta+2z.
\]
Using $C_1\ge1369+I+\mu+\mathcal H$, $\mu\ge63$, $\rho\ge50$, and \eqref{eq:p1-hole-strip},
\[
C_1+Q_0-4\mathcal H
\ge1959+3I+\delta+2z-3\mathcal H
\ge903+6I+4\delta+2z>0.
\]
Hence $V_0/2<L_1$, as is $J$. All actual final subset shifts therefore join, proving
\[
[\mathcal H,S-\mathcal H]_{\Z}\subseteq B_D.
\]
The two supports are symmetric by complementing labeled selections. Their holes lie in the same lower strip $[1,\mathcal H-1]_{\Z}$ and its disjoint upper reflection. Since $B_E\subseteq B_D$, the common strip lemma now gives
\begin{equation}\label{eq:p1-localized}
z,w\text{ even},\quad0\le w\le z\le14,\quad
\eta_s=\eta_t=w/2,\quad \mathcal A\le373.
\end{equation}
Every interval used here is genuinely filled. No activity is inferred merely from overlap of unfilled hulls.

\subsection{Integrality and the impossibility of two caps}
The earlier bound for $K-\delta$, together with $\tau\ge63+2z$ and \eqref{eq:p1-localized}, gives
\[
K-\delta\ge18+I+2z-w>0.
\]
An inactive leading resource would cost $K$, so both are active and their shifts $\lambda_v$ are integral. At most one own singleton is missing; each side therefore has at least three integral own centers. These determine three outer magnitudes, and the identity $\lambda_v=v_1-v_2-v_3-v_4$ determines the fourth. All outer coordinates are integral.

We require one additional four-edge estimate:
\begin{equation}\label{eq:p1-four}
2X_v+X_o+g_v\ge256-\tau-\delta\qquad(v=s,t).
\end{equation}
Delete the largest three core entries per sign to obtain $64$ points. Its span is at most $63+\tau/4$, because the retained side maximum is at most $r_v-3+\tau_v/4$. Thus two copies at gap $q$ overlap in at least $64-\tau/4-q$ points. Select a two-edge arm with gaps $X_v,g_v$, its unused $X_v$ star edge on the opposite target, and either $X_o$ star edge. Their four gaps sum to $2X_v+X_o+g_v$.

Here is the precise charge if the one missing unit is a selected vertex. On each target, positive restricted copies lie at or above $\lambda_v+M_v-R_o$, while negative restricted copies lie at or below $\lambda_v-M_o+R_v$. These zones are disjoint since $M>R$. At a point, the selected restricted edges therefore lie in at most one arm. Inside the target an empty vertex is absent and induced-forest counting applies. Outside, at most one empty vertex raises the edge count by at most one, which is covered by the retained membership count. Group only full retained resources, count each leading once, and apply the original deficit. This proves \eqref{eq:p1-four} without restoring capacity to the empty unit. The same argument, or Lemma~\ref{lem:one-empty-leaf}, proves \eqref{eq:star4}.

Let $T_v$ mean that $O_v$ is active or that both $\beta_v,H_v$ are active. The raw gaps give $t_v\le56+\eta_v$ whenever $T_v$ holds. For an active $O_v$, leading activity supplies the lower gap guard even if $H_v$ is missing. For the second alternative, beta activity supplies that guard. If both $T_s,T_t$ held, Lemma~\ref{lem:two-caps} with $j=1$ would give
\[
4I+4\tau+\delta\le328+6w,
\qquad \tau+\delta\le139+6w-6z-4I\le139.
\]
At least one $C7_v$ is active, so \eqref{eq:e-c7}, \eqref{eq:p1-four}, and that side's cap instead give
\[
\tau+\delta\ge201-\eta_v=201-w/2\ge194.
\]
This contradiction proves that the two cap predicates cannot coexist.

\subsection{The unique remaining form}
If both $O$ units were retained, both cap predicates would hold. Thus $O_v$ is the unique missing whole unit. In B, the retained $\beta_v,H_v$ force the forbidden second cap, so B is excluded. In A, $\beta_v$ must be empty and the retained beta pair must contribute on $u=o$. All other singletons are retained. Lemma~\ref{lem:first-gap-seven} applies because \eqref{eq:p1-range} gives the sizes, positive $z$ gives a first gap, and $I+\tau\le\mathcal A\le373$. Lemma~\ref{lem:one-hole-terminal} now gives
\[
74\le3w/2\le21,
\]
a contradiction. This proves Proposition~\ref{prop:p1} in both physical orientations and for every original paired choice.


\section{The positive hole case with five empty units}\label{sec:p5}
\begin{proposition}[P5]\label{prop:p5}
No original configuration has $1\le z\le14$ and $j=5$.
\end{proposition}
\begin{proof}
The original budget and $\mathcal T\ge146-2z\ge118$ give
\[
\delta+\rho\le79-I.
\]
If the smaller side were at most $20$, then $I\ge210$ and the generic paired defect floor would be $50$, impossible. Thus both sides are at least $21$ and $\rho\ge59$. Consequently
\begin{equation}\label{eq:p5-domain}
\delta\le20-I,\qquad I=k(k+1)\le20,\qquad30\le r_s,r_t\le39.
\end{equation}
The boundary $k=4,I=20,\delta=0$ is included. Positive holes also give $\mu\ge63$ and $\tau\ge63+2z$.

Set $\mathcal H=h(P_s)+h(P_t)=\kappa+z$. The stronger structural budget gives
\[
390+\delta+\rho+\kappa+146-3z\le587-I,
\]
so
\begin{equation}\label{eq:p5-strip}
\mathcal H\le-8+4z-I-\delta\le48.
\end{equation}
Each core mass is at least $T_{30}=465>\mathcal H$. Hence the common strip lemma gives the filled interval $[\mathcal H,C_0-\mathcal H]_{\Z}$ in the labeled absolute-value core support.

Let $U_0,V_0$ be the total absolute masses added in the first and second four-point peel layers, respectively; put $C_1=C_0+U_0$ and $S=C_1+V_0$. The quantitative layer rows at $j=5$ give
\[
C_0-U_0\ge598+2I+\delta-2z,
\qquad C_1-V_0\ge848+2I+\delta+2z.
\]
Subtracting twice \eqref{eq:p5-strip} yields
\[
C_0-U_0-2\mathcal H\ge614+4I+3\delta-10z\ge474,
\]
\[
C_1-V_0-2\mathcal H\ge864+4I+3\delta-6z\ge780.
\]
At each layer, the actual filled interval and its translate by the whole added mass therefore overlap. Selecting none or all of that layer proves successively
\[
[\mathcal H,C_1-\mathcal H]_{\Z}\subseteq B_{D_1},\qquad
[\mathcal H,S-\mathcal H]_{\Z}\subseteq B_D.
\]
Complement symmetry and the same lower strip now imply $w\le z$. Evenness follows too, but is not needed.

A rough elementary source estimate suffices, so we give it instead of using the sharper offset theorem. Write the peel offsets as $a_v<b_v<c_v<e_v$ and put
\[
u=\sum_v(a_v+b_v)\ge6,\quad C=c_s+c_t,\quad
J=\sum_v(e_v-c_v),\quad\mathcal A=I+\tau+\delta.
\]
The original identity $S=2K+1098+w-\delta$ gives
\[
2K=127-w+\mathcal A+4R+u+2C+J.
\]
Deleting the two source maxima gives $K\le2R+C+168$. For the full source, all sums except same-sign pairs of the four peeled points lie in a hull of $2R+C+J+1$ positions; there are at most $20$ exceptional unordered pair labels. Thus $K\le2R+C+J+21$. With $m_0=127-w$, comparison gives
\[
J+u\le336-m_0-\mathcal A,\qquad
J\ge m_0+\mathcal A+u-42.
\]
Combining them and using $u\ge6$ proves $\mathcal A\le56+w$. Hence
\[
I+63+2z+\delta\le I+\tau+\delta\le56+w\le56+z,
\]
which would require $7+z+I+\delta\le0$. This is impossible. The argument never uses an outer-coordinate integrality assumption or a chosen paired owner, so it covers all original branches and placements.
\end{proof}


\section{The high hole family}\label{sec:h}
The high-hole route uses two additional, precisely scoped defect results. They are finite-backed inputs; their finite bases and supplied inference certificates are part of the companion verification boundary. The general induction and inverse-theorem implications remain mathematical inputs unless explicitly proved in the foundation sources.

\begin{theorem}[Dense-hole defect inputs]\label{thm:dense-input}
For the primitive core sides in the original structural domain, the paired defect satisfies
\[
\begin{array}{c|cccc}
z&15\text{--}16&17\text{--}20&21\text{--}48&\ge49\\\hline
\rho&\ge300&\ge345&\ge391&\ge436.
\end{array}
\]
Separately, if $X\subset\Z_{>0}$ is primitive, $31\le\card X\le38$, and its subset-sum support has at least $15$ hull holes, then
\[
\rho(X)\ge8\card X-112.
\]
The general maximum-removal input, for a side of size $r\ge15$, is
$\max X\le r+\rho(X)+h(X)+2$.
\end{theorem}
The paired floors apply to all physical sizes $15\le r_s,r_t\le54$ with sum $69$, not just balanced sizes. The direct bound is used only on its displayed $31$--$38$ domain.

\subsection{The whole family has at most one empty unit}
The structural budget and Theorem~\ref{thm:dense-input} give
\begin{center}
\begin{tabular}{@{}lll@{}}\toprule
Range&Scalar consequence&Reason\\\midrule
$15\le z\le16$&$c\le173-I$&$\rho\ge300$, $\mathcal T\ge146-2z$\\
$17\le z\le20$&$c\le136-I$&$\rho\ge345$, $\mathcal T\ge146-2z$\\
$21\le z\le48$&$c\le146-I$&$\rho\ge391$, $\mathcal T\ge146-2z$\\
$z\ge49$&$c+z\le151-I$&$\rho\ge436$, $\mathcal T\ge z$\\\bottomrule
\end{tabular}
\end{center}
Suppose $c\ge155$. Only the first band could survive. There $I\le18$, so $k\le3$ and $31\le r_s,r_t\le38$. The direct theorem applies to each core and gives $\rho\ge8\cdot69-224=328$. With $\mathcal T\ge114$, the budget gives $c\le145-I$, a contradiction. Therefore
\begin{equation}\label{eq:h-c154}
c\le154,\qquad j\le1.
\end{equation}
Every $78$-point ordered source satisfies $K=\card{D+D}\ge155$, by the two monotone endpoint chains. Since $\delta\le c<K$, neither leading resource can be empty. Thus both leading shifts are integral. This reduction directly removes every $j\ge2$; no interval catalogue is needed.

\subsection{Localization on the entire high hole domain}
Use the same source peel notation $a_v<b_v<c_v<e_v$ and put
\[
u=\sum_v(a_v+b_v),\quad C=c_s+c_t,\quad
J=\sum_v(e_v-c_v),\quad\mathcal H=h(P_s)+h(P_t).
\]
The mass identity and two-maximum envelope give
\[
2K=-263-w+I+\tau+c+4R+u+2C+J,
\qquad J+u\le599+w-I-\tau-c.
\]
The quantitative outer-layer row, for actual $\mu,\rho$, is
\[
2C+J\le C_0+u-336-2I-c-\mu-\rho-2z.
\]
Adding the bound for $J$ and using $C_0=1225+I+\tau$ yields
\[
2(e_s+e_t)\le1488+w-2I-2c-\mu-\rho-2z.
\]
The maximum-removal input gives $R\le73+\rho+\mathcal H$. Since $\mu\ge\rho$ and $w\le4z$, we obtain
\begin{equation}\label{eq:h-M}
M\le817+\mathcal H+z-I-c.
\end{equation}
The two structural kappa bounds give
\[
\mathcal H\le587-I-c-\rho+z,
\qquad\mathcal H\le441-I-c-\rho+4z.
\]
Together with the four band caps, these imply the following uniform bounds. All rows keep both physical size orientations.
\begin{center}\small
\begin{tabular}{@{}lrrrrr@{}}\toprule
$z$ range&$I$ cap&minimum side&$\mathcal H$ cap&$C_0-2\mathcal H$ floor&$M$ cap\\\midrule
$15$--$16$&$173$&$22$&$205$&$1320$&$1038$\\
$17$--$20$&$136$&$23$&$176$&$1394$&$1013$\\
$21$--$48$&$146$&$23$&$242$&$1374$&$1107$\\
$\ge49$&$102$&$25$&$302$&$1359$&$1270$\\\bottomrule
\end{tabular}
\end{center}
For example, the width floor is $1225+I+\mu-\mathcal H\ge1225+\rho-\mathcal H$. The last row also uses $z\le151$. In each row the individual core mass is greater than $\mathcal H$, since $T_{22}=253$, $T_{23}=276$, and $T_{25}=325$.

The common strip lemma gives $[\mathcal H,C_0-\mathcal H]_{\Z}\subseteq B_E$. Every one of the eight peeled entries is at most $M$, which is strictly less than the initial width in its row. Adjoining an entry $x$ joins this interval and its translate by $x$; their overlap or integer adjacency gives the enlarged interval with the same endpoint strip. Inducting through all eight entries proves
\[
[\mathcal H,S-\mathcal H]_{\Z}\subseteq B_D.
\]
Complement symmetry and inclusion $B_E\subseteq B_D$ on the common lower strip therefore prove
\begin{equation}\label{eq:h-localized}
z,w\text{ even},\qquad0\le w\le z,\qquad\eta_s=\eta_t=w/2.
\end{equation}

\subsection{Excluding zero and one empty unit}
At $j=0$, all own centers are active, so all outer coordinates are integral. Both leading resources are active by \eqref{eq:h-c154}; the four star leaves and both $O$ centers are retained. The raw gaps and Lemma~\ref{lem:two-caps} give
\[
4I+4\tau+c\le562+6w.
\]
But $\tau=\mu+\mathcal H\ge\rho+2z\ge300+2z$ and $w\le z$. The inequality would require $638+2z+4I+c\le0$, impossible.

Now let $j=1$. Each side has at least three active own centers; integral leading shifts determine the remaining outer magnitudes, so all are integral. If no $O$ is missing, both raw caps hold. Should $H_v$ be the missing singleton, retained $\own2_v<H_v$ supplies the lower gap guard. At most one star leaf is missing, so Lemma~\ref{lem:one-empty-leaf} applies. The two-cap inequality becomes
\[
4I+4\tau+\delta\le328+6w,
\]
which would require $872+2z+4I+\delta\le0$. Therefore the unique missing unit is some $O_v$.

The sharp source envelope at $j=1$ gives
\[
I+\tau+\delta\le359+w\le359+z.
\]
Using $\tau\ge\rho+2z$, we have $I+\delta+\rho+z\le359$. Every $z\ge17$ is impossible already by the next paired floor $\rho\ge345$; later floors are stronger. Evenness from \eqref{eq:h-localized} then leaves
\begin{equation}\label{eq:h1-domain}
z=16,\quad w\le16,\quad I+\delta\le43,\quad
28\le r_s,r_t\le41,\quad I+\tau+\delta\le375.
\end{equation}
If $\beta_v$ were active, the retained $H_v$ would supply the second cap on side $v$, while $O_o$ supplies it on side $o$. The same two-cap contradiction excludes this. Branch B has $\beta_v$ as a retained singleton and is therefore impossible. In A, $\beta_v$ is empty and the beta pair contributes on $u=o$. Every other singleton is retained.

The first-gap and final-form lemmas apply by \eqref{eq:h1-domain}. They give
\[
74\le3w/2\le24,
\]
again impossible. Together with \eqref{eq:h-c154}, this excludes the entire original high-hole family.


\section{Completion of the proof}\label{sec:conclusion}
The case division is by actual integer invariants, not by a representative list. The common partition gives precisely $Z_0,\ldots,Z_5$, $P_0,\ldots,P_5$, and $H$. The main-text arguments and the complete case appendices exclude every row:
\begin{center}\small
\begin{tabular}{@{}lll@{}}\toprule
Domain&Exhaustive alternatives&Contradiction mechanism\\\midrule
$Z_0,P_0$&$j=0$, $w\le56$&Pair1 has no actual contributor\\
$Z_1$&Pair1 empty or retained&Raw gaps; $72+I\le57$\\
$Z_2$&Complete necessary activity cover&Six retained rooted edges\\
$Z_3$&Empty beta or opposite missing $H$&All endpoint weight cases\\
$Z_4$&Complete A-E, A-H and B cover&Endpoint cards; ownerless final family\\
$Z_5$&A0, A1 and B&Retention templates and endpoint transport\\
$P_1$&Unique missing $O$ after cap reduction&$74\le3w/2\le21$\\
$P_2$&First gaps and final A forms&Rooted trim-six contradiction\\
$P_3$&First gaps and A0/A1 endpoint forms&Profile-free bridge and full endpoint cases\\
$P_4$&A0, A1 and B&One-hole source, exact credits and forests\\
$P_5$&All original placements&$7+z+I+\delta\le0$\\
$H$&$c\le154$, then $j=0$ or $1$&Two-cap obstruction; $74\le3w/2\le24$\\\bottomrule
\end{tabular}
\end{center}
The longer zero-hole derivations are in Appendices~\ref{sec:z-activity}--\ref{sec:z5}; the longer positive-hole derivations are in Appendices~\ref{sec:p2}--\ref{sec:p4}. Their classifiers and endpoint cards are part of the proof.

These alternatives preserve both physical orientations, the equality boundary in branch A, true whole-unit emptiness, and every zero-residual boundary. A choice of a contributing paired component is never counted as a second unit.

\begin{proof}[Proof of Theorem~\ref{thm:main}]
If a feasible set had more than $85$ elements, the retained upper bound $N\le86$ would make its size exactly $86$. The rational-kernel reduction produces an integral feasible set of that size containing zero. The stated structural theorem and original covering budget apply to it. The exhaustive partition places it in one of the rows just excluded, a contradiction. The construction in Section~\ref{sec:lower} proves the lower bound $71$.
\end{proof}
The conclusion is a mathematical proof relative to the expressly retained input theorems. It is not a claim that the entire chain has been formalized in a proof assistant, nor that finite tests replace its infinite arguments.



\section{The verification program and its exact scope}\label{app:verification}
The companion archive is \texttt{subset\_sums\_upper85\_verification\_20261005.zip}. It is independent of the original research folders. Extract it, open a terminal in the extracted program folder, and run one of the following commands with Python 3.10 or later:
\begin{quote}\ttfamily
python verify.py -{}-full\\
python verify.py -{}-quick
\end{quote}
On Windows, the supplied \texttt{run\_windows.bat} launches the full check; alternatively use \texttt{py -3 verify.py -{}-full}. The standard library is sufficient. No solver, package installation or network connection is required. The program writes \texttt{coverage.json} and \texttt{verification\_log.txt}; these identify every executed assertion and its scope. The quick mode explicitly skips the two large dense finite bases. Use full mode for the complete finite suite.

\subsection{What is actually checked}
The principal checks are as follows. Counts identify the stated finite domains; they are not counts of all feasible source configurations.
\begin{enumerate}
\item \textbf{The lower construction.} Exact subset-sum and fourfold-sum sets are formed for the stated $71$-element set. Its $3349$ subset sums are checked against its $3495$ fourfold sums, and all twelve interval rows are reconstructed.
\item \textbf{The fixed cover.} All $110$ signed fixed-label formulas and the $26$ disjoint negative-sum groups are reconstructed with integer coefficient vectors. This checks the fixed allowance $84$. The completeness of the geometric cover and half-line pairing is the written proof.
\item \textbf{Structural finite arithmetic.} The non-AP deleted-interval base consists of $2340$ sets, $910$ having positive holes. The structural inverse arithmetic includes $17$ inverse rows, four endpoint repairs, $13$ dense-diameter rows and $30$ oriented overlap alternatives. Four fixed positive-integer witnesses are expanded over $152$ systems: $72$ nonprimitive, $72$ half-step and eight size-$14$ boundary systems. The actual-set reductions to those systems remain written arguments.
\item \textbf{High-hole profile traces.} The four required thresholds $15,17,21,49$ and all sizes $15$ through $54$ give $160$ claims and $5200$ AP/no-AP branches. The checker verifies $95701$ coordinate increases and $1461470$ future-index conditions, using the old vector for every simultaneous update and the exact strict inequalities. These are supplied proof traces, not a search for stronger profiles.
\item \textbf{The two full-mode dense bases.} Every bounded-density boundary case for deficit $0$ through $9$ is generated: $227239$ cases, plus $1728$ small-size cases. Their lists overlap in $1078$ sets, so the union has $227889$ distinct sets. The restricted direct-defect calculation examines all $43758$ eight-deletion and $48620$ nine-deletion low patterns; $151$ and $1108$ pass the hole filter. It then reconstructs $91875$ complete sets in the three proved deletion families for sizes $14$ through $38$. Exact subset sums and distinct-pair defects are recomputed with independent representations.
\item \textbf{Common and terminal arithmetic.} Pointwise forest and missing-leaf memberships, exact quadratic minima, first-gap baseline inequalities, scalar partition ranges, and the explicitly listed endpoint-card recombinations and contradiction margins are checked with integers or rational numbers. These checks assume the row applicability, activity and retention proved in the relevant chapter. Deliberately corrupted data and strict-boundary controls must be rejected.
\end{enumerate}
The coverage map pairs these tasks with the report's mathematical identifiers. In particular, \texttt{G-FIRSTGAP} corresponds to Lemma~\ref{lem:first-gap-seven}, \texttt{G-ONEHOLE} to Lemma~\ref{lem:one-hole-terminal}, \texttt{G-STAR4} to Lemma~\ref{lem:one-empty-leaf}, and \texttt{G-TWOCAPS} to Lemma~\ref{lem:two-caps}. The exact list of terminal cards checked, including any additions beyond this summary, is in \texttt{COVERAGE\_MAP.json}; a row marked as a written assumption is not silently treated as verified.

\subsection{What remains a mathematical input or written proof}
A successful full run validates the finite domains just described. It does not by itself prove the earlier upper bound $86$, rationalization, the structural theorem's applicability to actual sets, the published additive-combinatorial inputs, the infinite non-AP reduction, bounded-density induction, scale-one arguments or complete activity classifications. Nor does it certify an arbitrary unlisted endpoint system merely because a similar one was checked. Those implications are supplied by the written proofs and the declared input theorems.

The package contains no global descriptor census, LP solver, endpoint-model search, profile discovery or claimed search over all real sets. A stored success field or a file hash is never accepted as the mathematical verification of a lemma. Hashes, where supplied in the source map, identify files only. The program reports explicit skips in quick mode and explicit failures rather than interpreting an unfinished check as success.



\clearpage
\appendix


\section{Common zero-hole source and activity reductions}\label{sec:z-activity}
This section proves the source and activity input used for $Z_2,Z_3,Z_4$.
It makes no positive-hole or first-gap assumption. Physical signs are never sorted by cardinality.
Write $U_v=\lambda_v+v_3$, $V_v=\lambda_v+v_4$ and retain the ordinary centre names $O_v,H_v,C_{4v},\ldots,C_{7v},\alpha_v,\beta_v$. Set
\[
 d_v=v_1-v_2,\ h_v=v_2-v_3,\ g_v=v_3-v_4,\ X_v=v_4-M_v,
 \quad X=X_s+X_t,\ G=g_s+g_t.
\]
Define $\tau_v=\sum P_v-T_{r_v}$ and $\tau=\tau_s+\tau_t$. Let $\xi_v=R_v-r_v$, so $0\le\xi_v\le\tau_v$ by the ordinary increasing-integer baseline, and put $\xi=R-69$, $A=I+\tau+\delta$.
The four source magnitudes outside the core on side $v$ are
$R_v+a_v,R_v+b_v,R_v+c_v,R_v+e_v=M_v$, with integral $1\le a_v<b_v<c_v<e_v$.
Put $L=e_s+e_t=M-R$, $U=a_s+b_s+a_t+b_t\ge6$, $C=c_s+c_t$, and $J=L-C$.

\subsection{One source calculation for the three rows}
The complete targets and original deficit give
\[
 F_v=[1,S_v]\cap\mathbb Z,\qquad S_s+S_t=2K+1488-78j-\delta,
 \qquad 2K=78j-263+A+4R+U+C+L.
\]
Delete both source maxima. The remaining 76-point sumset occupies a central interval of $2R+C+1$ positions plus at most 12 same-sign peeled-pair labels. Reinserting the maxima adds at most $2\cdot76+3=155$ labels. Thus
\[
 K\le2R+C+168,\qquad J+U\le599-78j-A.
\]
The full-source central/fringe envelope and Lemma~\ref{lem:common-offsets} give
\[
 K\le2R+L+1+B_s+B_t,\quad
 2B_v\le a_v+b_v+8,\quad 2B_v-(a_v+b_v+c_v)\le3.
\]
Comparison with the exact mass row yields
\begin{equation}\label{eq:z-common-source}
\begin{aligned}
 J&\ge A+78j-281, &L&\ge A+78j-271,\\
 A&\le437-78j, &K-\delta&\ge78j-123+I+\tau,\\
 K-M&\ge R+78j-425+A.&&
\end{aligned}
\end{equation}
For $j=2,3,4$ the cap on $A$ is respectively $281,203,125$.
Since $I\le A$, the physical ranges are respectively $18\!:\!51$, $21\!:\!48$, $24\!:\!45$.
The native losses of the two leadings sum to at most $\delta$; since $K>\delta$, both are active. Therefore $\lambda_s,\lambda_t$ are integral before any outer-integrality assertion.

\paragraph{Actual hits and raw gaps.}
An active ordinary centre $u$ is integral and satisfies $1-M_v\le u\le S_v+M_o$. Once the outer coordinates are integral, an empty centre satisfies $u\le-M_v$ or $u\ge S_v+1$, using respectively the actual source points $M_v$ and $0$.
Write $T_v=(O_v\text{ active})\lor(\beta_v,H_v\text{ active})$.
If $O_v$ is active, leading activity supplies the lower guard $H_v+M_v>\lambda_v+2M_v\ge1$, while $O_v-M_o\le S_v$ supplies the upper guard. Thus the first raw gap is licensed even when $H_v$ is empty, and before whole-outer integrality. The raw-55 gap count gives
\begin{equation}\label{eq:z-caps}
 T_v\Longrightarrow d_v-M\le56,\qquad
 \beta_v,O_v\text{ active}\Longrightarrow d_v-M\le57/2.
\end{equation}
For the second bound use the two disjoint intervals
$(\alpha_v+M_v,H_v-M_o)$ and $(H_v+M_v,O_v-M_o)$.
Beta activity implies $h_v-d_v\ge X_v+1$, so $\alpha_v>U_v$ and the leading maximum is below the first gap. These two gaps use one set of 55 fixed labels. The first bound uses just the appropriate one gap. Both arguments are valid even with real outer coordinates: a real open interval of length $q$ contains at least $q-1$ integers.
The leading maximum and the elementary four-summand high-tail count give
\begin{equation}\label{eq:z-tail}
 d_s+d_t\ge2K+G+2X-\delta,
 \qquad d_s+d_t+G+\delta-2K\ge806-78j.
\end{equation}
For the tail bound, a subset sum above $d_v+g_v$ after adding $v_1+v_2+v_3+v_4$ requires at least three occurrences from $\{v_1,v_2\}$ in its four-term representation. There are at most $4(78+6)+5=341$ such values per target.

\subsection{Restricted arms and empty vertices}\label{sec:separated-arms}
The following pointwise forest argument applies to any actual target, including a holed target. Only the subsequent source-span estimates are specialized to the current zero-hole domain.
For physical sign $v$ the candidate arms are
\[
 \mathcal L_v-V_v-U_v,\qquad \mathcal L_o-C_{7o}-C_{6o},
\]
with successive gaps $X_v,g_v$. Positive restricted supports on target $v$ lie at or above $\lambda_v+M_v-R_o$; negative ones lie at or below $\lambda_v-M_o+R_v$. These zones are disjoint since $M>R$.
Choose at most one arm in each zone, allowing at most one empty ordinary vertex per arm, and count only retained full resources. At target points an empty vertex contributes nothing and edge incidence is at most retained multiplicity minus one. Outside the target all incident edges lie in one arm, so at most one empty vertex is present and edge incidence is at most retained multiplicity. Consequently the restricted edge intersections are charged to the grouped deficiency of retained resources alone. No empty vertex is assigned capacity. This is the separated-arm extension of Lemma~\ref{lem:common-forest}.

Trimming $k$ largest core magnitudes on each sign yields
\[
 |Q_k|=70-2k,\quad\operatorname{span}Q_k\le69-2k+\frac{\tau}{k+1},\quad
 |Q_k\cap(Q_k+q)|\ge70-2k-\frac{\tau}{k+1}-q.
\]
Here $q$ must be integral. A suitable two-edge arm always exists for $j\le3$. For $j=4$ its failure would use all four holes $\{V_v,U_v,C_{7o},C_{6o}\}$ and leave both $O$ active; that possibility is excluded below before the arm is used. Combining a two-edge physical-$v$ arm, its unused $X_v$ star and an $X_o$ star gives
\begin{equation}\label{eq:z-four}2X_v+X_o+g_v\ge256-\tau-\delta.\end{equation}
Two two-edge arms and the unused $X_v$ star give, for $j=3,4$,
\begin{equation}\label{eq:z-five}X_o+2X_v+G\ge310-\tau-\delta.\end{equation}
Four full-core stars give $2X+4R+\delta\ge556$. These three statements require the indicated displacements to be integral, not any unproved activity of their ordinary leaves.

\subsection{Outer integrality and the no-two-cap argument}
A nonintegral $v_3$ forces singleton holes $U_v,C_{6o}$, and a nonintegral $v_4$ forces $V_v,C_{7o}$. All eight witnesses are distinct.
For $j=2$, at most one lower coordinate can be nonintegral. Its two witnesses exhaust the holes, but the retained $H_v,\alpha_v$ have nonintegral difference $\lambda_v+v_3+v_4$, a contradiction. With the lower coordinates integral, nonintegral $v_2$ would force the three holes $O_v,H_v,\alpha_v$.
For $j=3$, a sole nonintegral lower coordinate requires its two witnesses and one of $H_v,\alpha_v$. If $\alpha_v$ is the third hole, retained $H_v,O_v$ force $v_2,v_1$ integral. If $H_v$ is the third hole, retained $\alpha_v,O_v$ together with $C_{5o}$ in A or $\beta_v$ in B force $v_2$ integral. Either is impossible. After lower integrality, nonintegral $v_2$ forces $O_v,H_v,\alpha_v$ and additionally $C_{5o}$ in A or $\beta_v$ in B. Thus all outer coordinates are integral for $j=2,3$.

If both $T$ hold, set $Y=2M-2K+112+\delta$. Equations \eqref{eq:z-caps}--\eqref{eq:z-tail} give $G+2X\le Y$ and $G\ge806-78j-Y$.
For $j=4$, \eqref{eq:z-common-source} gives $Y\le338-2R-2I-2\tau-\delta\le200$, whereas the two preceding inequalities imply $2Y>494$. This excludes two caps before outer integrality.
For $j=2,3$, combine instead with the four-core-star row to obtain
\[
 234j-562+4I+4\tau+\delta\le0.
\]
This immediately contradicts $j=3$. For $j=2$ it gives $4I+4\tau+\delta\le94$. If a $C_{7v}$ is active, its identity $d_v=2v_4+o_4+g_v+C_{7v}$ and \eqref{eq:z-four} give $d_v-M\ge257-\tau-\delta$, contradicting the cap. If neither $C_7$ is active, they are exactly the two holes and both lower $C_6$ are active. Thus both empty $C_7$ are HIGH. A HIGH cross centre gives $\lambda_v>S_v+M_o$, and its leading loses at least $r_v+82\ge100$ points, contradicting $\delta\le94$. Hence the no-two-$T$ conclusion holds in all three rows.

It remains to justify outer integrality for $j=4$. Two nonintegral lower coordinates consume all four holes and leave both $O$ active. For a sole nonintegral lower coordinate on side $v$, let its residue be $\theta\ne0$ and the residue of $v_2$ be $b$. The opposite top coordinates are integral, since otherwise $O_o,H_o,\alpha_o$ add three holes to its two witnesses. The residues of $H_v,O_v,\alpha_v,\beta_v,C_{5o}$ are
\[
 b,\quad b+\theta,\quad b-\theta,\quad b-2\theta,\quad-b.
\]
If $b\ne0$, in A the two witnesses plus $H_v,C_{5o}$ leave both $O$ active. In B, either empty $O_v$ or empty $O_o$ would be the fourth hole and would leave $\alpha_v,\beta_v$ both retained, forcing $b=\theta=2\theta$. Therefore both $O$ are actually active in B also. This is impossible. Hence $b=0$, and the exact holes are the two witnesses, $O_v$, and $\alpha_v$.
If $v_4$ is the nonintegral coordinate, use a two-edge opposite arm, its unused star, and the aggregate edge $\mathcal L_v-U_v$ of integral gap $X_v+g_v$. Their four gaps give $X_v+g_v+2X_o+g_o\ge131+I$. Retained $O_o,C_{6o}$ then force $d_o-M\ge132+I>56$.
If $v_3$ is nonintegral, four integral-gap edges give $2X_o+X_v+g_o\ge131+I$. Active $\beta_o$ would contradict the double-gap cap using retained $C_{7o}$. In B it is retained. In A, if inactive, the retained beta pair supplies active $\beta_v$, and retained $H_v,O_o$ give two $T$ predicates. Finally, after lower integrality a nonintegral top pair forces $O_v,H_v,\alpha_v$ and $C_{5o}$ in A or $\beta_v$ in B. Every lower singleton is retained, so \eqref{eq:z-five} and active $O_o,C_{6o}$ force $d_o-M\ge186+I>56$. This completes integrality, including the branch-B residue case.

\subsection{Activity implications and owner grammar}
For $j=3,4$, equations \eqref{eq:z-five} and the identity
$d_v=o_4+2v_4+G+C_{6v}$ show that active $C_{6v}$ gives
$d_v-M\ge108+I$ or $186+I$. Thus $O_v,C_{6v}$ cannot both be active, nor can $\beta_v,H_v,C_{6v}$ all be active.
Similarly active $C_{7v}$ gives $d_v-M\ge54+I$ or $132+I$ by \eqref{eq:z-four}, which excludes simultaneous $\beta_v,O_v,C_{7v}$.
If a cross centre is empty but a lower cross centre is active, it is HIGH. Its leading loss is at least $r_v+82$, and each retained higher singleton loses at least $r_v+5$. Thus $C_{5v}\Rightarrow C_{6v}$ follows from
\[
 j=3:\ \delta\ge4r_v+97,\ I+4r_v\ge134;
 \qquad j=4:\ \delta\ge3r_v+92\ge164>125.
\]
For $j=4$ the same $3r_v+92$ estimate proves $C_{4v}\Rightarrow C_{5v}$, and $2r_v+87\ge135>125$ proves $C_{6v}\Rightarrow C_{7v}$. For $j=2$ in A, empty $C_{5v}$ above active $C_{4v}$ leaves at least five higher retained singletons, giving $\delta\ge6r_v+107$; this contradicts $A\le281$ because $I+6r_v\ge198$.

Finally put $W=v_3+v_4+o_3+o_4>2M$. The identities
\[
 C_{4s}+C_{4t}=-(s_2+s_3+s_4+t_2+t_3+t_4),\quad
 \beta_s+\beta_t+C_{5s}+C_{5t}=-2W,
\]
\[
 C_{4v}+\beta_v+C_{5o}=-2v_3-2v_4-o_2-o_3-o_4,\quad
 C_{4v}+C_{5o}=-(\beta_s+\beta_t)-2W-d_o
\]
compared with the actual lower activity endpoints prove sole ownership of retained pair1, prohibit the indicated triple, and in B prove sole, matching ownership of the two retained pairs. They also prohibit all four beta/$C_5$ components being active. These are actual-configuration conclusions, so later cards may charge a pair through its sole owner's native resource.


\section{Two missing units: the zero-hole case $Z_2$}\label{sec:Z2}
Throughout this chapter physical signs are $s,t$; $o$ means the sign opposite $v$.
We write $U_v=\lambda_v+v_3$, $V_v=\lambda_v+v_4$ for the ordinary centres otherwise called own2 and own3, and $P=\{C_{4s},C_{4t}\}$ for pair1. The symbol $\xi_v=R_v-r_v$ and $\xi=R-69$ denotes radial excess; it is never the structural defect $\rho$.
The leading full resource is $\mathcal L_v=\lambda_v+(D_v+D_v)$, with native loss $\ell_v=K-|\mathcal L_v\cap F_v|$.
Set $X_v=v_4-M_v$, $g_v=v_3-v_4$, $h_v=v_2-v_3$, $d_v=v_1-v_2$, $X=X_s+X_t$, and $G=g_s+g_t$.

\subsection{Source rows and activity}
Here $18\le r_v\le51$, $r_s+r_t=69$, $F_v=[1,S_v]\cap\mathbb Z$, and the source rows are
\begin{equation}\label{eq:z2-source}
 A=I+\tau+\delta\le281,\quad K-\delta\ge33+I+\tau,
 \qquad K-M\ge R-269+I+\tau+\delta .
\end{equation}
The zero-hole activity lemma supplies integral outer coordinates, active full leadings, and the following elementary consequences of the original cover:
\[
 u\text{ active}\Longrightarrow 1-M_v\le u\le S_v+M_o,
 \qquad u\text{ empty}\Longrightarrow u\le-M_v\ \text{or}\ u\ge S_v+1.
\]
Put $T_v=(O_v\text{ active})\lor(\beta_v,H_v\text{ active})$.
Then $T_s,T_t$ cannot both hold; $T_v$ gives $d_v-M\le56$.
The literal identities imply that $P$ has a sole actual owner when retained, that active $\beta_v,C_{5o}$ forbid owner $v$ of $P$, and that in branch B the two retained pairs have the same sole owner.
The leading maximum gives
\begin{equation}\label{eq:z2-lead}
 d_v-M\ge K-M+2X_v+g_v-\ell_v.
\end{equation}
These activity assertions refer to actual hits, never to a syntactically selected paired witness.

For clarity the exact necessary two-hole cover is displayed here. Every unit outside the hole set is retained. In A put $B=\{\beta_s,\beta_t\}$; in B put $C=\{C_{5s},C_{5t}\}$.
\begin{center}\small
\begin{tabular}{lll}
Type&Whole holes&Actual owners\\\hline
A-OO&$O_s,O_t$&$B:b$, $P:o(b)$\\
A-H&$O_v,H_v$&$B:b$, $P:o(b)$\\
A-empty-$B$&$O_v,B$&$P:a$\\
A-empty-$P$&$O_v,P$&$B:o$\\
A-$C_5$&$O_v,C_{5v}$&$B:o$, $P:o$\\
A-other&$O_v,W$, $W\in\mathcal W_v$&$B:o$, $P:v$\\
B-H&$O_v,H_v$&$P,C:a$\\
B-$\beta$&$O_v,\beta_v$&$P,C:a$
\end{tabular}
\end{center}
Here $\mathcal W_v=\{H_o,U_s,U_t,V_s,V_t,C_{5o},C_{6s},C_{6t},C_{7s},C_{7t},\alpha_s,\alpha_t\}$.
Indeed at least one $O$ is absent. If both are absent, both $H$ are present and the no-two-$T$ rule forces A with a sole beta owner. If only $O_v$ is absent, $T_o$ forces either $H_v$ absent or $\beta_v$ inactive. The second hole is then precisely one of the displayed paired or singleton alternatives. The owner restrictions follow from the literal identities and the cross-activity implication $C_{4v}\Rightarrow C_{5v}$. This also rules out a genuinely split beta pair.

\subsection{Rank overlap and one original budget}
Delete the largest $k$ core magnitudes on each sign, retaining zero, to obtain $Q_k$. Monotonicity of the core rank deviations gives
\[
 |Q_k|=70-2k,\qquad \operatorname{span}Q_k\le69-2k+\frac{\tau-\xi}{k},\qquad
 |Q_k\cap(Q_k+q)|\ge f_k-q,
 \quad f_k=70-2k-\frac{\tau-\xi}{k}.
\]
In particular $f=f_4=(179-\tau+R)/4$.
The two arms associated to physical sign $v$ are
\[
 \mathcal L_v-V_v-U_v,\qquad \mathcal L_o-C_{7o}-C_{6o},
\]
with gaps $X_v,g_v$; the shortcuts have gap $X_v+g_v$.
Apply the pointwise forest lemma only to retained full resources. Designated leading-root losses and native losses of unused distinct singletons can be included in the same original deficit:
\begin{equation}\label{eq:z2-rooted}
 \sum_{e}|A_e\cap A'_e|+\sum_{\text{designated roots}}\ell_v+
 \sum_{\text{unused }Z}\Gamma_Z\le\delta.
\end{equation}
A full leading is counted once even when two restricted copies certify different edges. From \eqref{eq:z2-source},
\begin{equation}\label{eq:z2-fbounds}
 K-M+3f-\delta\ge-14+I+\tau/4,\qquad
 K-M+2f-\delta\ge-76+I+\tau/2.
\end{equation}

\subsection{No simultaneous beta and $O$ hit}
Suppose $\beta_v,O_v$ are active. Then $h_v-d_v\ge X_v+1$ and $\alpha_v>U_v$.
The two disjoint open gaps
\[
 (\alpha_v+M_v,H_v-M_o),\qquad(H_v+M_v,O_v-M_o)
\]
lie within the complete target whenever nonempty and contain no nonfixed contribution. The same 55 raw fixed labels cover both. Therefore
\begin{equation}\label{eq:z2-doublegap}d_v-M\le57/2.\end{equation}
The no-two-$T$ rule makes $O_o$ a hole. Write $W$ for the one other whole hole. The following four cards exhaust it.
\begin{enumerate}
\item $W\notin\{V_v,C_{7v},C_{7o}\}$. Choose a complete two-edge physical-$v$ arm and the opposite one-edge star; one of $U_v,C_{6o}$ is retained. The three gaps total $2X_v+g_v$, and designate $\mathcal L_v$ as root. The unused retained $C_{7v}$ satisfies $C_{7v}\le-M_v+24$: this follows by inserting \eqref{eq:z2-doublegap} in $C_{7v}=d_v-2v_4-o_4-g_v$. Its core translate has at most 20 positive points, and all four negative source peels miss. Thus $\Gamma_{C_{7v}}\ge54$. Equations \eqref{eq:z2-lead}--\eqref{eq:z2-fbounds} imply $d_v-M\ge40+I+\tau/4$.
\item $W=V_v$. Leading activity makes $V_v>-M_v$, hence its emptiness makes it HIGH. Then $\alpha_v,H_v,O_v$ are HIGH retained resources, each losing at least $r_v+5$ native points. Use the shortcut $\mathcal L_v-U_v$ and $\mathcal L_o-C_{7o}$, root $\mathcal L_v$, and these three unused losses. This gives
\[
 d_v-M\ge-61+I+3r_v+\tau/2\ge40,
 \quad I+3r=(r-33)^2+101.
\]
\item $W=C_{7v}$. The lower retained $C_{6v}$ makes $C_{7v}$ HIGH; consequently $U_v,H_v,O_v$ each lose at least $r_v+5$. Use $\mathcal L_v-V_v$ and the full negative arm on $o$, root $\mathcal L_v$. The result is $d_v-M\ge1+I+3r_v+\tau/4\ge102$.
\item $W=C_{7o}$. Its lower retained $C_{6o}$ makes it HIGH. Thus $\lambda_o>S_o+M_v$, and $\mathcal L_o$ contains the entirely HIGH set $\lambda_o+[D_o\cup(M_o+D_o)]$ of at least $r_o+82$ points. Moreover $V_o,U_o,H_o$ are retained HIGH resources, each losing at least $r_o+5$. Use $\mathcal L_v-V_v$ and the shortcut $\mathcal L_o-C_{6o}$, with both leading roots designated. Then
\[
 d_v-M\ge21+I+4r_o+\tau/2\ge155,
 \quad I+4r=(r-32)(r-33)+134\ge134.
\]
\end{enumerate}
Every lower bound contradicts \eqref{eq:z2-doublegap}. Paired resources are not charged in any of the four cards.

\subsection{The final six edges}
The preceding exclusion removes A-empty-$P$, A-$C_5$, A-other, B-H and B-$\beta$, and removes the beta-owner-$o$ alternative of A-H. Precisely three forms remain:
\[
 \{O_v,B\};\qquad \{O_s,O_t\}\ (B:b,P:o(b));\qquad
 \{O_v,H_v\}\ (B:v,P:o).
\]
All their $\alpha,C_5,U,V,C_7$ singletons are retained. At least one $T_u$ holds: use the retained opposite $O$ in the first and third forms, and $\beta_b,H_b$ in the second. The identity
\[
 \alpha_o+C_{5u}=d_u-2(M+X)-G
\]
and activity give $2X+G\le54$. On each target use the retained tree $\mathcal L_v-V_v-U_v$ together with $\mathcal L_v-C_{7v}$, now with $Q_6$. Its six edges across the two targets have total gap $2X+G$. Hence
\[
 2X+G\ge6f_6-\delta=348-\tau+\xi-\delta,
 \qquad \tau+\delta\ge294+\xi.
\]
Together with $I+\tau+\delta\le281$ this gives $13+I+\xi\le0$, impossible. Thus $Z_2$ is excluded in both branches, all actual owner choices, and both physical size orientations.


\section{Three missing units: $Z_3$}\label{sec:Z3}
Here $21\le r_v\le48$, $A\le203$, $L\ge A-37$, $J\ge A-47$, and
\begin{equation}\label{eq:z3-source}
 S_s+S_t=2K+1254-\delta,\quad
 2K=-29+A+4R+U+C+L,\quad J+U\le365-A.
\end{equation}
All coordinates, leadings and activity implications used below were proved in Section~\ref{sec:z-activity}.

\subsection{Exhaustive retention reduction}
At least one $O$ is absent. If both are absent there is one further hole. A $C_6$ hole in A would force the same $C_5$ singleton hole, exceeding the budget. Thus both $C_6$ are active. An empty beta pair gives exactly
\[
 \text{A0}:\quad\{B,O_s,O_t\},\qquad P\text{ sole-owned by }s
\]
after reflecting the actual owner. If beta is retained, its active component $\beta_s$ forces $H_s$ absent. Opposite $H_t,C_{6t}$ are active, so $\beta_t$ is inactive, and active $C_{5t}$ forces pair1 owner $t$. This gives exactly
\[
 \text{A1}:\quad\{O_s,O_t,H_s\},\qquad B:s, P:t.
\]
In B, both $O$ absent would require a further hole in each of the disjoint triples $\{\beta_v,H_v,C_{6v}\}$, which is impossible.
If only $O_o$ is absent and $O_v$ active, then $C_{6v}$ and hence $C_{5v}$ are inactive. In A these are the other two holes. Beta is then sole on $v$, while $\beta_v,O_v,C_{7v}$ are all active, impossible. In B there is one further hole $Z\in\{\beta_o,H_o,C_{6o}\}$. Both pairs are retained and sole on $o$; $Z=C_{6o}$ contradicts active $C_{5o}$, while either other choice leaves $\beta_v,O_v,C_{7v}$ active. This exhausts B and proves the two displayed A forms are exhaustive. Every unlisted whole unit in either form is retained.

\subsection{Source deletion and endpoint tools}
Delete only $M_s$. The residual 77-point source has central sum interval of $2R+c_s+e_t+1$ positions and at most $6+10$ exceptional pair labels. Reinsertion adds at most 78 values, hence
\[
 K\le2R+c_s+e_t+95.
\]
Comparing its two orientations with \eqref{eq:z3-source} yields, for $j_v=e_v-c_v$,
\[
 |j_s-j_t|\le219-A-U,\quad j_v\ge A-133+U/2,\quad
 e_v\ge\max(4,A-127).
\]
Deleting the largest two source values on $s$ instead gives
$K\le2R+b_s+e_t+169$: there are 76 remaining values, $3+10$ exceptional pairs and $2\cdot76+3=155$ reinsertion labels. Comparison with \eqref{eq:z3-source} and $J+U\le365-A$ gives
\begin{equation}\label{eq:z3-side}e_s\le366-A-a_s-a_t-b_t\le362-A.\end{equation}
These rows use only source data and hold in either physical orientation.

\begin{lemma}[Complete-target rank endpoints]\label{lem:z-rank-endpoint}
For an active native ordinary resource $Z+D_v$ with complete target $[1,S_v]_{\Z}$, let $\Gamma_Z=l_Z+h_Z\le77$ be its lower and upper clips. Then
\[
 Z\ge74-r_v-l_Z-b_Z,\qquad
 Z\le S_v+r_o-73+h_Z+b_Z,
\]
The symbol $b_Z$ is chosen separately for the selected lower or upper endpoint; the lower error uses source side $v$ and the upper error uses side $o$.
A normal endpoint has error zero. At an outside-hull endpoint with selected clip $k<74$,
\[
 (74-k)b_Z\le\tau_{\mathrm{source}},\qquad b_Z\le\xi_{\mathrm{source}}.
\]
At selected clip $k=74,\ldots,77$,
\begin{equation}\label{eq:z3-tail-error}
 b_Z\le\xi_{\mathrm{source}}+e_{\mathrm{source}}-4.
\end{equation}
\end{lemma}
\begin{proof}
For a lower endpoint with $Z\ge1$, at most $Z-1$ negative source values can give positive images below $Z$. Counting the $r_o+4$ negative source values therefore gives $Z\ge74-r_v-l_Z$, with zero error. If $Z\le0$, write $l_Z=74-r_v+j$. Activity supplies a next positive source magnitude $y_{j+1}$ and
\[
Z\ge1-y_{j+1}=74-r_v-l_Z-b_Z,\qquad b_Z=y_{j+1}-(j+1).
\]
If $l_Z<74$, that next rank is in the core. The last $r_v-j=74-l_Z$ core deviations are at least $b_Z$, proving both bounds for a light error. If $l_Z\ge74$, the next rank is the $k=l_Z-73$th peel, so
\[
b_Z=\xi_v+\operatorname{offset}_k-k\le\xi_v+e_v-4
\]
by strict integral ordering of the four offsets. Apply the same lower argument to $S_v+1-Z$ and the reflected source $-D_v$ for the upper endpoint. Completeness of the target gives $\Gamma_Z=l_Z+h_Z$; no interior-hole loss is omitted.
\end{proof}
This generic complete-target argument is independent of the specialized positive-hole endpoint lemma. It retains every peeled rank and every outside-hull error.

\subsection{A0: the five weighted endpoints}
Here pair1 is sole on $s$. Both empty $O$ lie above retained lower centres and hence are HIGH. The eight distinct fixed labels $O_v+v_i$ ($v=s,t$, $i=1,\ldots,4$) exceed their targets and are disjoint from the 26 nonpositive baseline witnesses. Therefore the fixed credit is $\eta=8$ and every compatible grouping has deficiency at most
\[
 \varepsilon=\delta-8,\qquad \tau+\varepsilon\le195-I.
\]
Use lower endpoints $C_{4s},\alpha_s,\alpha_t,C_{5t}$ and upper endpoint $H_s$, with respective weights $1,2,1,1,1$. The five resources are distinct; the weight two is algebraic only. The error source of $C_{4s},\alpha_s$ is $s$; the other three use $t$. Call $k\le72$ light and $k\ge73$ heavy in this card. A light error satisfies both $b_Z\le\xi_v$ and $b_Z\le k$: an outside-hull clip has $26\le k\le72$, and
\[
 k(75-k)-216=(k-3)(72-k)\ge0,
 \qquad \tau_v\le195-k\le k(74-k).
\]
For every heavy clip, including the core-maximum rank $k=73$, the error bound \eqref{eq:z3-tail-error} remains valid because $e_v\ge4$. The selected clips sum to at most $\varepsilon\le195$, so there are at most two heavy endpoints.

Trim four core ranks per sign and set $f=(179-\tau+R)/4$. On each target use the retained tree $\mathcal L_v-V_v-U_v$ and $\mathcal L_v-C_{7v}$. Across targets the six gaps total $2X+G$. Its grouped deficiency $D_F$ satisfies
\[
 D_F\ge6f-(2X+G),\qquad N+D_F\le\varepsilon,
\]
where $N$ is the sum of the five endpoint native deficiencies. No endpoint is a forest vertex. The literal identity
\[
 H_s=C_{4s}+2\alpha_s+\alpha_t+C_{5t}+4(M+X)+2G
\]
gives $S_s\ge236+4(M+X)+2G-E-B$, where $E,B$ are the weighted clips and errors. The envelope $K\le M+R+21$ then gives
\[
 S_t+2M+12f\le2R+1060-\delta+E+2D_F+B.
\]
Keep the endpoint rebates within the same budget:
\[
 E+2D_F+B\le2\varepsilon+2b_{\alpha_s}
       +\sum_{Z\ne\alpha_s}(b_Z-k_Z).
\]
Light summands in this sum are nonpositive. Write
$B_* =2b_{\alpha_s}+\sum_{Z\ne\alpha_s,\, Z\text{ heavy}}(b_Z-k_Z)$.
With $r=r_t$, use
$S_t+2M\ge T_r+4r+\tau_t+10+2R+2L+4\xi_t$.
The exact master row is
\begin{equation}\label{eq:z3a0master}
 T_r+4r+\tau_t+2L-3\tau+3\xi_s+7\xi_t
 \le290+\delta+B_*.
\end{equation}
The integer minimum used in every case is
\begin{equation}\label{eq:z3-min3}
 Q_3(r):=T_r+4r+3I=641+\frac{(r-29)(7(r-29)+1)}2\ge641.
\end{equation}

The following six cases are exhaustive. In the table $c\ge73$ denotes another heavy clip; $H\ge146$ denotes the sum of two heavy clips. Radial error coefficients are absorbed componentwise by $3\xi_s+7\xi_t$. This is legitimate even for the two source-$s$ endpoints: there is only one such endpoint other than $\alpha_s$, namely $C_{4s}$.
\begin{center}\small
\begin{tabular}{p{.31\linewidth}p{.42\linewidth}r}
Endpoint case&Resulting upper inequality&Lower bound\\\hline
$\alpha_s$ light, none heavy&$Q_3+\tau_t+2\delta\le567$&$657$\\
$\alpha_s$ light, one heavy&$Q_3+\tau_t+2\delta+c\le725$&$876$\\
$\alpha_s$ light, two heavy&$Q_3+\tau_t+2\delta+H\le891$&$1095$\\
$\alpha_s$ heavy, no other&$Q_3+\tau_t+2\delta\le739$&$803$\\
$\alpha_s$ heavy, other on $t$&$Q_3+\tau_t+2\delta+c\le811$&$1022$\\
$\alpha_s$ heavy, other $C_{4s}$&$Q_3+\tau_t+2\delta+c\le894$&$1022$
\end{tabular}
\end{center}
Here are the offset substitutions, so the table does not conceal a tail estimate. In the first three rows use respectively $2L\ge2A-74$, $2L-e_v\ge L+4\ge A-33$, and $2L-e_v-e_w\ge0$, followed by $\tau+\delta\le203-I$. In the fourth and fifth rows cancel the weighted $\alpha_s$ error, then use respectively $2e_t\ge2A-254$ and $e_t\ge A-127$. In the last row the retained offset is $2e_t-e_s\ge3A-616$, by \eqref{eq:z3-side}. A heavy clip forces $\delta\ge8+73=81$ and two force $\delta\ge154$; no heavy clip forces $\delta\ge8$. These are the displayed lower bounds. Each strictly exceeds its upper bound, proving A0 impossible.

\subsection{A1: opposite native owners}
Now $B$ is sole on $s$, $P$ sole on $t$, and $O_s,O_t,H_s$ are the holes. The eight $O$ labels and three $H_s+s_i$ labels ($i=2,3,4$) are HIGH. A twelfth is $\alpha_s+s_2$: beta activity gives $s_2>d_s$, and $\alpha_s=H_s-d_s$. Thus
\[
 \eta=12,\quad\varepsilon=\delta-12,\quad\tau+\varepsilon\le191-I.
\]
The endpoints are lower $\beta_s,C_{4t}$ and upper $H_t$. Put
$B=\Gamma_{\beta_s}+\Gamma_{C_{4t}}$ and $H=\Gamma_{H_t}$, so $B+H\le\varepsilon$.
For this card light means $k<74$ and heavy $k\ge74$; at most two are heavy. Every light error is at most $2\Gamma_Z$, including outside-hull centres: on $26\le k\le73$, $k(149-2k)\ge219>192$, while $\tau_v\le192-k$. Heavy errors are \eqref{eq:z3-tail-error}; the error source is $s$ for $\beta_s,H_t$ and $t$ for $C_{4t}$.

Use $f=\max(f_0,f_2,f_6)$, where
\[
 f_0=139-R,
 \quad f_k=70-2k-\frac{\tau-\xi}{k}\quad(k=2,6).
\]
Each forest uses an actual trimmed set attaining this maximum. In particular
\begin{equation}\label{eq:z3-uniform}
 f\ge(202-\tau)/3,\qquad4\tau+12f\ge808.
\end{equation}
The ten-edge forest consists of
\[
\begin{array}{ll}
s:\ \mathcal L_s-V_s-U_s,&\mathcal L_s-C_{7s}-C_{6s}-C_{5s},\\
t:\ \mathcal L_t-V_t-U_t-H_t,&\mathcal L_t-C_{7t}-C_{6t}.
\end{array}
\]
Its gap sum is $2(X_s+X_t+g_s+g_t+h_t)$ and it omits both paired endpoint resources. The four-edge physical-$s$ double arm, and its two first edges, omit both pairs and $H_t$. Their separate budget inequalities are
\[
 X_s+X_t+g_s+g_t+h_t\ge5f-(\varepsilon-B)/2,
\]
\[
 X_s+g_s\ge2f-(\varepsilon-B-H)/2,
 \qquad X_s\ge f-(\varepsilon-B-H)/2.
\]
Add the first, twice the second and three times the third:
\[
 6X_s+X_t+3g_s+g_t+h_t\ge12f-3\varepsilon+3B+\tfrac52H.
\]
The endpoint identity is
\[
 H_t=2d_s+\beta_s+C_{4t}+2s_4+t_4+g_s+g_t+h_t.
\]
Use $d_s\ge K+g_s+2X_s-\ell_s$ twice, and retain the native rebate
$E_{\mathrm{clip}}+2\ell_s\le2\varepsilon-(B+H)$.
After source-mass substitution the master row is
\begin{equation}\label{eq:z3a1master}
 T_r+4r+5\xi_s+\tau_s+R+L+2e_s+12f+60+4B+\tfrac72H
 \le1096+4\delta+E_{\mathrm{err}},\quad r=r_s.
\end{equation}
Here $Q_4(r)=T_r+4r+4I=665+(r-30)(9r-273)/2\ge665$ for integral $r$.
Subtract all light errors using $2\Gamma_Z$; the remaining coefficients of every light native loss are nonnegative. A heavy beta or pair1 endpoint contributes at least $4\cdot74=296$, and a heavy $H_t$ endpoint at least $7\cdot74/2=259$. Cancel its full error, including the $e_v-4$ term. The exhaustive result is
\begin{center}\small
\begin{tabular}{p{.35\linewidth}p{.38\linewidth}r}
Case&Contradictory bound&Minimum\\\hline
One heavy, source $s$&$Q_4+4\xi_s+\tau_s+4\tau+R+12f\le1415$&$1542$\\
One heavy, source $t$&$Q_4+6\xi_s+\tau_s+4\tau+12f\le1467$&$1473$\\
Two heavy, both source $s$&$Q_4+3\xi_s+\tau_s+4\tau+R+12f\le1119$&$1542$\\
Two heavy, different sources&$Q_4+5\xi_s+\tau_s+4\tau+12f\le1208$&$1473$\\
All light, $A\le131$&$Q_4+5\xi_s+\tau_s+4\tau+R+12f\le1458$&$1542$\\
All light, $A\ge131$&$Q_4+5\xi_s+\tau_s+4\tau+R+12f\le1530$&$1542$
\end{tabular}
\end{center}
For the first heavy row use $L+e_s\ge L+4$ and $L\ge A-37$; for the second, $L+2e_s-e_t=3e_s\ge12$ and $R-\xi_t=69+\xi_s$. For the two-heavy rows the remaining offsets are respectively $L$ and $2e_s\ge8$. Then use $\delta\le203-I-\tau$.
In the all-light rows, errors remove at most $2(B+H)$, leaving $2B+3H/2\ge0$ in \eqref{eq:z3a1master}. If $A\le131$, use $e_s\ge4$, $L\ge A-37$ and $\delta\le131-I-\tau$. If $A\ge131$, use $e_s\ge A-127$, $L\ge A-37$ and $\delta\le203-I-\tau$. The shared boundary $A=131$ is included. The minima follow from $Q_4\ge665$, $R\ge69$ and \eqref{eq:z3-uniform}. Thus A1 is impossible, completing $Z_3$ with every endpoint position and both physical owner orientations.


\section{Four missing units: $Z_4$}\label{sec:Z4}
The physical ranges are $24\le r_v\le45$. The source interface proved above is
\begin{equation}\label{eq:z4-source}
 A\le125,\quad L\ge A+41,\quad J\ge A+31,\quad J+U\le287-A,
 \quad S_s+S_t=2K+1176-\delta.
\end{equation}
In particular all outer coordinates are integral, both leadings are active, and
$K-M\ge R-113+A$.

\subsection{Both $O$ units are absent}
The four-edge bound gives $2X_v+X_o+g_v\ge131+I$. Thus activity of $C_{7v}$ implies $d_v-M\ge132+I$. Comparison with the raw caps proves
\[
 O_v,C_{7v}\text{ cannot both hit};\qquad
 \beta_v,H_v,C_{7v}\text{ cannot all hit}.
\]
We first exclude simultaneous $\beta_v,O_v$ activity. The opposite $O_o$ is empty and lies above a retained one of $H_o,U_o,V_o,C_{7o}$, because four such singletons cannot all be among the remaining three holes. Hence $O_o\ge S_o+1$. Its four fixed labels $O_o+o_i$ are HIGH, giving fixed credit $\eta=4$ and budget $\varepsilon=\delta-4$.
Choose a two-edge physical-$v$ arm and its unused $X_v$ star, allowing at most one empty vertex per arm as proved above. With the rank-three trim the three-edge bound is
\[
 2X_v+g_v\ge192-3\tau/4-D_F,\qquad D_F\le\varepsilon.
\]
The separate leading-loss bound is $\ell_v\le\varepsilon$. Using two copies, not one copy, of this remaining budget gives
\[
 d_v-M\ge K-M+192-3\tau/4-D_F-\ell_v
 \ge R+79+I+\tau/4-\delta+2\eta.
\]
The double-gap cap then forces $R+2I+5\tau/4+2\eta\le149/2$, but its left side is at least $77$. Thus $\beta_v,O_v$ cannot both hit.

If $O_v$ were active, then $O_o,C_{7v},C_{6v}$ are holes, and $C_{5v},\beta_v$ are inactive. In A the fourth hole is $C_{5v}$, leaving the retained beta pair sole on $o$; in B it is the singleton $\beta_v$. Either way $\beta_o,H_o,C_{7o}$ all hit, a contradiction. Therefore both $O$ are absent. Their eight HIGH labels give a common fixed credit $\eta=8$ for the remaining configurations.

\subsection{The exact retention table}
Put $\mathcal N=\{U_s,U_t,V_s,V_t,\alpha_s,\alpha_t\}$ and
$\mathcal W=\mathcal N\cup\{P,H_s,H_t,C_{5s},C_{5t}\}$.
Every unlisted unit in the following table is retained; $o$ is opposite the named $v$ or $r$.
\begin{center}\small
\begin{tabular}{p{.13\linewidth}p{.40\linewidth}p{.36\linewidth}}
Type&Exact whole holes&Actual pair status\\\hline
A-E&$\{B,O_s,O_t,Z\}$, $Z\in\mathcal W$&$P$ empty if $Z=P$; sole owner $o$ if $Z=C_{5v}$; either owner otherwise\\
A-H&$\{O_s,O_t,H_v,Z\}$, $Z\in\mathcal W\setminus\{H_v\}$&$B:v$; $P$ empty if $Z=P$, owner $v$ if $Z=C_{5o}$, owner $o$ otherwise\\
B&$\{O_s,O_t,X_r,Y_o\}$, $X_r\in\{\beta_r,H_r\}$, $Y_o\in\{\beta_o,H_o\}$&Both pairs retained, sole owner $r$
\end{tabular}
\end{center}
Here is a direct exhaustion. If beta is empty in A, one further hole $Z$ remains. A $C_6$ hole requires the same $C_5$ hole; a $C_7$ hole requires a $C_6$ hole. Thus $Z\in\mathcal W$. If beta hits on $v$, then either $H_v$ or $C_{6v}$ is absent. The latter forces $C_{5v}$ absent, leaving active $\beta_v,H_v,C_{7v}$, impossible. Thus $H_v$ is absent and the final hole is $Z\in\mathcal W\setminus\{H_v\}$. Beta cannot hit on both sides: the only additional choice would be $Z=H_o$, leaving both $C_5$ active, contrary to the four-component identity. The displayed pair1 owner rules follow from the cross chain and the forbidden triple $C_{4v},\beta_v,C_{5o}$.
In B each disjoint triple $\{\beta_v,H_v,C_{6v}\}$ needs a hole. The two remaining holes therefore occupy these triples and both pairs are retained. Their common owner $r$ has active $C_{5r},C_{6r}$; the opposite $C_6$-hole choice leaves $\beta_o,H_o,C_{7o}$ active. This proves exactly the B row. Both physical owner choices and equality in branch A are included.

\paragraph{Excluding two beta holes in B.}
An empty $\beta_v$ opposite active $C_{4o}$ must be LOW. To see this, assume $\beta_v\ge S_v+1$. The three-edge bound without fixed credit is
$2X_v+g_v\ge192-3\tau/4-\delta$.
The opposite leading's lowest possible hit and $C_{4o}$ activity give
$v_1\le S_o+2M_v+M_o-K+\ell_o$.
The HIGH beta gives $v_1\ge S_v+1+2d_v+2M_v+2X_v+g_v$.
Use $d_v\ge K+2X_v+g_v-\ell_v$, $\ell_s+\ell_t\le\delta$, and the exact mass row to obtain
\[
 2S_v+K+3g_v+6X_v\le M_o+1175+\delta.
\]
Insert $K-M\ge R-113+A$, $S_v\ge T_r+\tau_v+4r+10$, $M_v\ge r+4$, and the three-edge bound. The result is
\[
 2T_r+9r+4I+7\tau/4+2\tau_v\le994,
 \qquad 2T_r+9r+4I=1223+(r-27)(5(r-27)+4)\ge1223,
\]
a contradiction. For the owner-side empty $\beta_r$, retained $H_r$ would give $2d_r\le M_o-1$ if beta were HIGH; active $C_{4r}$ instead gives $d_r>M$, so this beta is LOW also. If both beta singletons were holes, their sum would be at most $-M<0$, contradicting B.
The remaining B cases are $HH$, $H_r\beta_o$, and $\beta_rH_o$.

\subsection{Reusable endpoint and forest notation for the cards}
For every card below put $f=(179-\tau+R)/4$ unless a different trim is stated.
Its four-rank-trim supports certify the positive arms $\mathcal L_v-V_v-U_v$ and negative arms $\mathcal L_v-C_{7v}-C_{6v}$, with gaps $(X_v,g_v)$ and $(X_o,g_o)$. A star is the first edge of an arm. If a tree system has $m$ edges of total gap $W_F$ and grouped deficiency $D_F$, then $W_F\ge mf-D_F$. All forests in the cards use retained vertices, except the explicitly identified one-empty-vertex card.
Let $N$ be the sum of native deficiencies of the distinct endpoint resources. For a forest disjoint from these resources, $N+D_F\le\varepsilon=\delta-\eta$. Overlapping forest systems use separate copies of this same inequality; they are never declared a single disjoint forest.
Lemma~\ref{lem:z-rank-endpoint} supplies clips $k_Z$, errors $b_Z$, and $\Gamma_Z\ge k_Z$. For $k<74$, $b_Z\le\xi_{\mathrm{source}}$; for $k\ge74$, $b_Z\le\xi_{\mathrm{source}}+e_{\mathrm{source}}-4$. There can be at most one heavy selected endpoint whenever $N\le\varepsilon\le117$. Every outside-hull light endpoint also satisfies $b_Z\le k_Z$: since $29\le k\le73$, $k(75-k)\ge146>117$, while $\tau_v\le117-k$. The next-rank inequality gives $(74-k)b_Z\le\tau_v\le k(74-k)$, hence $b_Z\le k$; normal errors are zero. This also covers $k=73$, so weighted light errors are at most the corresponding weighted clips. Lower error sources are the target sign and upper error sources the opposite sign.

The source deletion arguments of Section~\ref{sec:Z3}, with the current mass coefficient 49 in place of $-29$, give
\begin{equation}\label{eq:z4-side}
 e_v\ge\max(4,A-49),\qquad e_v\le284-A.
\end{equation}
Explicitly the one-maximum envelope is $K\le2R+c_v+e_o+95$, giving $j_v-j_o\le141-A-U$ and $e_v\ge A-49$. The same-side two-maximum envelope is $K\le2R+b_v+e_o+169$, giving $2e_v\le576-2A-2a_v-2a_o-2b_o\le568-2A$.
The integer minima used below are
\begin{align*}
 Q_3&=T_r+4r+3I\ge641,& Q_4&=T_r+4r+4I\ge665,\\
 Q_5&=T_r+4r+5I=680+\tfrac12(r-31)(11(r-31)+1)\ge680,\\
 Q_{5/2}&=T_r+4r+\tfrac52I=623+3(r-28)^2\ge623,\\
 Q_{9/4}&=T_r+4r+\tfrac94I=612+\tfrac14(r-27)(11(r-27)-9)\ge612.
\end{align*}
The products prove the minima for all integers, without replacing the physical sizes by a sorted owner convention.


\subsection{Card F: a weighted five-endpoint transport}\label{sec:z4-cardF}
This card has three applications:
\begin{enumerate}
\item A-E with pair1 sole on $s$ and $Z\in\{U_s,U_t,V_s,V_t,H_t\}$;
\item A-H with beta sole on $t$, pair1 sole on $s$, and $Z\in\{U_s,U_t,V_s,V_t\}$;
\item A-H with beta sole on $s$ and $P$ empty.
\end{enumerate}
Use only the common fixed credit $\eta=8$, even when a stronger credit is available.
In the first two applications select lower endpoints $C_{4s},\alpha_s,\alpha_t,C_{5t}$ and upper endpoint $H_s$, with weights $1,2,1,1,1$. They are retained in every listed application. The identity is
\[
 H_s=C_{4s}+2\alpha_s+\alpha_t+C_{5t}+4(M+X)+2G.
\]
Use on each target $\mathcal L_v-V_v-U_v$ and $\mathcal L_v-C_{7v}$. There is at most one empty vertex, the displayed $U$ or $V$ hole; the separated-arm proof permits even an empty middle $V$ and assigns it zero capacity. The upper endpoint $H_s$ is not a forest vertex.
In the third application instead use lower endpoints $\beta_s,C_{5s},C_{5t},\alpha_t$ and upper endpoint $H_t$, with weights $1,1,1,2,1$. Its identity is
\[
 H_t=\beta_s+C_{5s}+C_{5t}+2\alpha_t+4(M+X)+2G.
\]
Both pairs in the identity are used only through an actually sole owner; the empty pair1 in the third application has no role. The same six-edge tree has every ordinary vertex retained. All endpoint labels are disjoint from its retained full resources.

Let $E,B$ be weighted clips and errors and $N$ the five unweighted native deficiencies. The endpoint constant is 236 in either orientation. Thus
\[
 S_{\rm endpoint}\ge236+4(M+X)+2G-E-B,
 \quad D_F\ge6f-(2X+G),\quad N+D_F\le\varepsilon,
 \quad E\le2N.
\]
Using the source envelope gives, with $r$ the opposite residual source side's cardinality and $\tau_{\mathrm{res}}$ its excess,
\begin{equation}\label{eq:z4Fmaster}
 T_r+4r+\tau_{\mathrm{res}}+2L-3\tau+2\eta\le228+\delta+B.
\end{equation}
For light endpoints $B\le E\le2\varepsilon$. A unique heavy endpoint has weight $q=1$ or 2, selected clip $c\in[74,77]$, and error at most $\tau_v+e_v$. Consequently
$B\le2\varepsilon-qc+q\tau+qe_v$ and $\delta\ge\eta+c$.
Use $e_v\le L$, whose remaining coefficient $2-q$ is nonnegative, then $L\ge A+41$ and $\tau+\delta\le125-I$. The complete scalar finish is
\begin{center}\small
\begin{tabular}{lll}
Case&Forced inequality&Contradiction\\\hline
No heavy&$Q_3+\tau_{\mathrm{res}}+4\eta\le271$&$673\le271$\\
Heavy weight $q$&$Q_{3+q}+\tau_{\mathrm{res}}+q\delta+4\eta+qc\le271+291q$&\(q=1:853\le562;\ q=2:1024\le853\)
\end{tabular}
\end{center}
This covers every outside-hull position and both error-source sides, and excludes all three applications.

\subsection{Card O: opposite native beta and pair1 owners}\label{sec:z4-cardO}
The exact predicate is: both $O$ and $H_s$ are empty; $\beta_s$ is a native whole resource (a sole beta pair in A or singleton in B); $C_{4t}$ is the sole pair1 resource; and all eight $U,V,C_6,C_7$ singletons are retained. There is no requirement on $\alpha$, $C_5$, or $H_t$ activity.
The eight $O$ HIGH labels, three $H_s+s_i$ labels and $\alpha_s+s_2$ give $\eta=12$, exactly as in Z3 A1; the last label is HIGH from actual beta activity, without requiring an alpha hit.
Use lower $\beta_s,C_{4t}$ and upper $V_t$ endpoints. Put
$B=\Gamma_{\beta_s}+\Gamma_{C_{4t}}$, $H=\Gamma_{V_t}$.
The three forest rows are
\[
 X_s+g_s\ge2f-(\varepsilon-B-H)/2,\quad
 X_s\ge f-(\varepsilon-B-H)/2,\quad
 X_s+X_t\ge2f-(\varepsilon-B)/2.
\]
Here the first forest is the physical-$s$ double arm, the second its two first edges, and the third the four stars. Only the last contains endpoint $V_t$, so its rebate is $B$ rather than $B+H$.
Use $f=\max(f_0,f_2,f_6)$ as in Z3 A1; then
$4\tau+10f\ge2020/3$. Adding three, two, and one copies of the displayed rows gives
\[
 6X_s+X_t+3g_s\ge10f-3\varepsilon+3B+\tfrac52H.
\]
The exact endpoint identity is
$V_t=2d_s+\beta_s+C_{4t}+2s_4+t_4+g_s$.
The leading maximum used twice retains the rebate
$E_{\rm clip}+2\ell_s\le2\varepsilon-(B+H)$.
Source substitution gives the master row
\[
 T_r+4r+5\xi_s+\tau_s+R+L+2e_s+10f+60+4B+\tfrac72H
 \le1018+4\delta+E_{\rm err},\qquad r=r_s.
\]
Light errors are at most $2\Gamma_Z$ by the rank argument already proved; heavy errors are $\xi_v+e_v-4$. Since $\varepsilon\le113$, at most one endpoint is heavy. Its error source is $s$ for $\beta_s$ or upper $V_t$, and $t$ for $C_{4t}$.
\begin{center}\small
\begin{tabular}{p{.30\linewidth}p{.43\linewidth}r}
Case&Forced upper row&Minimum\\\hline
Heavy, source $s$&$Q_4+4\xi_s+\tau_s+R+10f+4\tau\le1025$&$4222/3$\\
Heavy, source $t$&$Q_4+6\xi_s+\tau_s+10f+4\tau\le1077$&$4015/3$\\
All light, $A\le53$&$Q_4+5\xi_s+\tau_s+R+10f+4\tau\le1068$&$4222/3$\\
All light, $A\ge53$&$Q_4+5\xi_s+\tau_s+R+10f+4\tau\le1140$&$4222/3$
\end{tabular}
\end{center}
For the source-$s$ heavy case retain credit 259 and $L+e_s\ge L+4$. For the source-$t$ heavy case retain credit 296 and use $L+2e_s-e_t=3e_s\ge12$. In the all-light cases use respectively $e_s\ge4$ or $e_s\ge A-49$, along with $L\ge A+41$; the split point 53 is included on both sides. Then substitute $\delta\le125-I-\tau$. Each row contradicts $Q_4\ge665$, $R\ge69$, and $4\tau+10f\ge2020/3$.

\subsection{Card S: same native beta and pair1 owner}\label{sec:z4-cardS}
The predicate is both $O$ and $H_s$ empty, retained CORE8 $U,V,C_6,C_7$, both alphas retained, and distinct native resources $C_{4s},\beta_s$ on $s$. The beta may again be a sole pair or a singleton. The eleven HIGH $O,H_s$ labels give $\eta=11$.
Use lower endpoints $C_{4s},\beta_s,\alpha_t$ and upper endpoint $\alpha_s$, all with weight one. Let $N$ sum their native losses. The error sources have two entries on each physical sign, so their total error obeys
\[
 B\le2\xi_s+2\xi_t\quad\text{(all light)},\qquad
 B\le2\xi_s+2\xi_t+e_v-4\quad\text{(one heavy)}.
\]
At most one is heavy since $N\le\varepsilon\le114$.
Use the six-edge tree of Card F, and the three-edge physical-$t$ system $\mathcal L_t-V_t-U_t$ plus $\mathcal L_s-C_{7s}$. Their total gaps are $2X+G$ and $2X_t+g_t$, with deficiencies $D_6,D_3$. The endpoint resources avoid both systems and the leadings. Therefore
\[
 N+\ell_t\le\varepsilon,\quad N+D_6\le\varepsilon,\quad N+D_3\le\varepsilon,
 \qquad E+\ell_t+D_6+D_3\le3\varepsilon-2N.
\]
The identity $\alpha_s=C_{4s}+\beta_s+\alpha_t+d_t+2(M+X)+G$ and the leading bound $d_t\ge K+2X_t+g_t-\ell_t$ give, after source substitution,
\[
 T_r+4r+\tau_t+L-9\tau/4+(9/4)\xi_s+(25/4)\xi_t+3\eta+2N
 \le472+2\delta+B,\quad r=r_t.
\]
The radial coefficients absorb the baseline errors. For light endpoints use $L\ge A+41$; for a heavy error on source $v$, cancel $e_v$ from $L$ and retain $e_o\ge4$. Then $\tau+\delta\le125-I$ gives respectively
\[
 Q_{9/4}+\tau_t+\delta/4+3\eta+2N\le2349/4
 \quad\text{or}\quad
 Q_{9/4}+\tau_t+\delta/4+3\eta+2N\le2981/4.
\]
The light lower bound is $612+33=645>2349/4$. In the heavy case $N\ge74$ and $\delta\ge85$, so it is $612+85/4+33+148=3257/4>2981/4$. This is an all-position same-owner theorem.

\subsection{Card E-H: empty beta and the owner's $H$ hole}\label{sec:z4-cardEH}
The holes are $B,O_s,O_t,H_s$, pair1 is sole on $s$, and every other unit is retained. Use $\eta=8$. The five endpoints are lower $C_{4s},\alpha_s,\alpha_t,C_{5t}$ and upper $V_s$, all weight one. Let $N$ be their native-loss sum and $h=\Gamma_{V_s}$.
Use a six-edge system consisting of both double arms on $t$ and the negative double arm on $s$, with gap $X_s+2X_t+g_s+2g_t$. Use separately the four stars, of gap $2X$. The first avoids all endpoints; the second contains $V_s$ only. Hence
\[
 N+D_6\le\varepsilon,\quad N-h+D_*\le\varepsilon,
 \quad W:=3X_s+4X_t+g_s+2g_t\ge10f-D_6-D_*.
\]
The endpoint identity is
\[
 V_s=C_{4s}+\alpha_s+\alpha_t+C_{5t}+3s_4+4t_4+g_s+2g_t.
\]
Put $J_*=E+B-2N+h$. Every light lower endpoint contributes at most zero to $J_*$, and the upper endpoint contributes at most its error. Thus
\[
 J_*\le\xi_t+e_t-4\quad\text{if no lower endpoint is heavy},
\]
\[
 J_*\le\xi_t+\xi_v+e_v-4-c\quad\text{if a lower endpoint on source }v
 \text{ has heavy clip }c.
\]
The first case deliberately includes a heavy upper endpoint. Endpoint transport and the two forest budgets give the master
\[
 T_r+4r+\tau_t+L+2e_t+5\xi_t+\tfrac32R-\tfrac52\tau+2\eta
 \le1067/2+\delta+J_*,\quad r=r_t.
\]
For no heavy lower endpoint, cancel $e_t$ from its error and use $L\ge A+41$, $e_t\ge4$, and $\tau\le125-I-\delta$. The result is
\[
 Q_{5/2}+\tfrac32\delta+\tau_t+2\eta\le1137/2,
\]
whose left side is at least $639>1137/2$. If a lower endpoint is heavy, the remaining offsets $L+2e_t-e_v$ are at least 12 for either source $v$, and the remaining radial coefficients are nonnegative. The same cap yields
\[
 Q_{5/2}+\tfrac32\delta+\tau_t+2\eta+c\le1453/2.
\]
Now $c\ge74$, $\delta\ge82$, so the lower bound is $623+123+16+74=836>1453/2$.

\subsection{Card E-$\alpha$: either missing alpha}\label{sec:z4-cardEA}
Let $u$ be the sole actual pair1 owner, $o$ the opposite sign, and keep $\eta=8$. The holes are $B,O_s,O_t$ and one alpha. Raw outer order, rather than activity of that missing resource, supplies $\alpha_v\ge1-v_4$. Start from
$H_u=C_{4u}+2\alpha_u+\alpha_o+C_{5o}+4(u_4+o_4)+2G$.

If $\alpha_u$ is missing, use lower $C_{4u},\alpha_o,C_{5o}$ and upper $H_u$, all weight one. The identity becomes
\[
 H_u\ge C_{4u}+\alpha_o+C_{5o}+2u_4+4o_4+2G+2.
\]
Use the eight-edge system of all four double arms and the two physical-$o$ stars. Their gap sums are $2X+2G$ and $2X_o$, so together $W=2X_u+4X_o+2G\ge10f-D_8-D_2$. Both avoid all four endpoints, giving $E+D_8+D_2\le2\varepsilon-N$. With $r=r_o$ the exact master is
\[
 T_r+4r+\tau_o+3e_o-\tfrac52\tau+\tfrac12\xi_u+\tfrac{13}2\xi_o+2\eta+N
 \le502+\delta+B.
\]
The baseline errors are $\xi_u+3\xi_o$. After subtraction, use $\xi_u\le\tau_u$ to bound the remaining excess term below by $-3\tau$. For all light endpoints use $3e_o\ge3A-147$. For a heavy source-$o$ endpoint cancel $e_o-4$ and use $2e_o\ge2A-98$. For a heavy source-$u$ endpoint, which can only be $C_{4u}$, use additionally $e_u\le284-A$. The resulting exhaustive rows are
\begin{center}\small
\begin{tabular}{lll}
All light&$Q_3+2\delta+2\eta+N\le649$&$673>649$\\
Heavy on $o$&$Q_3+2\delta+2\eta+N\le721$&$895>721$\\
Heavy on $u$&$Q_4+3\delta+\tau+2\eta+N\le929$&$1001>929$
\end{tabular}
\end{center}
Here heavy means $N\ge74$, $\delta\ge82$.

If $\alpha_o$ is missing, use lower $C_{4u},\alpha_u,C_{5o}$ and upper $H_u$, with weights $1,2,1,1$. Raw order gives
\[
 H_u\ge C_{4u}+2\alpha_u+C_{5o}+4u_4+3o_4+2G+1.
\]
Use the six-edge system of Card F and the five-edge system with both positive double arms plus the negative star on target $o$. Their gaps total $4X_u+3X_o+2G$, and each avoids the endpoints. Since $E\le2N$, their two separate budgets give $E+D_6+D_5\le2\varepsilon$. The master is
\[
 T_r+4r+\tau_o+2L-\tfrac{11}4\tau+\tfrac{23}4\xi_o+\tfrac{11}4\xi_u+2\eta
 \le367+\delta+B.
\]
Baseline errors are $2\xi_o+3\xi_u$. The remaining excess expression is at least $\tau_o/4-3\tau$. In the light case use $2L\ge82+2A$; in a heavy case of weight $q=1$ or 2 use $qe_v\le qL$, leaving coefficient $2-q\ge0$. The cap gives
\[
 Q_3+2\delta+\tau_o/4+2\eta\le410\quad\text{(light)},
\]
\[
 Q_3+2\delta+\tau_o/4+2\eta+4q\le410+166q\quad\text{(heavy)}.
\]
The lower bounds are $673>410$, $825>576$ for $q=1$, and $829>742$ for $q=2$. Thus both alpha-hole positions, every peeled rank, and both physical source sides are covered.

\subsection{Card E-$C_5$: the forced opposite $C_5$ hole}\label{sec:z4-cardEC}
Here $P$ is sole on $s$ and the holes are $B,O_s,O_t,C_{5t}$. Use $\eta=8$ and the lower endpoints $C_{4s},\alpha_s,\alpha_t,C_{6t}$ and upper $H_s$, all weight one. Their identity is
\[
 H_s=C_{4s}+\alpha_s+\alpha_t+C_{6t}+3s_4+4t_4+2G.
\]
For the seven-edge system use both double arms on $s$, the positive double arm on $t$, and the negative star on $t$. Its gap is $2X_s+2X_t+g_s+2g_t$. For the four-edge system use the positive double arm and negative star on $s$, and the positive star on $t$. Its gap is $X_s+2X_t+g_s$. Both exclude endpoint $C_{6t}$ and all other endpoints. Consequently
\[
 3X_s+4X_t+2G\ge11f-D_7-D_4,
 \qquad E+D_7+D_4\le2\varepsilon-N.
\]
Transport gives, for $r=r_t$,
\[
 T_r+4r+\tau_t+L+2e_t-\tfrac{11}4\tau
 +\tfrac74\xi_s+\tfrac{27}4\xi_t+2\eta+N
 \le368+\delta+B.
\]
The baseline errors are $2\xi_s+3\xi_t$. The remaining excess expression is at least $-3\tau$, since its difference from $-3\tau$ is
$(\tau_s-\xi_s)/4+5\tau_t/4+15\xi_t/4\ge0$.
If light, use $L\ge A+41$ and $2e_t\ge8$. If heavy on source $t$, subtract $e_t-4$ and use $L+e_t\ge A+45$. Both give
\[
 Q_3+2\delta+2\eta+N\le569,
\]
with lower bound respectively $673$ and $895$. If heavy on source $s$, cancel $e_s$ using $L=e_s+e_t$, retain $3e_t\ge12$, and obtain
\[
 Q_3+2\delta+2\eta+N\le727<895.
\]
This covers the source-$s$ lower endpoints $C_{4s},\alpha_s$, and the source-$t$ endpoints $\alpha_t,C_{6t}$ and upper $H_s$.

\subsection{Dispatch of the complete retention table}
All applications use actual sole ownership and do not assign losses to an inactive paired component.
\begin{itemize}
\item In A-E, pair1 owner $s$, Card F excludes $U_s,U_t,V_s,V_t,H_t$. The other retained-pair1 choices are $H_s$, either alpha, or $C_{5t}$, excluded respectively by Cards E-H, E-$\alpha$, E-$C_5$. A same-side $C_{5s}$ hole is impossible because the cross chain would contradict pair1 ownership. The only remaining A-E choice is $Z=P$.
\item In A-H with beta owner $s$, Card F excludes the four $U,V$ choices. If $Z$ is $H_t,C_{5s},\alpha_s,$ or $\alpha_t$, pair1 is sole on $t$ and Card O applies. If $Z=C_{5t}$, pair1 is sole on $s$ and Card S applies. If $Z=P$, the third application of Card F applies. These exhaust A-H.
\item The B two-beta-hole case was excluded above. In B-$HH$ and B-$\beta_rH_o$, use Card O with $s=o,t=r$: beta on $o$ is a retained singleton and pair1 is sole on $r$. In B-$H_r\beta_o$, use Card S with $s=r,t=o$. All required CORE8 and alpha resources are retained in those exact hole sets.
\end{itemize}
Therefore the only remaining whole-$Z_4$ configuration would be branch A with exact holes $\{B,P,O_s,O_t\}$. Both pairs are genuinely empty and have no owner. The fourteen retained singleton resources are $H,U,V,C_5,C_6,C_7,\alpha$ on both signs.


\subsection{The ownerless final family: fixed credit and core gaps}\label{sec:z4-final}
Only the exact holes $\{P,B,O_s,O_t\}$ remain. All fourteen singletons
$H,U,V,C_5,C_6,C_7,\alpha$ are active. No empty pair is assigned an owner, endpoint, or forest role.
Write $p_v=v_4=M_v+X_v$. The active $C_5,C_7,\alpha$ give
\begin{equation}\label{eq:z4-final-centres}
\begin{split}
 d_v&\ge M+2X_v+X_o+G+h_o+1,\\
 d_v&\ge M+2X_v+X_o+g_v+1\ge M+132+I,\qquad
 h_v\ge X_v+1.
\end{split}
\end{equation}
Also $X+g_v\ge66$: twice this integral quantity strictly exceeds
$2X_v+X_o+g_v\ge131+I$.

\paragraph{A ten-label gap and eighteen extra misses.}
Below $H_v$ the largest ordinary centre is at most $\max(U_v,\alpha_v)$.
The leading maximum is below $U_v+M_v$. The raw-55 gap therefore gives
$\min(d_v,h_v)\le M+56$. As $d_v\ge M+132+I$, we have
$h_v\le M+56<d_v$ and $d_v-h_v-X_v\ge X+g_v-55\ge11$.
In the gap $(U_v+M_v,H_v-M_o)$ exactly the following ten fixed labels remain possible:
\[
 C^v_{33},C^v_{34},C^v_{44},C^v_{15},C^v_{26},C^v_{27},C^v_{28},
 U^v_2,U^v_3,V^v_2.
\]
These are fixed inventory labels; $U^v_i,V^v_i$ are distinguished by their indexed superscripts from the ordinary centres $U_v,V_v$.
The omitted labels have the following complete explanation:
$C_{11},\ldots,C_{14}$ are HIGH;
$C_{22},C_{23},C_{24}$ and $C_{16},C_{17},C_{18}$ exceed the upper endpoint;
$C_{25}$ is below the lower endpoint since $o_1>d_o>h_v$;
$C_{3i},C_{4i}$ ($5\le i\le8$) and the ten negative-negative $C$ labels lie below;
so do $U_i,V_i$ ($5\le i\le8$), $U_4=h_v$, $V_3=g_v+h_v-d_v$, $V_4=h_v-d_v$, and all five fixed shapes $F_i\le g_v$.
Here $d_v-o_2+M_o>0$ follows from $C_{5v}$ activity, and $g_v<U_v+M_v$ follows from \eqref{eq:z4-final-centres}.
The ten-label count improves the gap cap to
\begin{equation}\label{eq:z4-h11}h_v\le M+11<d_v.\end{equation}
Combining it with \eqref{eq:z4-final-centres} gives
$d_v-g_v-h_v\ge123$ and $\alpha_v-o_3\le11-X-g_o\le-55$.
Both beta centres are LOW because $\beta_v=h_v-d_v-p_v<-M_v$.
Both pair1 centres are LOW also: $C_{4v}=V_v-v_4-o_1<V_v-M_o\le S_v$, excluding the HIGH alternative for this empty centre.

Per sign the following sixteen distinct fixed labels are nonpositive:
\[
 C_{55},C_{56},C_{57},C_{58};\quad
 V_3,V_4,V_5,V_6,V_7,V_8;\quad
 U_5,U_6,U_7;\quad F_3,F_4,F_5.
\]
Their 32 labels occupy 25 of the 26 baseline groups: among the four unreflected rows they meet $C_{58}$, $U_6$, $V_6$, but not $C_{25}$; they meet both reflections of each of the eleven other rows. The two $U_7$ labels belong to no baseline group.
There are three further LOW labels in already occupied groups: at least one of each reflected pair $(V^s_2,V^t_2)$, $(C^s_{35},C^t_{35})$, $(C^s_{45},C^t_{45})$. Indeed their sums obey
\[
 V^s_2+V^t_2\le20-3X-G<0,\quad
 C^s_{35}+C^t_{35}=-2(p_s+p_t)-G-h_s-h_t<0,
\]
\[
 C^s_{45}+C^t_{45}=-2(p_s+p_t)-2G-h_s-h_t<0.
\]
For the first strict sign, add the two four-edge rows to obtain $3X+G\ge262+2I$; the relevant groups are already occupied respectively by $C_{56}$, $V_3$, $V_8$.
Thus 35 LOW labels occupy only 25 groups, and the untouched $C_{25}$ group supplies another LOW witness. Together with the eight HIGH $O$ labels, there are 44 misses instead of the baseline 26. The genuine extra fixed credit is therefore
\begin{equation}\label{eq:z4eta18}\eta=18,\qquad\varepsilon=\delta-\eta\ge0,\qquad\varepsilon\le107.\end{equation}
The eight HIGH labels are included in 18 and are not added a second time.

\paragraph{Core-hit versus truncated $H$ gap.}
For fixed physical $v$ set
\[
 a=U_v+R_v=d_v-X_v-e_v,\qquad b=H_v-R_o,
 \quad T=\min(b,S_v+1),\quad\mathcal J=(a,T)\cap\mathbb Z.
\]
The $C_5$ bound gives $U_v>M_o$, so $a,b>0$.
The $H_v$ resource has a core hit exactly when $b\le S_v$, because its lowest core image is $b$ and belongs to the core translate.
If it has a core hit, the signed gap count is $h_v-R-1$; if not, the full $H_v$ activity yields signed count at least $h_v-(R_v+M_o)$.
Every ordinary centre except $H_v$ is at most $U_v$ by \eqref{eq:z4-h11}; its core contributions are below $a$. The full $C_5,C_6,C_7$ resources lie below $a$ as well. Initially at most 12 ordinary positive-peel labels from $U_v,V_v,\alpha_v$, four negative-peel labels from $H_v$, and ten positive-peel-pair leading labels can enter $\mathcal J$.
Among fixed labels the only possible eleven are
\[
 C_{33},C_{34},C_{44},C_{15},C_{26},C_{27},C_{28},U_2,U_3,U_4,V_2.
\]
The elimination is the preceding complete gap-label list with its lower endpoint now $U_v+R_v$: $C_{25}-a=h_v-o_1-R_v<0$; $C_{16},C_{17},C_{18}>b$ from $d_v-o_2+R_o>0$; $C_{3i},C_{4i}$ and negative-negative pairs are below $a$; $U_{5:8},V_{5:8},V_3,V_4,F_{1:5}$ are below $a$; and the seven listed HIGH or upper $C$ labels are outside the gap. This excludes exactly 44 fixed labels.
The first total $12+4+10+11=37$ proves
\[
 h_v\le R+38\ \text{with a core hit},\qquad
 h_v\le R_v+M_o+37\ \text{without one}.
\]
In particular the latter bound holds in both cases since $e_o\ge4$.
It removes all four alpha-peel contributions, as
$a-(\alpha_v+M_v)=d_v-h_v-e_v\ge95+I>0$.
It also removes $U_3,U_4,V_2$, since
\[
 a-U_3\ge X+g_o+h_o-36\ge31>0,
 \qquad U_4<U_3,\quad V_2<U_3.
\]
The remaining count is $8+4+10+8=30$. Consequently
\begin{equation}\label{eq:z4-coreH}
 h_v\le R+31\quad\text{if }H_v\text{ has a core hit},\qquad
 h_v\le R_v+M_o+30\quad\text{otherwise}.
\end{equation}
In the latter case $H_v-R_o>S_v$; the 74 source points at or above $-R_o$ all give upper misses. Thus its actual native loss is at least 74. This localization uses membership in the complete target, not an untruncated hull-gap assumption.

\subsection{Four final cases, with every alpha endpoint}
For each retained $\alpha_v$, let $\Gamma_v$ be its native deficiency, $k_v$ its lower clip, and define the actual error
$\zeta_v=\max(0,74-r_v-k_v-\alpha_v)$.
Then
\[
 N:=\Gamma_s+\Gamma_t\ge k:=k_s+k_t,\quad
 \alpha_s+\alpha_t\ge79-k-B,\quad B=\zeta_s+\zeta_t.
\]
For $k_v<74$, the next-rank argument gives $\zeta_v\le\xi_v$; a normal positive alpha gives zero error. Since $N\le\varepsilon\le107$, at most one lower alpha clip is heavy ($\ge74$). No bound on its possibly unbounded peeled-rank error will be needed.
Use $f=62-(\tau-\xi)/4$ and designate each full leading as root. The two physical-$v$ stars are $\mathcal L_v-V_v$ and $\mathcal L_o-C_{7o}$; they are disjoint from both alphas and both $H$ resources. Across both signs their four edges form two three-vertex trees. Write $\ell=\ell_s+\ell_t$. The rooted forest lemma gives the two separately valid rows
\begin{equation}\label{eq:z4stars}
 2X_v\ge2f+N+\ell-\varepsilon,\qquad
 2X\ge4f+N+\ell-\varepsilon.
\end{equation}
If an $H_v$ has no core hit, its actual upper clip $u\ge74$ is distinct from all star and alpha resources, giving instead
\begin{equation}\label{eq:z4starsH}
 N+u\le\varepsilon,\qquad2X\ge4f+N+u+\ell-\varepsilon.
\end{equation}

\paragraph{Both $H$ core-hit; both alpha clips light.}
Now $B\le\xi$ and \eqref{eq:z4-coreH} bounds $h_s+h_t\le2R+62$.
Using $\alpha_s+\alpha_t=h_s+h_t-M-X$ gives $L+X\le R-17+k+B$.
Insert $L\ge A+41$ and the four-star row to get exactly
\[
 226+2I+\tau+\delta+\eta+N+\ell\le\xi+2k+2B.
\]
Use $N\ge k$, $k\le\varepsilon=\delta-\eta$, $B\le\xi\le\tau$ and $\tau+\delta\le125-I$. This forces
$4I+4\eta+\ell\le24$, contradicting $\eta=18$.

\paragraph{Both $H$ core-hit; one alpha clip heavy.}
Say $k_v\ge74$; the opposite clip is light. The opposite core-H cap and its alpha endpoint give
$e_o+X_o\le26+\xi_v+k_o+\zeta_o$.
Use $e_o\ge4$, the two physical-$o$ stars, $N\ge k_v+k_o$, and $k_o\le\varepsilon-k_v$:
\[
 2f+2k_v+\ell\le44+2\xi_v+2\zeta_o+2\varepsilon\le44+2\xi+2\varepsilon.
\]
Since $2f=124-\tau/2+\xi/2$, this implies
$228+2\eta+\ell\le\tau/2+3\xi/2+2\delta\le250-2I$,
or $2I+2\eta+\ell\le22$, again impossible. The heavy alpha error was never bounded or discarded.

\paragraph{Exactly one $H$ has no core hit.}
Its upper count $u\ge74$ gives $N\le33$, so both alpha clips are light and $B\le\xi$. The two H caps give
$h_v+h_o\le2R+e_o+61$, hence $e_v+X\le R-18+k+B$.
Insert $e_v\ge4$ and \eqref{eq:z4starsH}; using $N\ge k$ and $k\le\varepsilon-u$ yields
\[
 4f+2u+\ell+44\le2R+2B+2\varepsilon.
\]
Therefore $302+2\eta+\ell\le\tau+3\xi+2\delta\le4\tau+2\delta$.
The cap gives $4\tau+2\delta\le500-4I-2\delta$, while the same resource budget gives $\delta\ge\eta+74$. Thus
$4I+4\eta+\ell\le50$, contradicting $\eta=18$.

\paragraph{Neither $H$ has a core hit.}
Their two distinct native deficiencies total at least $148>125\ge\delta$, impossible.
These four cases exhaust all H core-hit statuses and all lower-alpha ranks. They include both-low alpha, source-gap, fringe, coincident-centre and outside-hull placements. The ownerless residual is excluded, and therefore every original $Z_4$ configuration is impossible.


% Self-contained mathematical Z5 section; common foundations precede this file.
\section{Five missing units: $Z_5$}\label{sec:z5}
We now assume $z=0$, $j=5$.  Throughout this section $s,t$ are
physical signs; a sign is renamed only to identify an actual owner or an
actual missing singleton.  In particular it is never renamed to sort the
core sizes. Write
\[
 \begin{gathered}
 \mathcal B=\beta_s+\beta_t,\quad K=|D+D|,\quad \mathcal A=I+\tau+\delta,\\
 a=s_4=M_s+X_s,\quad b=t_4=M_t+X_t,\quad X=X_s+X_t,\quad G=g_s+g_t.
 \end{gathered}
\]
Here $a,b$ without subscripts are outer magnitudes, whereas
$a_v<b_v<c_v<e_v$ denote the four source offsets above $R_v$.
We abbreviate $P_v=C4_v$ and, in the proof cards only,
$U_v=\mathrm{own2}_v$, $V_v=\mathrm{own3}_v$.
The symbols $U_i^v,V_i^v,C_{ij}^v,F_i^v$ with a numerical subscript
always denote \emph{fixed labels} of the common covering lemma; $F_i^v$
are its five fixed exceptional labels. They are not ordinary resources.

\subsection{Source bounds and the actual-retention classification}
Since $w\le4z=0$, both targets are $F_v=[1,S_v]\cap\mathbb Z$.
The exact original slack identity and the structural budget give
\begin{equation}\label{eq:z5-source}
 S=2K+1098-\delta,\qquad
 \delta+\rho_s+\rho_t+\kappa\le51-I.
\end{equation}
Consequently $28\le r_s,r_t\le41$. The source-only envelope lemma,
with $m_0=127$, gives
\begin{equation}\label{eq:z5-envelope}
 \mathcal A\le47,\quad M\ge R+119+\mathcal A,\quad
 K-M\ge R-35+\mathcal A,\quad K\le M+R+21.
\end{equation}
In particular $R\le69+\tau$ and $R+\delta\le116-I$.
We shall sometimes use the deliberately weaker inequalities
\begin{equation}\label{eq:z5-weak}
 R+\delta\le124-I,\qquad 2\tau+3\delta\le225-3I.
\end{equation}
All occurrences of $\delta$ in these source inequalities retain their
original meaning.

Here are two useful sharpenings needed in the cards:
\begin{equation}\label{eq:z5-weighted}
 M\ge R+121+\mathcal A,\qquad
 M+M_v\ge R+R_v+158+2I+\tau+2\delta.
\end{equation}
We include their source-only proof. Put $L=e_s+e_t$, and let $B_v$
count distinct sums of two offsets exceeding $e_v$, repetitions allowed;
put $Q_v=2B_v-a_v-b_v-c_v$. The four-offset lemma gives
$2B_v\le a_v+b_v+8$ and $Q_v\le3$. Its equality case is
\[
 Q_v=3\quad\Longleftrightarrow\quad
 (a_v,b_v,c_v,e_v)=(1,2,4,5)\text{ or }(1,3,5,6).
\]
For completeness, if $c\le a+b$, the bound $B\le c+1$ gives $Q\le2$.
Otherwise $2a,a+b\le e$, and $a+b+c\ge14$ gives $Q\le2$.
The remaining pairs $(a,b)$ and possible $c$ are
$(1,2):4\ldots10$, $(1,3):5\ldots9$, $(1,4),(2,3):6\ldots8$,
and $(1,5),(2,4):7$; the respective bounds on $B$ are $5,6,6,6,7,7$.
Directly comparing the at most six non-$e$ pair sums with the four
$e$-containing sums gives the two equality cases above.

The central interval plus fringes and the exact mass give
$L\ge125+\mathcal A-Q_s-Q_t$. Both $Q_v=3$ would give $L\le12$,
contrary to $L\ge119+\mathcal A$. Thus $L\ge120+\mathcal A$.
If equality holds, one side has $e_b\in\{5,6\}$ and the other has
$Q_a=2$, $e_a\ge114+\mathcal A$. Since
$a_a+b_a+c_a=2B_a-2\le18$, we have $c_a\le15$, hence
$e_a>2c_a$, $B_a=4$, and $(a_a,b_a,c_a)=(1,2,3)$.
The central interval then misses
$[2R_a+7,M_a-M_b-1]$, whose length is at least
$101+\mathcal A-R\ge32+I+\delta>0$. Equality in the
central-interval bound required no such hole. This proves the first row
of \eqref{eq:z5-weighted}.

For the second, distinguish sign $s$. Apart from the one cross-maximal
sum $M_t-M_s$, no sum lies strictly between $2R_t+2c_t$ and
$M_t-R_s-c_s$. Thus the central interval has at least
$H=\max(0,e_t-R-2c_t-c_s-2)$ holes and
\[
 L\ge125+\mathcal A-Q_s-Q_t+2H.
\]
If $e_t\ge2c_t$, then $B_t=4$, $Q_t\le5-c_t$.
The inequality $2c+4\max(0,C-2c)\ge C$ gives
$L+e_s\ge238+2\mathcal A-R-(2Q_s+c_s)$.
Since $Q_s\le\min(3,8-c_s)$, $2Q_s+c_s\le11$.
Therefore $L+e_s\ge227+2\mathcal A-R
\ge158+2I+\tau+2\delta$.
If $e_t<2c_t$, then $Q_t\le17-c_t<17-e_t/2$, so
$L+e_s>210+2\mathcal A$, a stronger bound. Reflection proves both signs.

Both leadings are active, since $K\ge155>\delta$, so their centers
$\lambda_v$ are integral. With native leading loss $\ell_v$,
\begin{equation}\label{eq:z5-leading}
 d_v\ge K+g_v+2X_v-\ell_v,\qquad
 d_v-M\ge34+g_v+2X_v.
\end{equation}
A retained singleton has at least $78-\delta\ge31$ hits, hence a core
hit. The same holds for a pair only after sole ownership is proved.
No such assertion is made for one component of a split pair.

The raw55 gaps now force
\begin{equation}\label{eq:z5-baseholes}
 O_s,O_t\text{ empty},\qquad \beta_v,H_v\text{ not both active}.
\end{equation}
Here are the details that precede outer integrality. Activity of $O_v$,
or of $\beta_v,H_v$ with $O_v$ empty, gives $d_v\le M+56$.
Together with \eqref{eq:z5-envelope} this gives
$g_v+2X_v\le22-I-\tau$. If $U_v$ or $C6_o$ were active,
$q=v_3-M_v=X_v+g_v$ would be integral and a leading/core intersection
would give $q\ge139-R-\delta\ge23+I$, impossible. The same argument
with $q=X_v$ excludes $V_v,C7_o$. Thus the cap forces four distinct
singleton holes $U_v,V_v,C6_o,C7_o$.
The opposite $O_o$ supplies a fifth hole. If $\beta_v,H_v$ were both
active, the double raw55 gap excludes $O_v$, giving six holes.
If $O_v$ were active, then in A the opposite $\beta_o$ and $H_o$
would both be active; in B the singleton $\beta_v$ would be a sixth
hole. This proves \eqref{eq:z5-baseholes}. The double gap uses the
same 55 labels for two disjoint intervals and gives
$d_v-M\le57/2$, contrary to \eqref{eq:z5-leading}.

Once a coordinate is integral, an empty own center above a retained own
center is HIGH, at least $S_v+R_o+1$; a larger retained singleton would
then lose at least 74 source points. Thus empty $V_v$ forces $U_v,H_v$
empty, and empty $U_v$ forces $H_v$ empty. An empty cross center is LOW:
otherwise the leading loses at least $r_v+82>\delta$ points. Hence
\begin{equation}\label{eq:z5-crosschain}
 P_v\text{ active}\Longrightarrow C5_v\text{ active}
 \Longrightarrow C6_v\text{ active}\Longrightarrow C7_v\text{ active}.
\end{equation}
For clarity, the actual-owner identities are
\[
\begin{split}
 P_s+P_t&=-(s_2+s_3+s_4+t_2+t_3+t_4)<-3M,\\
 \beta_s+\beta_t+C5_s+C5_t&=-2(s_3+s_4+t_3+t_4)<-4M,\\
 P_s+\beta_s+C5_t&=-2s_3-2s_4-t_2-t_3-t_4<-4M_s-3M_t.
\end{split}
\]
Any active center on side $v$ is greater than $-M_v$. The first
identity proves sole pair1 ownership; the second excludes simultaneous
activity of all four centers, so split beta forces a $C5$ hole; the
third excludes the three displayed active centers. These are activity
statements about the original whole units, not selected witnesses.

In A0 the beta pair is empty. Its three base holes are
$\{\mathrm{pair2},O_s,O_t\}$, leaving two more. A nonintegral $v_3$
would force $U_v,C6_o$ empty; a nonintegral $v_4$ would force
$V_v,C7_o$ empty. Those two holes leave $H_v,\alpha_v$ and the other
lower own center retained, contradicting
$v_3=H_v-V_v-\alpha_v$ or $v_4=H_v-U_v-\alpha_v$.
Then a nonintegral $v_2$ would force the three holes
$H_v,\alpha_v,C5_o$. All eight outer coordinates are therefore integral.
The own and cross chains prove that all four extreme leaves
$V_s,V_t,C7_s,C7_t$ are retained. The extra holes are precisely
$\{U_v,H_v\}$, $\{C6_v,C5_v\}$, or two labels from
$\{H_s,\alpha_s,H_t,\alpha_t,\mathrm{pair1},C5_s,C5_t\}$,
except two $C5$ holes. If pair1 is retained, normalize $s$ to its actual
sole owner; $C5_s$ is retained and the thirteen rows below exhaust the
possibilities. If it is empty, full physical reflection leaves exactly
one additional hole $H_s,\alpha_s$, or $C5_s$.

In A1, a retained beta pair also has a sole actual owner. Indeed suppose
both components are active. The holes are $O_s,O_t,H_s,H_t$ and exactly
one $C5_t$; pair1 has sole owner $s$. All lower own centers and both
alphas are retained, proving all outer coordinates integral. Let $N$ be
the total number of beta hits; the whole-pair budget gives
$78-\delta\le N\le78$ and $\beta_s+\beta_t\ge N-M$.
Using the endpoints of $P_s,\alpha_s$, the leading inequalities and
$\alpha_s-P_s=\beta_s+\beta_t+2d_t+2(a+b)+G$ gives
\[
 S_t+M+2X_s+6X_t+g_s+3g_t\le1019+R+2\delta.
\]
The four retained extreme leaves give
$2X_s+6X_t\ge1112-8R-3\delta$, while the four ordinary pairs give
$g_s+3g_t\ge556-4R-3\delta/2$.
Also $S_t+M\ge T_r+4r+R+129+\mathcal A+\tau_t$, where $r=r_t$.
Using $R\le69+\tau$ and $\mathcal A\le47$ yields
$T_r+4r+12I\le567$. But
\begin{equation}\label{eq:z5-q717}
 T_r+4r+12(r-34)(r-35)
 =717+\tfrac12\{25(r-33)^2+3(r-33)\}\ge717
\end{equation}
for integer $r$. In this argument $N$ was charged only to the whole
pair, never to a selected component.

Normalize A1 so $\beta_s$ is its actual sole owner. The base holes are
$O_s,H_s,O_t$. A nonintegral lower coordinate again costs two extra
holes. On side $t$, the retained $H_t,\alpha_t$ and other lower own
center give a contradiction. On side $s$, the retained $\alpha_s$ and
$\beta_s$ give integral $d_s$, $s_3+s_4$, and $s_2$, also a
contradiction. Nonintegral $t_2$ costs $H_t,\alpha_t,C5_s$.
Thus all coordinates are integral. Empty $V_t$ would cost $U_t,H_t$;
empty $V_s$ would cost $U_s$ and $\alpha_s$ (since actual $\beta_s$
puts $\alpha_s$ above $U_s$); empty $C7_v$ costs $C6_v,C5_v$.
All four extreme leaves are again retained. With pair1 empty the final
extra hole is exactly
\begin{equation}\label{eq:z5-a1list}
 \eta\in\{\alpha_s,H_t,\alpha_t,C5_s,C5_t,U_s\}.
\end{equation}
The retained-pair1 alternatives will be treated below, including the
same-owner pattern with empty $\alpha_s,C5_t$.

\subsection{A reusable endpoint and loss calculus}
For a retained singleton or actual sole pair component $Z$ on target $v$,
let $l_Z,u_Z$ be its actual lower and upper clipping counts. Then
\begin{equation}\label{eq:z5-end}
 Z\ge74-r_v-\tau_v/23-l_Z,\qquad
 Z\le S_v+r_o+\tau_o/23+u_Z-73.
\end{equation}
To prove the first bound, delete the four positive peeled points and 22
largest positive core points. The remaining 52 source points have maximum
at most $r_v-22+\tau_v/23$; at least $52-l_Z>0$ images are positive
integers. Reflect the argument with the 52 highest points for the second
bound. This covers BAD centers as well as normal ones. If
$1\le Z\le S_v$, the sharper bounds omit both $\tau/23$ errors,
by counting the $r_o+4$ negative and $r_v+4$ positive source points.
A BAD center has at most $\max(r_s,r_t)+4$ hits and therefore deficiency
at least $74-\max(r_s,r_t)$.

For a group of $m$ retained translates let $\Gamma$ be $78m$ minus
its actual target union size. For a common $q$-point source subset $Q$
and one chosen lower or upper clipping count $u$,
\begin{equation}\label{eq:z5-joint}
 \Gamma\ge u+mq-\left|\bigcup_i(Z_i+Q)\right|.
\end{equation}
Pointwise, inside the target multiplicity minus one pays the overlap
excess; outside it multiplicity also pays the one clipping indicator.
A subgroup has deficiency at most that of the whole group. Thus a pair
inside a triple never receives a second independent allocation.
If just one member of a two-resource pair is retained, its common-$Q$
intersection with the empty translate lies outside the target and is
paid by its actual off-target loss; the empty member receives no capacity.

Put $\tau'=\tau-R+69\ge0$. Delete the four largest core magnitudes on
each sign to get $Q_{62}$ with span at most $61+\tau'/4$.
Indeed the monotone core deviations $x_i-i$, excluding the maximum,
have sum $\tau'_v$, so $x_{r_v-4}\le r_v-4+\tau'_v/4$.
Set $f=62-\tau'/4$. For a pair of gap $g$ and a triple of span $g+h$,
\begin{equation}\label{eq:z5-pairs}
 g\ge f-\Gamma,\qquad g+h\ge f+62-\Psi.
\end{equation}
The four disjoint ordinary groups
\[
 \mathcal G_s=\{\{U_s,V_s\},\{C6_t,C7_t\}\},\qquad
 \mathcal G_t=\{\{U_t,V_t\},\{C6_s,C7_s\}\}
\]
have summed deficiencies $\Omega_s,\Omega_t$ and give
$g_v\ge f-\Omega_v/2$, $\Omega_s+\Omega_t\le\epsilon$.
Here $\epsilon=\delta-e$ if compatible fixed credit $e$ is proved,
and otherwise $e=0$. The separately valid stars are
\begin{equation}\label{eq:z5-stars}
 2X_v+2R+\epsilon\ge278,\qquad
 2X+4R+\epsilon\ge556.
\end{equation}
Each comes from the leading hubs and the extreme leaves; opposite hubs
inside one leading are disjoint because $M>R$. Combining independently
valid scalar rows does not declare their resource groupings disjoint.
Only a single displayed regrouping may share one $\epsilon$ budget.

All endpoint transports in the A0 table are obtained from
\begin{equation}\label{eq:z5-identities}
\begin{aligned}
 P_s&=d_s-d_t-2a-b-G-h_t,& C5_t&=d_t-2b-a-G-h_s,\\
 C6_t&=d_t-2b-a-G,& C7_t&=d_t-2b-a-g_t,\\
 U_s&=d_s-a,&U_t&=d_t-b,\\
 V_s&=d_s-a-g_s,&V_t&=d_t-b-g_t,\\
 H_s&=d_s+h_s-a,&H_t&=d_t+h_t-b,\\
 \alpha_s&=h_s-a,&\alpha_t&=h_t-b.
\end{aligned}
\end{equation}
This also provides a direct algebra check for every coefficient below.

\subsection{A0 with pair1 retained: the thirteen proof cards}
In the following exhaustive table $s$ is the actual pair1 owner and the
base holes are $\{\mathrm{pair2},O_s,O_t\}$. Superscripts on endpoint
lists are multiplicities. An upper endpoint has its coefficient one
unless indicated. Every listed resource is retained in its row.
\begin{center}\small
\begin{tabular}{c|l|l|l|c}
 & Extra holes & Upper endpoints & Lower endpoints & Card\\\hline
1&$\alpha_s,\alpha_t$&$H_s,H_t$&$P_s,(C5_t)^2$&D\\
2&$U_s,H_s$&$V_s,V_t$&$P_s,(C5_t)^2,\alpha_s^2,\alpha_t$&A\\
3&$U_t,H_t$&$V_s,V_t$&$P_s,(C5_t)^2,\alpha_s^2,\alpha_t$&A\\
4&$H_s,H_t$&$V_s,V_t$&$P_s,(C5_t)^2,\alpha_s^2,\alpha_t$&A\\
5&$\alpha_t,H_s$&$V_s,H_t$&$P_s,(C5_t)^2,\alpha_s^2$&B\\
6&$\alpha_t,H_t$&$H_s,U_t$&$P_s,(C5_t)^2,\alpha_s^3$&C\\
7&$\alpha_s,H_t$&$H_s,U_t$&$P_s,(C5_t)^2,\alpha_t$&E\\
8&$\alpha_s,H_s$&$U_s,H_t$&$P_s,(C5_t)^2,\alpha_t^2$&F\\
9&$C5_t,\alpha_s$&$H_s,H_t$&$P_s,(C6_t)^2,\alpha_t^2$&G\\
10&$C5_t,\alpha_t$&$H_s,H_t$&$P_s,(C6_t)^2,\alpha_s$&H\\
11&$C5_t,H_t$&$H_s,U_t$&$P_s,(C6_t)^2,\alpha_s,\alpha_t$&I\\
12&$C5_t,C6_t$&$H_s,H_t$&$P_s,(C7_t)^2,\alpha_s,\alpha_t^2$&J\\
13&$C5_t,H_s$&$U_s,H_t$&$P_s,(C6_t)^2,\alpha_t^2$&K
\end{tabular}
\end{center}
Here and below the pair resource $P_s$ is charged once, irrespective of
its endpoint coefficient; its unused opposite component has no capacity.
Set $r=r_s$, $k=r_t=69-r$. For example, upper $H_s,H_t$, lower
$P_s$, and twice lower $C5_t$ give exactly
\[
 S\ge3a+4b+3G+3h_s+2h_t+161+r
 -(2\tau_s+3\tau_t)/23-E,
\quad E=u_{H_s}+u_{H_t}+l_{P_s}+2l_{C5_t}.
\]
Replacing an upper $H_v$ by $U_v$ or $V_v$ subtracts $h_v$ or
$h_v+g_v$, respectively. Replacing $C5_t$ by $C6_t$ removes its
$h_s$ term; replacing it by $C7_t$ also removes its $g_s$ term.
Each lower alpha adds the corresponding lower bound to $h_v-a_v^{\rm out}$.
Thus the following cards are derivations from explicit endpoints, rather
than a dispatch to unnamed cases.

\paragraph{Cards A, B, C: retained owner alpha and opposite cross.}
For A, the endpoint transport, with maximum loss coefficient two, is
\[
 3M+5X+2G\le2R+826+4\tau/23+\delta.
\]
Use $3M\ge3R+363+3\mathcal A$,
$5X\ge1390-10R-5\delta/2$, and
$2G\ge248-\tau'-\delta$. Then
\[
 1106+3I+42\tau/23\le8R+3\delta/2,
 \qquad 228+22I+84\tau/23+13\delta\le0.
\]
For B the transport is
\[
 2M+M_s+5X_s+4X_t+2g_s+3g_t+2h_t
 \le2R+831+r+(4\tau_s+3\tau_t)/23+\delta.
\]
Use the two groups in $\mathcal G_s$, the triple $\{H_t,U_t,V_t\}$,
and $\{C6_s,C7_s\}$. Their distinct-label deficiencies give
$2g_s+3g_t+2h_t\ge434-5\tau'/4-2\delta$.
Also $2M+M_s\ge2R+246+r+2I+2\tau+2\delta$ and
$5X_s+4X_t\ge1251-9R-5\delta/2$. Hence
\[
 4055/4+2I+(53\tau_s+57\tau_t)/92\le31R/4+7\delta/2,
 \quad211+39I+(53\tau_s+57\tau_t)/23+17\delta\le0.
\]
For C, the transport is
\[
 2M+2M_s+6X_s+4X_t+3G+h_t
 \le2R+757+2r+(5\tau_s+3\tau_t)/23+2\delta.
\]
Drop $h_t>0$ and use $2M+2M_s\ge2R+250+2r+2I+2\tau+2\delta$,
$6X_s+4X_t\ge1390-10R-3\delta$, and
$3G\ge372-3\tau'/2-3\delta/2$. This gives
\[
 2303/2+2I+(13\tau_s+17\tau_t)/46\le17R/2+9\delta/2,
 \quad195+21I+(13\tau_s+17\tau_t)/23+8\delta\le0.
\]
All three final inequalities are impossible.

\paragraph{Cards D, E, F: the BAD cases and their normal refinements.}
These three cards must first establish normality. Their respective
transports are
\begin{align*}
D:\quad&M_s+2M_t+3X_s+4X_t+3G+3h_s+2h_t
 \le2R+979-r+(2\tau_s+3\tau_t)/23+E_D-\delta,\\
E:\quad&M+2M_t+3X_s+5X_t+3G+3h_s
 \le2R+836+2k+(2\tau_s+4\tau_t)/23+E_E-\delta,\\
F:\quad&M+3M_t+3X_s+6X_t+3G+2h_s
 \le2R+762+3k+(2\tau_s+5\tau_t)/23+E_F-\delta.
\end{align*}
The loss expressions, in the order of the endpoint table, are
\[
\begin{split}
E_D&=u_{H_s}+u_{H_t}+l_P+2l_{C5_t},\\
E_E&=u_{H_s}+u_{U_t}+l_P+2l_{C5_t}+l_{\alpha_t},\\
E_F&=u_{U_s}+u_{H_t}+l_P+2l_{C5_t}+2l_{\alpha_t}.
\end{split}
\]
For D group $\{H_s,U_s,V_s\}$, $\{H_t,U_t,V_t\}$ and
$\{C6_s,C7_s\}$, with deficiencies $\Psi_s,\Psi_t,\Omega$;
leave $P,C5_t$ as separate resources. With their chosen upper losses,
\eqref{eq:z5-joint} on the raw core gives
\[
 3G+3h_s+2h_t\ge1184-6R-3\Psi_s-2\Psi_t-\Omega
                  +3u_{H_s}+2u_{H_t}.
\]
Use $M_s+2M_t\ge R+125+k+\mathcal A$ and
$3X_s+4X_t\ge973-7R-2\delta$. Each remaining resource-deficiency
coefficient is at most three. Thus
\begin{equation}\label{eq:z5-Dboot}
 15I+(21\tau_s+20\tau_t)/23+11\delta\le364.
\end{equation}
It forces $\delta<34$, $I<25$, hence sizes $30\ldots39$.
Every retained resource then has more than 44 hits, whereas a BAD center
has at most 43. All used endpoints are normal. Both empty alphas are LOW
because they lie below normal $H_v$; translating $M_v$ gives $h_v\le X_v$.
Apply the joint inequality to each triple and its contained
$\{H_v,U_v\}$ pair with the \emph{same} $\Psi_v$. It gives
$g_v+2h_v\ge186-\tau'/2-2\Psi_v+2u_{H_v}$.
The normal transport therefore implies
$124+2R+I\le2\tau+4\delta$, hence $\delta\ge37+4I$.
Together with \eqref{eq:z5-Dboot} this is
\[
 43+59I+(21\tau_s+20\tau_t)/23\le0.
\]

For E group $\{H_s,U_s,V_s\}$ and the two groups of $\mathcal G_t$,
with deficiencies $\Psi,\Omega$; let $u=u_{H_s}$, $w=u_{U_t}$.
The other resources $P,C5_t,\alpha_t$ are separate.
Raw-core joint bounds give
$3G+3h_s\ge1044-6R-3\Psi-3\Omega/2+3u+3w/2$.
Together with
$M+2M_t\ge R+162+2k+2I+\tau+2\delta$ and
$3X_s+5X_t\ge1112-8R-5\delta/2$, this gives
\[
 17I+(21\tau_s+19\tau_t)/23+25\delta/2\le378.
\]
Thus $\delta<31$, $I<23$, sizes $30\ldots39$, and every used center
is normal. Empty $\alpha_s$ lies below normal $H_s$, so $h_s\le X_s$.
The triple, its contained pair, and the two $g_t$ groups give
\[
 3g_s+6h_s+3g_t\ge744-9\tau'/4-6\Psi-3\Omega/2+6u+3w/2.
\]
In the normal transport use this and $5X_t\ge695-5R-5\delta/2$.
The largest distinct-resource coefficient is six, yielding
$2439/4+2I\le15R/4+5\tau/4+11\delta/2$.
The two rows of \eqref{eq:z5-weak} give $33+61I+\delta\le0$.

For F use all four groups $\mathcal G_s,\mathcal G_t$, with summed
$\Omega$, and choose $w=u_{U_s}$ in its group. Leave
$H_t,P,C5_t,\alpha_t$ separate. For any $f=2|Q|-\operatorname{span}(Q)-1$,
\[
 3G\ge6f-3\Omega/2+3w/2.
\]
Empty $H_s$ is HIGH above retained $U_s$, so
$h_s\ge74-\tau_t/23-w$. Using the raw core,
$M+3M_t\ge R+166+3k+2I+\tau+2\delta$, and
$3X_s+6X_t\ge1251-9R-3\delta$, the transport gives
\[
 18I+(21\tau_s+16\tau_t)/23+13\delta\le347.
\]
Thus $\delta<27$, $I<20$, sizes $31\ldots38$, and all used centers
are normal; now $h_s\ge74-w$.
The HIGH-alpha obstruction proved below applies to the retained triple
$H_t,\alpha_t,C5_t$, so the empty $\alpha_s$ is LOW and $h_s\le X_s$.
With $Q_{62}$, replace $3X_s+2h_s$ by $5h_s\ge370-5w$ and use
$6X_t\ge834-6R-3\delta$. Joint accounting gives
$1753/2+2I-\tau/2\le11R/2+6\delta$.
Then \eqref{eq:z5-weak} gives $553+33I+\delta\le0$.

\paragraph{Cards G, H, I, J: missing opposite cross.}
For G, group $\{H_s,U_s,V_s\}$, $\{C6_t,C7_t\}$ and the two
$g_t$ pairs, with deficiencies $\Psi,\Lambda,\Omega$; use
$u=u_{H_s}$ and $l=l_{C6_t}$ in the first two groups.
The endpoint transport is
\[
 M+3M_t+3X_s+6X_t+3G+h_s
 \le2R+762+3k+(2\tau_s+5\tau_t)/23+E-\delta,
\]
where $E=u_{H_s}+u_{H_t}+l_P+2l_{C6_t}+2l_{\alpha_t}$.
The raw joint bound
$3G+h_s\ge904-6R-\Psi-2\Lambda-3\Omega/2+u+2l$
gives $18I+(21\tau_s+18\tau_t)/23+14\delta\le425$.
It forces normality exactly as above ($\delta<31$, sizes $30\ldots39$),
and $h_s\le X_s$. For the final grouping leave $C6_t$ separate and
use the triple and its contained pair to get
\[
 3g_s+4h_s+3g_t\ge620-7\tau'/4-4\Psi-3\Omega/2+4u.
\]
The normal transport yields $2949/4+2I-3\tau/4\le21R/4+4\delta$,
and hence $15+67I+19\delta\le0$.

Cards H, I, J need no normality assumption. Their transports are
\begin{align*}
H: &2M+4X+3G+2h_t\le2R+905+3\tau/23+\delta,\\
I: &2M+M_t+4X_s+5X_t+3G
 \le2R+831+k+(3\tau_s+4\tau_t)/23+\delta,\\
J: &2M+2M_t+4X_s+6X_t+g_s+3g_t
 \le2R+757+2k+(3\tau_s+5\tau_t)/23+\delta.
\end{align*}
For H, use the two $g_s$ pairs, the triple $\{H_t,U_t,V_t\}$,
and $\{C6_s,C7_s\}$ to obtain
$3G+2h_t\ge496-3\tau'/2-2\delta$.
Use also $2M\ge2R+242+2\mathcal A$, $4X\ge1112-8R-2\delta$.
For I use
$2M+M_t\ge2R+k+279+3I+2\tau+3\delta$,
$4X_s+5X_t\ge1251-9R-5\delta/2$, and
$3G\ge372-3\tau'/2-3\delta/2$.
For J use
$2M+2M_t\ge2R+2k+316+4I+2\tau+4\delta$,
$4X_s+6X_t\ge1390-10R-3\delta$, and
$g_s+3g_t\ge248-\tau'-3\delta/2$.
The last bound uses the \emph{one-empty} form for $\{C6_t,C7_t\}$.
The respective intermediate and terminal rows are
\begin{align*}
H: &1683/2+2I+17\tau/46\le13R/2+3\delta,
 &71+17I+17\tau/23+7\delta&\le0,\\
I: &1935/2+3I+(17\tau_s+15\tau_t)/46\le15R/2+2\delta,
 &75+21I+(17\tau_s+15\tau_t)/23+11\delta&\le0,\\
J: &1128+4I+(20\tau_s+18\tau_t)/23\le9R+3\delta/2,
 &24+26I+(40\tau_s+36\tau_t)/23+15\delta&\le0.
\end{align*}

\paragraph{Card K: the last retained-pair1 family.}
The fixed HIGH set consists of $C_{11},C_{12},C_{13},C_{14}$ on both
signs and $C_{22}^s,C_{23}^s,C_{24}^s$. The LOW set is
\[
 \{C_{55}^t,C_{56}^t,C_{57}^t,C_{58}^t,C_{66}^t,C_{67}^t,C_{68}^t\}
 \cup\{V_i^v: i=5,6,7,8,\ v=s,t\}.
\]
The empty own centers are HIGH and the empty opposite crosses and beta
components are LOW, so these signs follow by adding positive or negative
outer magnitudes. The 15 LOW labels touch 12 baseline groups; the 11
HIGH labels lie outside those groups. Thus $e=11+15-12=14$.
Put $\epsilon=\delta-14$. If $m=\max(r_s,r_t)$, a BAD deficiency
exceeds $\epsilon$, since
$74-m-(51-I-14)=(m-35)^2+2>0$. All used endpoints are normal.
Their transport is
\[
 M+3M_t+3X_s+6X_t+3G\le2R+762+3k+E-\delta,
 \quad E=w+u+p+2c+2q,
\]
where $(w,u,p,c,q)=(u_{U_s},u_{H_t},l_P,l_{C6_t},l_{\alpha_t})$.
Group the four ordinary pairs, summed deficiency $\Omega$, and leave
$H_t,P,\alpha_t$ separate. Joint clipping in $\{U_s,V_s\}$ and
$\{C6_t,C7_t\}$ gives
$3G\ge372-3\tau'/2-3\Omega/2+3w/2+3c/2$.
Use $M+3M_t\ge R+166+3k+2I+\tau+2\delta$ and
$3X_s+6X_t\ge1251-9R-3\epsilon$. Their combination yields
\[
 45I+29\delta\le467.
\]
Since $I$ is zero or at least two and $\delta\ge14$ is integral,
$I=0$ and $\delta\in\{14,15,16\}$, so $\epsilon\le2$.
For native zero-loss endpoints the same transport gives
$S\ge8M+3X_s+6X_t+3G+5$. With at most two clipped native points,
the third extreme source magnitude $R_v+b_v$ is still present.
Put $H_2=\max(e_s-b_s,e_t-b_t)$. At most two endpoint positions can
have clipping, each with coefficient at most two, so rigorously
\[
 S\ge8M+3X_s+6X_t+3G+5-4H_2.
\]
Choose $j$ attaining $H_2=e_j-b_j$, with opposite $i$.
The weighted source inequality gives
\[
 6M-4H_2=6R+2(L+e_i)+2e_i+4b_j
 \ge6R+332+4I+2\tau+4\delta.
\]
Use this, the stars and $3G\ge834-6R-3\epsilon/2$ in the preceding
native row. Then $R+\delta\le124-I$ gives
\[
 15I+2\tau+16\delta\le82+9\epsilon/2.
\]
Its left side is at least 224, its right side at most 91. This excludes
card K and finishes all thirteen retained-pair1 cases.

\subsection{A0 with both pairs empty: the remaining three cards}
After full reflection, the exact holes are
$\{\mathrm{pair2},\mathrm{pair1},O_s,O_t,\eta\}$ with
$\eta=H_s,C5_s$, or $\alpha_s$. There is no pair owner in these cases.
The thirteen retained singletons consist of six on $s$ and seven on $t$.
Define the fixed-label sets
\[
\begin{split}
\mathcal H_0&=\{C_{11}^v,C_{12}^v,C_{13}^v,C_{14}^v:v=s,t\},\\
\mathcal L_0&=\{C_{55}^v,C_{56}^v,C_{57}^v,C_{58}^v,
                       V_5^v,V_6^v,V_7^v,V_8^v:v=s,t\}.
\end{split}
\]
The empty $O$ centers are HIGH, the empty pair components LOW.
If $\alpha_s$ is empty, it lies below retained $H_s$ and is LOW;
if $C5_s$ is empty it lies below retained $C6_s$ and is LOW.
Thus $\mathcal H_0$ is HIGH, $\mathcal L_0$ LOW. The latter has 16
labels meeting exactly 14 baseline groups. The following table specifies
the full initial ledgers and their endpoint/star budgets; all numbers
follow from the 26 explicit groups of the common lemma.
\begin{center}\small
\begin{tabular}{c|l|l|r r r r}
$\eta$&HIGH additions&LOW additions to $\mathcal L_0$&$|\mathcal H|$&$|\mathcal L|$&$g$&$e$\\\hline
$H_s$&$C_{22}^s,C_{23}^s,C_{24}^s$&$F_{3,4,5}^{s,t}$&11&22&20&13\\
$C5_s$&none&$C_{66}^s,C_{67}^s,C_{68}^s,F_{3,4,5}^{s,t}$&8&25&20&13\\
$\alpha_s$&none&$U_{5,6,7,8}^s,F_{3,4,5}^{s,t},V_{2,3,4}^s$&8&29&23&14
\end{tabular}
\end{center}
Here $F_3=g-d$, $F_4=g-d-h$, $F_5=g-d-2h$ are negative by the
leading inequality. In the alpha case $h_s\le X_s$, so
$d_s>g_s+2h_s$, proving $V_2^s,V_3^s,V_4^s<0$.
Adding the six $F$ labels touches six new groups in the first two rows,
but only five in the alpha row because its $U_8^s$ already touched
$\{C_{44}^s,U_8^s,F_5^s\}$. The three $V$ additions there touch
only one new group. These facts prove every count in the table without
assuming normality. For all three rows
\[
 \epsilon=\delta-e\ge0,\qquad
 74-m-(51-I-e)=(m-35)^2+e-12>0
\]
with $m=\max(r_s,r_t)$. Therefore all retained centers are normal.
The same original cover licenses the $\epsilon$ versions of joint groups
and stars; structural equations continue to use $\delta$.

For normal retained same-target centers $B,D$, their actual group
$\Gamma$ satisfies the following useful packing inequality:
\begin{equation}\label{eq:z5-pack2}
 B+D\ge2(74-r_v)+f-2\Gamma.
\end{equation}
Indeed, put $m=\min(B,D)$, $q=|B-D|$, and let $l$ be the lower
clipping of the smaller center. Joint accounting gives
$\Gamma\ge l+f-q$, while $m\ge74-r_v-l$ and $l\le\Gamma$.
This proves \eqref{eq:z5-pack2}, including coincident numerical centers.

Group all six retained side-$s$ resources with its leading. Its actual
union deficiency is at most $\epsilon$, so
\begin{equation}\label{eq:z5-side}
 S_s\ge K+468-\epsilon,\qquad S_t\le K+630-e.
\end{equation}
This is a replacement in the whole original cover. A set being contained
in a target alone would not give the required upper bound on deficiency.
Write $A=\alpha_s$, $B=\alpha_t$, $D=C5_t$.
The identity $H_t=2(a+b)+G+A+B+D$, its normal upper endpoint, and
\eqref{eq:z5-side} give
\begin{equation}\label{eq:z5-sideT}
 M+2X+G+A+B+D\le R+578+r_s-e+u_{H_t}.
\end{equation}
If $A$ is retained, lower it normally and pack $B+D$ by
\eqref{eq:z5-pack2}; the three resource groups are disjoint from $H_t$.
Then
$M+2X+G+f\le R+494-e+2\epsilon$.
With $M\ge R+121+\mathcal A$, the four-star row, and
$G\ge2f-\epsilon/2$, this yields
\[
 1269+4I+\tau+4\delta+4e\le13R+14\epsilon.
\]
Consequently the exact terminal card for either retained-$A$ case is
\begin{equation}\label{eq:z5-corrected17}
 \boxed{17I+\tau+3\delta+18e\le343.}
\end{equation}
In particular the coefficient is $3$. To finish we prove $e=17$ within
each of these two actual families.

\paragraph{The $H_s$ hole.}
First use only $e=13$. If $A\ge b$, use
\eqref{eq:z5-pack2} at $B,D$ in \eqref{eq:z5-sideT} to obtain
\[
 M+M_t+2X_s+3X_t+G+f\le R+499+r_t-e+2\epsilon.
\]
Use the weighted source row, stars, and $G\ge2f-\epsilon/2$.
After $R+\delta\le124-I$ the result is
$25I+\tau+9\delta+20e\le155$.
Its left side is at least $377$, so $A<b$.
The five additional LOW labels are
\[
 U_7^s=A-t_3,\quad U_8^s=A-b,\quad
 V_2^s=2A-2b-g_t-B-C5_s,\quad
 V_4^s=A-a-2b-G-B-C5_s,\quad F_5^s.
\]
Starting with the base HIGH11, LOW16, touched14 ledger, the first four
are negative by $A<b$ and normal $B,C5_s>0$, and $F_5^s<0$ by the
leading inequality. $V_2^s,V_4^s$ lie in already touched groups,
$U_7^s$ in none, and $U_8^s,F_5^s$ touch one new group.
Thus HIGH11, LOW21, touched15 gives $e=17$.
Now \eqref{eq:z5-corrected17} and $\delta\ge e$ give $357\le343$.

\paragraph{The $C5_s$ hole.}
Use only $e=13$ and suppose $A\ge b/2$. The same packing gives
\[
 M+M_t/2+2X_s+5X_t/2+G+f\le R+499+r_t-e+2\epsilon.
\]
Use half each of the two source bounds in \eqref{eq:z5-weighted},
the stars and the gaps. Simplification gives
$21I+\tau+6\delta+19e\le259+2r_t$.
At $e=13$, this says
$21I+\tau+6\epsilon\le2r_t-66$.
Since $r_t\le41$, $I$ is zero or at least two, and $\epsilon$ is
integral, necessarily $I=0$, $r_t\in\{34,35\}$ and $\epsilon=0$.
With zero native endpoint loss, upper $H_s,H_t$ and lower
$C6_s,A,B,D$ give
\[
 S\ge7M+4X_s+3X_t+2G+5.
\]
The stars, raw-core gaps and $M\ge R+121+\mathcal A$ imply
$999+5I+5\tau+6\delta\le8R$, and then
$7+13I+5\tau+14\delta\le0$, impossible.
Therefore $A<b/2$. Put $W=C6_s>0$; the exact identity
$d_s=2a+b+G+W$ now makes the four new labels
\[
 V_2^s=2A-b-g_t-W,\quad V_4^s=A-a-b-G-W,\quad
 U_7^s=A-t_3,\quad U_8^s=A-b
\]
strictly negative. In the displayed direct ledger, the first, second and
fourth already have touched groups (the fourth via $F_5^s$), and the
third has none. Thus HIGH8, LOW29, touched20 gives $e=17$.
Again the corrected row \eqref{eq:z5-corrected17} gives $357\le343$.

\paragraph{The $\alpha_s$ hole: the fixed-credit argument.}
We give the full tail and order arguments because scalar $\delta\ge38$
alone would not license the reduced-budget proof. Start with $e=14$.
Then $h_s\le X_s$, and separately the leading, stars and ordinary pairs
give
\begin{equation}\label{eq:z5-tailbase}
 d_v\ge104+g_v+2X_v,\quad X_v,g_v\ge29,\quad h_s\ge29.
\end{equation}
For the last bound use the retained pair $\{H_s,U_s\}$ alone; no
independent allocation is added to a group already used elsewhere.

For either physical sign put $T=v_1+v_2+v_3+v_4$ and
$Q=2v_1+v_3+M_v$, so $Q-T=d_v-X_v$.
Every integer $u\in(d_v-X_v,S_v]$ has $T+u\in4A_v$.
Call $v_1,v_2$ top entries. With at most one top entry, a sum is at
most $v_1+3v_3<Q$ by \eqref{eq:z5-tailbase}. With exactly two,
it is at most $v_1+v_2+2v_3<Q$ unless both are $v_1$; then the
other two must be $v_3$ or $v_4$ to exceed $Q$, giving pullbacks
$d_v-g_v,d_v,d_v+g_v$. Exactly three top entries give
$O_v,H_v,\alpha_v,\beta_v$ plus either a source entry or one of
$v_3,v_4,-o_1,-o_2,-o_3,-o_4$. Exactly four give
$(j-1)d_v+2h_v+g_v$, $j=0,\ldots,4$.
This exhausts all representations and leaves exactly 32 non-source
labels: 24 three-top, five four-top and three two-top types. In the
original inventory they are
\[
 \{C_{13},C_{14},C_{15},C_{16},C_{17},C_{18},
 C_{23},C_{24},C_{25},C_{26},C_{27},C_{28},
 U_{3..8},V_{3..8},V_2,U_2,C_{22},C_{12},C_{11},C_{44},C_{34},C_{33}\}.
\]
Numerical coincidences can only decrease their coverage.

On side $s$ only $H_s+D_s$ among the four source families remains,
ending at $B_s=H_s+M_s=d_s-X_s+h_s\le d_s$.
The tail $(d_s-X_s,S_s]$ has at most $X_s$ points from this family
and at most 26 non-source labels: all six beta types lie below its
start. Hence $S_s\le d_s+26$. The identity
$O_s-t_2-d_s=C5_s+s_3+h_s>26$ excludes
$C_{16}^s,C_{17}^s,C_{18}^s$ above this endpoint.
For $d_s+1,d_s+2$, every other listed type is at most $d_s$ or at
least $d_s+29$, except the single score $Y=O_s-t_1$.
(The alpha and low four-top scores obey $h_s+g_s,2h_s+g_s<d_s$;
positive $H_s$ scores and $d_s+g_s$ have gaps at least 29.)
Thus $S_s\le d_s+1$. Above $B_s$ only $d_s-g_s,d_s,Y$ remain;
the first two are at least 29 apart, so they and $Y$ cannot cover three
consecutive integers. Therefore
\[
 S_s\le H_s+M_s+2.
\]
This proves seven new HIGH labels
$C_{22}^s,C_{23}^s,C_{24}^s,C_{16}^s,C_{17}^s,C_{18}^s,C_{33}^s$.

On side $t$ every actual nonfixed resource ends by
$B_t=H_t+M_t$: the leading ends at $d_t-g_t-2X_t$; own2 and own3
end at $d_t-X_t,d_t-g_t-X_t$; alpha ends at $B_t-d_t$; the three
retained crosses have centers below $H_t$; the remaining components are
empty. In the 32-label list, six alpha types, six beta types and four
negative-$H_t$ types lie below $B_t$. Consequently $S_t\le B_t+16$.
As $X_t\ge29$, the labels $C_{22}^t,C_{23}^t,C_{24}^t$ are HIGH.
The identity $O_t-s_2-B_t=C5_t+M_t+2X_t+g_t>16$ proves
$C_{16}^t,C_{17}^t,C_{18}^t$ HIGH as well.
We have obtained HIGH21, LOW29, touched23, hence $e=27$, without
using the new credit to prove any of its signs.

Put $B=\alpha_t$, $C=C5_s$, $D=C5_t$. At $e=27$, grouping whole
physical targets also gives $S_s\le K+552-e$, $S_t\le K+630-e$.
The normal upper endpoints of $H_s,H_t$ therefore give
\begin{align*}
 M_t+X_s+2X_t+G+h_s+B+C&\le R+500+r_t-e+u_{H_s},\\
 M_t+X_s+2X_t+G+h_s+B+D&\le R+578+r_s-e+u_{H_t}.
\end{align*}
If $B\ge a$ or $D\ge a$, lower the other retained center normally
to obtain the common weaker row
$M+2X+G+h_s\le R+573-e+\epsilon$.
The triple $\{H_s,U_s,V_s\}$ and the two $g_t$ groups yield
$G+h_s\ge186-\tau'/2-\epsilon$.
Together with the stars, the sharp mass bound, and
$2R+I+\delta\le191$ proved below, this gives
\[
 11I+2\tau-\delta+16e\le315.
\]
Since $\delta\le47-I-\tau$, its left side at $e=27$ is at least
$385$, impossible. Hence $B,D<a$.
The nine new LOW labels
\[
 U_{5,6,7,8}^t,\quad V_3^t,V_4^t,\quad C_{66}^t,C_{67}^t,C_{68}^t
\]
follow directly: for example
$V_3^t=B-a-b-g_s-h_s-D<0$ and
$V_4^t=B-a-b-G-h_s-D<0$.
Only the groups $\{C_{36}^s,U_4^s,U_5^t\}$ and
$\{C_{35}^t,V_3^t,C_{46}^s\}$ are newly touched. Thus
HIGH21, LOW38, touched25 gives $e=34$.

Next $C<b$. If $C\ge b$, it is the smallest retained ordinary center
on $s$; all nonfixed ordinary hits start at $C-M_t$, and the leading
starts even higher, at
$\lambda_s-2M_t=C+2X_t+B+g_t>C$.
The prefix $[1,C-M_t-1]$ is therefore fixed-only.
Already at $e=14$, $X_t\ge29$, so it contains at least 28 integers.
The tail proof made 11 side-$s$ labels HIGH, and the initial $e=14$
ledger made 18 side-$s$ labels LOW. Only $55-11-18=26$ fixed labels
remain, a contradiction. Thus $C<b$, proving four further LOW labels
\[
 C_{66}^s=C-t_2,\quad C_{67}^s=C-t_3,\quad C_{68}^s=C-b,\quad
 C_{35}^s=C-D-2b-g_t-h_s.
\]
Their groups were already touched. The final ledger is therefore
\begin{equation}\label{eq:z5-e38}
 |\mathcal H|=21,\quad |\mathcal L|=42,\quad g(\mathcal L)=25,
 \qquad e=38,\quad\epsilon=\delta-38\ge0.
\end{equation}
Every ledger uses the original fixed inventory and hence preserves the
same-cover group and star licenses.

\paragraph{The $\alpha_s$ hole: the final transport.}
We supply the extra source inequalities used above and here. Delete one
source maximum at a time. The 77-point remainder has central interval
and fringe bound $K\le2R+e_t+c_s+95$ and its reflection; reinsertion
costs 78. Combining with
$2K=127+\mathcal A+4R+J+U+2C$ gives
\[
 |J_s-J_t|+U\le63-\mathcal A,\qquad J\ge109+\mathcal A,
 \qquad J_v\ge26+\mathcal A,\quad e_v\ge29+\mathcal A.
\]
Here $U\ge6$. For the penultimate core magnitude $p_v$,
$p_v\le2r_v-1+\tau_v-R_v$. Delete both source maxima and use the
central interval with endpoints $R_v+c_v+p_v$. The four remaining
points beyond each smaller core contribute at most 20 pair labels and
reinsertion costs 155. Hence $K\le C+312+\tau$.
The mass equation and $J\ge109+\mathcal A$ give
\begin{equation}\label{eq:z5-radial}
 2R+I+\delta\le191.
\end{equation}

Group $\{B,D\}=\{\alpha_t,C5_t\}$ with deficiency $\Gamma$,
$\{V_s,U_s,H_s\}$ with $\Psi$, and the two $g_t$ pairs with
summed $\Omega$; leave $H_t$ and $C=C5_s$ separate.
These eleven singleton labels are distinct. The joint packing argument
and subgroup domination give
\[
 2B+D\ge3(74-r_t)+f-3\Gamma,\quad
 g_s+h_s\ge124-\tau'/4-\Psi,\quad 2g_t\ge2f-\Omega.
\]
The exact identities
$H_s=a+2b+G+h_s+B+C$ and
$H_t=a+2b+G+h_s+B+D$ and the normal endpoints now yield
\[
 2M_t+2X_s+4X_t+2(G+h_s)+f
 \le2R+836+2r_t-\delta+3\Gamma+l_C+u_{H_s}+u_{H_t}.
\]
In this same grouping $u_{H_s}\le\Psi$, so the total loss after
inserting the gap bounds is at most
$3\Gamma+3\Psi+\Omega+\Gamma_C+\Gamma_{H_t}\le3\epsilon$.
Thus
\[
 2M_t+2X_s+4X_t+434-5\tau'/4
 \le2R+836+2r_t-\delta+3\epsilon.
\]
Use the separately valid stars
$2X_s+4X_t\ge834-6R-2\epsilon$ and
$M_t\ge r_t+29+\mathcal A$. This gives
\[
 1615+8I+3\tau+20e\le27R+8\delta.
\]
Double this row and apply \eqref{eq:z5-radial}:
\begin{equation}\label{eq:z5-alpha-final}
 \boxed{43I+6\tau+11\delta+40e\le1927.}
\end{equation}
There is no additional $M$ in the transport. With
$\delta\ge e\ge38$, its left side is at least $1938$, impossible.
This finishes all three pair1-empty A0 cases, with all BAD centers already
excluded by their initial fixed-credit argument.

\subsection{A1: endpoint transports, the exceptional pattern, and two final families}
We retain the actual beta owner $s$ proved above. This subsection first
proves four transport obstructions and one branch-independent endpoint
lemma, then explicitly routes every alternative in
\eqref{eq:z5-a1list}. A quadratic minimum over real numbers used below
is stronger than the required minimum over integer physical sizes.

\paragraph{The HIGH-alpha obstruction.}\label{par:z5-highalpha}
Suppose $\alpha_s\ge S_s+R_t+1$, with $H_t,\alpha_t,C5_t$ retained.
The identity $\alpha_s+C5_t=d_t-2(a+b)-G$ and the lower $C5_t$
endpoint give
\[
 d_t\ge S_s+R_t+2(a+b)+G+75-r_t-\tau_t/23-l_{C5_t}.
\]
The span $H_t-\alpha_t=d_t$ gives
$S_t\ge d_t+78-\tau/23-u_{H_t}-l_{\alpha_t}$.
These are three distinct resources, so the total loss is at most
$\delta$. Use the source mass, $K\le M+R+21$, the four-star row,
and $S_s\ge T_r+\tau_s+4R_s+10$, $r=r_s$. Successively,
\[
\begin{split}
 2S_s+R_t+2X+G&\le2R+987+r_t+(\tau+\tau_t)/23,\\
 2T_r+8r+5I+4\delta-2\tau/23&\le1100,\\
 141r^2-7935r+114895+95\delta&\le0.
\end{split}
\]
The last quadratic has real minimum $612185/188>0$. Thus the
assumed HIGH alpha is impossible. This statement requires no beta
retention and is the obstruction used in card F.

\paragraph{Both pairs retained, with opposite owners.}
Write $a_0$ for the actual pair1 owner and $b_0$ for the opposite actual
beta owner, and put $r=r_{b_0}$. (These subscripts are physical signs.)
The identities
\[
 P_{a_0}=d_{a_0}-2d_{b_0}-\beta_{b_0}-2(s_4+t_4)-G,
 \qquad V_{a_0}=d_{a_0}-(a_0)_4-g_{a_0}
\]
use the retained lower endpoints of $P_{a_0},\beta_{b_0}$ and upper
endpoint of $V_{a_0}$. Two leading-$b_0$ inequalities then give
\[
 S_{b_0}+2g_{b_0}+4X_{b_0}+(a_0)_4+2(b_0)_4+G-g_{a_0}
 \le946+r+\tau/23+\tau_{b_0}/23+\delta.
\]
For the loss accounting the three endpoint losses and one leading loss
sum to at most $\delta$, so using the leading twice leaves at most one
$\delta$ after the structural subtraction.
Half of the four-star plus five halves of the side-$b_0$ star yields
$X_{a_0}+6X_{b_0}\ge973-7R-3\delta$.
Use the positive mass bound and $g_{b_0}\ge139-R-\delta/2$ to get
\[
 466+T_r+4r\le9R+\tau/23+11\delta/2,
 \qquad440r^2-28566r+466105+164\delta\le0.
\]
Its quadratic minimum is $1082111/440>0$, so opposite owners are
impossible. No $U_{a_0}$ endpoint was required.

\paragraph{Both pairs retained, with the same owner and retained alpha.}
Let both owners be $s$, with $\alpha_s$ retained, and put $r=r_t$.
The identities
$\alpha_s-P_s=\mathcal B+2d_t+2(a+b)+G$,
$\beta_t=\alpha_t-d_t$, and $\alpha_t>-b$ give, using the
lower $\beta_s$ endpoint, the upper $\alpha_s$ endpoint and lower
$P_s$ endpoint,
\[
 S_t+g_t+2X_t+2a+b+G\le K+1015-r+\tau/23+\tau_s/23.
\]
One leading-$t$ loss and these three native resource losses fit one
$\delta$. Put $u=M-R-8$. With the mass bound,
$2X_s+3X_t\ge695-5R-3\delta/2$, and the depth-five one-empty pair
bound $g_s+2g_t\ge180-\tau/2-\delta$, this implies
\[
 T_r+4r+u\le147+5R+5\delta/2+27\tau/46.
\]
Already the weaker $u\ge89+I+\tau+\delta$ gives
$13r^2-819r+12924+19\tau/23+7\delta\le0$.
The quadratic minimum is $99/4>0$, a contradiction.

\paragraph{A branch-independent four-endpoint lemma.}\label{par:z5-estar}
Suppose all eight lower own/cross singletons are retained, pair1 has
actual sole owner $s$, and $\alpha_t$ is retained. The exact identity
\[
 U_s=P_s+\alpha_t+C6_t+2a+4b+2G
\]
uses four distinct retained native resources. Three lower endpoints and
one upper endpoint give, with $r=r_t$ and total loss at most $\delta$,
\begin{equation}\label{eq:z5-estar-transport}
 S_t+2M_t+2X_s+4X_t+2G
 \le914+2r+2R+(\tau_s+3\tau_t)/23.
\end{equation}
For this lemma use the alternative span bound
$\operatorname{span}(Q_{62})\le61+\tau/5$; it follows directly from
the five terminal monotone deviations on each side. Put $f_0=62-\tau/5$.
The stars give $2X_t\ge2f_0-\delta$ and $2X\ge4f_0-\delta$;
the four ordinary groups give $2G\ge4f_0-\delta$. Therefore
\[
 2X_s+4X_t+2G\ge620-2\tau-3\delta.
\]
All three copies of $\delta$ are displayed, since the star and ordinary
groupings are separately valid, not jointly disjoint. With
$S_t+2M_t\ge T_r+\tau_t+6r+18$, $R\le69+\tau$, and
$\tau+\delta\le47-I$, \eqref{eq:z5-estar-transport} yields
\begin{equation}\label{eq:z5-estar-final}
 23T_r+92r+95I\le13987.
\end{equation}
But for integer $r$,
\[
 23T_r+92r+95(r-34)(r-35)
 =15355+\tfrac12\{213(r-30)^2-123(r-30)\}\ge15355.
\]
This contradiction uses no endpoint, loss, capacity or sign of
$\beta,\alpha_s,H_s$.

If pair1 is retained in A1, opposite owners and same owners with retained
$\alpha_s$ have just been excluded. Otherwise the three-active-center
identity forces $C5_t$ empty. The exact remaining holes are
$O_s,H_s,O_t,\alpha_s,C5_t$; every other resource is retained.
This is the exceptional pattern sometimes denoted $E^*$, and the
branch-independent lemma excludes it. Thus every retained-pair1 A1 case
has been treated.

\paragraph{Pair1 empty: the four excluded placements.}
The six possibilities are exactly \eqref{eq:z5-a1list}.
If $\eta=\alpha_s$, its emptiness and its position above active
$\beta_s$ force it HIGH, while $H_t,\alpha_t,C5_t$ are retained.
The HIGH-alpha obstruction applies.
For $\eta=C5_s$ or $U_s$, the retained endpoints
$\alpha_s,\beta_s,H_t,\alpha_t,C5_t$ and two leading-$s$ bounds
use the identity $\beta_s+C5_t=d_t-d_s-2(a+b)-G$ to give
\[
 3g_s+g_t+2M+6X_s+2X_t\le863+3\tau/23+\delta.
\]
The endpoint loss expression is
$u_{\alpha_s}+2l_{\beta_s}+u_{H_t}+l_{\alpha_t}+l_{C5_t}+2\ell_s$,
whose coefficients are at most two on distinct original resources.
Use the three retained ordinary groups
$\{C6_t,C7_t\}$, $\{U_t,V_t\}$, $\{C6_s,C7_s\}$.
The raw core and a depth-five core respectively give
\[
 3g_s+g_t\ge556-4R-3\delta,\qquad
 3g_s+g_t\ge240-2\tau/3-3\delta.
\]
The source bound and stars give
$2M+6X_s+2X_t\ge1330-6R+2I+2\tau-\delta$.
Combine separately with the two gap rows and use the two bounds on $R$:
\[
 8I+83\tau/69+\delta\le37,\qquad
 333+2I\le187\tau/23+5\delta.
\]
Eliminating $\tau$ gives $6882+4654I+146\delta\le0$.
No missing own-$s$ pair has been used.

For $\eta=H_t$, use lower $\alpha_t,C5_s,\beta_s,C5_t$ and upper
$\alpha_s,V_t$. Their literal relations and endpoints give
\begin{align*}
 d_s&\ge2(a+b)+G+79-\tau/23-l_{\alpha_t}-l_{C5_s},\\
 d_t&\ge d_s+2(a+b)+G+79-\tau/23-l_{\beta_s}-l_{C5_t},\\
 S_s&\ge d_s+78-\tau/23-u_{\alpha_s}-l_{\beta_s},\\
 S_t&\ge d_t-t_3+73-r_s-\tau_s/23-u_{V_t}.
\end{align*}
The six distinct resources have endpoint coefficient at most two.
Adding the last two rows and substituting the first two yields
\[
 4M_s+3M_t+6X_s+5X_t+3g_s+2g_t
 \le2R+752+r_s+5\tau/23+\delta.
\]
Use, respectively,
$4M_s+3M_t\ge3R+331+r_s+3I+3\tau+3\delta$,
$6X_s+5X_t\ge1529-11R-3\delta$, and
$3g_s+2g_t\ge695-5R-3\delta/2$.
Then
\[
 18I+64\tau/23+25\delta/2\le57,\qquad
 768+3I\le281\tau/23+5\delta/2.
\]
Elimination gives $4418+700I+447\delta\le0$.
Thus only $\eta=\alpha_t$ and $\eta=C5_t$ remain.

\paragraph{The final $\alpha_t$-empty family.}
The exact holes are $O_s,H_s,O_t,\mathrm{pair1},\alpha_t$.
Use upper $\alpha_s,H_t$ and lower $\beta_s,C5_t,C5_s$, with
coefficients on lower beta and $C5_s$ equal to two. The identities
\[
 C5_s=d_s-2a-b-G-h_t,\qquad
 \beta_s+C5_t=d_t-d_s-2(a+b)-G
\]
and the spans $\alpha_s-\beta_s=d_s$, $H_t=d_t+h_t-b$ give
\begin{equation}\label{eq:z5-a1alphaT}
 4M_s+M_t+6X_s+3X_t+3G+3h_t
 \le2R+762+3r_s+(2\tau+3\tau_s)/23+\delta.
\end{equation}
The five resources are distinct, so the pre-subtraction loss is at most
$2\delta$. Group $\{U_s,V_s\}$ and $\{C6_t,C7_t\}$ with summed
$\Omega$, and $\{H_t,U_t,V_t\}$ with $\Psi$.
These seven labels are distinct, with $\Omega+\Psi\le\delta$.
The raw core gives $3G+3h_t\ge1044-6R-3\delta$.
Put $\xi_s=R_s-r_s\ge0$. The weaker source bound with constant 109
already gives
$4M_s+M_t\ge R+121+3r_s+3\xi_s+I+\tau+\delta$;
the stars give $6X_s+3X_t\ge1251-9R-3\delta$.
Together these imply
\[
 17I+3\xi_s+(18\tau_s+21\tau_t)/23+10\delta\le330.
\]
Thus $\delta\le33$, $I<20$, sizes $31\ldots38$, and every used
resource has at least 45 hits, exceeding the BAD maximum 42.
All used centers, including the actual sole beta component, are normal.
In particular empty $\alpha_t=H_t-d_t<S_t$ cannot be positive
because $0\in D$; nor can $\alpha_t+M_t$ be positive because it
would be a target hit. Hence
\begin{equation}\label{eq:z5-a1alpha-low}
 \alpha_t\le-M_t,\qquad h_t\le X_t.
\end{equation}
This derives LOW for the entire family, after normality and before use.

Repeat \eqref{eq:z5-a1alphaT} with normal endpoints, removing its
$\tau/23$ error. The same groups with $Q_{62}$ give
$3G+3h_t\ge558-3\tau'/2-3\delta$.
Using the mass and star rows above gives
$2129+2I+6\xi_s\le17R+\tau+12\delta$, and then
\begin{equation}\label{eq:z5-a1small}
 41I+12\xi_s+13\delta\le183.
\end{equation}
The contained pair $\{H_t,U_t\}$ has deficiency at most $\Psi$,
so $h_t\ge f-\Psi$. Together with the triple and the two $g_s$
pairs it gives
$3g_s+3g_t+6h_t\ge744-9\tau'/4-6\delta$.
Use \eqref{eq:z5-a1alpha-low} and the side-$s$ star in the normal
transport to obtain
$3127+4I+12\xi_s\le19R+5\tau+36\delta$.
The source inequalities imply
\[
 417+61I+24\xi_s\le19\delta.
\]
Multiply this by 13 and \eqref{eq:z5-a1small} by 19. After division
by 12 the contradiction is
\begin{equation}\label{eq:z5-a1alpha-final}
 162+131I+45\xi_s\le0.
\end{equation}

\paragraph{The final $C5_t$-empty family.}
The exact holes are $O_s,H_s,O_t,\mathrm{pair1},C5_t$.
The eleven HIGH labels are $\mathcal H_0$ together with
$C_{22}^s,C_{23}^s,C_{24}^s$.
The eleven LOW labels are $C_{55},C_{56},C_{57},C_{58}$ on both
signs and $C_{66}^t,C_{67}^t,C_{68}^t$.
Empty crosses are LOW and empty own centers HIGH as already proved.
The LOW set touches nine baseline groups, so $e=11+11-9=13$.
Put $\epsilon=\delta-13$. Normality follows because
$74-m-(38-I)=(m-35)^2+1>0$; it applies also to sole-owned beta.

Write $b_0=\beta_s$, $z_0=C5_s$, $q_0=\alpha_t$, $w_0=C6_t$.
The identity $d_s=2(a+b)+G+z_0+q_0$ gives
\begin{equation}\label{eq:z5-tailbeta}
 b_0+s_3<d_s-X_s.
\end{equation}
To prove it, suppose the reverse. Then
$b_0\ge M_s+2M_t+2X_t+g_t+z_0+q_0$ and
$\alpha_s\ge3M_s+4M_t+2X_s+4X_t+g_s+2g_t+2z_0+2q_0$.
Use upper $\alpha_s$, lower $z_0,q_0$, and the source mass of $t$.
Their endpoint losses have expression $E=2l_{z_0}+2l_{q_0}+u_{\alpha_s}$.
The four ordinary groups, disjoint from these resources, give
$g_s+2g_t\ge417-3R-\Omega$ and $E+\Omega\le2\epsilon$.
The weighted source bound and stars now yield
\[
 T_{r_t}+4r_t+12I+\tau+\tau_t+9\delta\le678,
\]
contrary to \eqref{eq:z5-q717}. This proves \eqref{eq:z5-tailbeta}.

Use the same actual 32-type upper-tail classification as above on side
$s$, with threshold $d_s-X_s$. The $O_s,H_s$ source families are
empty; by \eqref{eq:z5-tailbeta} so is the beta contribution to this
tail, and all six beta fixed types lie below it. Only the retained
$\alpha_s$ source family and at most 26 fixed labels remain. Hence
\[
 S_s\le d_s-X_s+104-\Gamma_{\alpha_s}.
\]
The inequality is automatic if $S_s\le d_s-X_s$.
Use the normal upper $\alpha_s=d_s+b_0$ and lower beta endpoints,
with $l_{\alpha_s}+u_{\alpha_s}\le\Gamma_{\alpha_s}$, to get
\begin{equation}\label{eq:z5-xs26}
 X_s\le26+l_{\beta_s}-l_{\alpha_s}.
\end{equation}
For the final local star, group only
$\{\mathcal L_s,V_s\}$ and $\{\mathcal L_t,C7_t\}$.
Their contained core translates have separation $X_s$, and the groups
are disjoint from beta and alpha. Thus the same reduced cover gives
\[
 2X_s+2R+\epsilon-\Gamma_\beta-\Gamma_\alpha\ge278.
\]
Insert \eqref{eq:z5-xs26}, use $l_\beta\le\Gamma_\beta$ and all
nonnegative losses, and then $\Gamma_\beta\le\epsilon$:
\[
 278\le52+2R+2\epsilon
       =26+2(R+\delta)\le274-2I.
\]
This excludes the second final family and completes A1.

\subsection{Branch B: use the actual C5 owner}
Now $\mathcal B=\beta_s+\beta_t>0$, so the second paired whole unit is
$\{C5_s,C5_t\}$ and both beta components are \emph{singleton} units.
The two $O$ holes leave three more. On each sign at least one of
$\beta_v,H_v$ is empty, so only one extra hole can lie outside these
two mandatory choices. A nonintegral lower coordinate would force two
such holes, impossible. With lower coordinates integral, a nonintegral
$v_2$ forces the three singleton holes $H_v,\alpha_v,\beta_v$ plus
the opposite mandatory hole. Thus all outer coordinates are integral.
Empty $V_v$ would force $U_v,H_v$ as well; empty $C7_v$ would force
$C6_v$ as well. Both contradict the available hole count, so the four
extreme leaves are retained.

The C5 pair is retained: otherwise it and the two mandatory beta/H holes
would exhaust the extras, leaving pair1 retained; its actual owner would
force C5 activity. Its owner is sole, since
\[
 C5_s+C5_t=-\mathcal B-2(s_3+s_4+t_3+t_4)<-4M,
\]
whereas two active centers would have sum greater than $-M$.
Normalize $s$ to this actual C5 owner. If pair1 is retained, it has the
same owner, by
$P_s+C5_t=-\mathcal B-2(s_3+s_4+t_3+t_4)-d_t$ and reflection.
The exhaustive dichotomy is whether this owner's $\alpha_s$ is retained
or empty; no size sorting or beta-pair semantics is used.

If $\alpha_s$ is retained, upper $\alpha_s$ and lower sole-owned
$C5_s$ give
$\alpha_s-C5_s\le S_s-78+\tau/23+E$.
But the literal identity is
\[
 \alpha_s-C5_s=\mathcal B+d_t+2(a+b)+G>d_t+2(a+b)+G.
\]
The strict sign belongs to branch B. With one leading-$t$ inequality,
$E+\ell_t\le\delta$, the mass identity and $K\le M+R+21$ give
\[
 S_t+M+2X_s+4X_t+g_s+2g_t<1041+R+\tau/23.
\]
For $Q_{62}$ with span at most $61+\tau/5$, two separately valid
stars and the four one-empty ordinary groups give
\[
 2X_s+4X_t+g_s+2g_t\ge558-9\tau/5-3\delta.
\]
The source mass gives
$S_t+M\ge T_r+4r+R+129+\mathcal A+\tau_t$, $r=r_t$.
Consequently
\[
 T_r+4r+I+\tau_t<354+97\tau/115+2\delta
                  \le448-2I,
\]
so $T_r+4r+3I<448$. For every integer $r$,
\[
 T_r+4r+3(r-34)(r-35)
 =641+\tfrac12\{7(r-29)^2+(r-29)\}\ge641,
\]
a contradiction. This proof needed only the four extreme leaves, not
retention of all eight lower resources or of any beta singleton.

If $\alpha_s$ is empty, it is the only extra hole outside the two
mandatory beta/H choices. Thus pair1, $\alpha_t$, and all eight lower
own/cross singletons are retained; pair1's actual owner is $s$.
The four resources $P_s,\alpha_t,C6_t,U_s$ therefore satisfy the
branch-independent endpoint lemma. It gives the unchanged contradiction
$15355\le13987$. In particular it does not borrow a beta-pair capacity
from branch A. Both alternatives are impossible. Since A and B include
respectively $\mathcal B\le0$ and $\mathcal B>0$, all of Z5 is excluded.


\section{Two missing units with positive core holes: $P_2$}\label{sec:p2}
We retain physical sign labels $s,t$; interchanging them never also sorts the
side sizes. Here $j=2$ and $r_s+r_t=69$. The original positive-hole cap is $\delta\le313-I-n(r_s)-n(r_t)$. A side of size at most 18 would give $I\ge272$ and a sum of defect floors at least 50, so $19\le r_s,r_t\le50$.
Write $C_v=\sum P_v$, $C_0=C_s+C_t$, $R_v=\max P_v$, $R=R_s+R_t$,
$M_v=\max D_v$, $M=M_s+M_t$, and $S=S_s+S_t$. Put
$\mathcal H=h(P_s)+h(P_t)$ for the sum of the two core support-hole counts.
The source envelope and layer inequalities give
\begin{align}
 I+\tau+\delta&\le281+w,& K-\delta&\ge33-w+I+\tau,\label{eq:p2-source}\\
 K-M&\ge R-269-w+I+\tau+\delta, &R&\le73+\rho+\mathcal H,\nonumber\\
 C_0-U_0&\ge364+2I+\delta-2z,& C_1-V_0&\ge492+\mu+\rho+2I+\delta+2z,\nonumber\\
 U_0+2\delta-6z&\le1146-2I.&&\nonumber
\end{align}
Here $U_0$ is the total of the first two peeled magnitudes on each side,
$V_0$ that of the last two, and $C_1=C_0+U_0$. Initially $w\le4z\le56$.
Individual primitivity gives $\mu\ge63$, $50\le\rho\le\mu$, and
$\tau=\mu+\mathcal H$. These statements concern the source alone.

\subsection{Localization and the first gap}
\begin{proposition}\label{prop:p2-local}
There is a common lower strip with hole sets $\Gamma_D\subseteq\Gamma_E$
such that
\[
 z=2|\Gamma_E|,\qquad w=2|\Gamma_D|,\qquad
 2\le w\le z\le14,
 \qquad [1,S_v]\setminus F_v=\{S_v-h:h\in\Gamma_D\}.
\]
Every integral center in $[1-M_v,S_v]$ is active. The first missing positive
absolute-core subset sum $m$ survives both peels and satisfies
$m<\min(R_s,R_t)$.
\end{proposition}
\begin{proof}
From \eqref{eq:p2-source},
$\mathcal H\le218+w-I-\delta\le274-I-\delta$.
For every integer $r$, writing $u=r-23$ gives
$T_r+I=408+u(3u+1)/2\ge408$; hence $C_v>\mathcal H$.
The common-strip counting lemma supplies
$[\mathcal H,C_0-\mathcal H]\subseteq B_E$, where $B_E$ is the
subset support of the absolute-core \emph{multiset}.

For clarity, the inner peel requires all sixteen choices of its four
occurrences. Sort these occurrences as $t_1\le\cdots\le t_4$ and set
$L=C_0-2\mathcal H$. The maximum difference between consecutive distinct
subset sums of positive weights in a chosen order is at most
\[
 \max_i\bigl(t_i-\sum_{h<i}t_h\bigr).
\]
Indeed adjoining $t_i$ unites a subset support and its translate: old gaps
persist or shorten, and the only possible new gap is between their hulls.
Here $L\ge1070+2I+\delta-w\ge1014$, whereas $U_0\le1230$.
The first three gap terms are at most $t_3\le U_0/2\le615$.
If $r_{\min}=34-k$, the three entries other than a largest one have total
at least $73+r_{\min}$. Therefore
\[
 L-\bigl(t_4-t_1-t_2-t_3\bigr)
 \ge2+4k^2+2k+3\delta>0.
\]
Thus all translates of $[\mathcal H,C_0-\mathcal H]$ by the actual inner
subset sums overlap, filling $[\mathcal H,C_1-\mathcal H]$.
For the outer peel,
\[
 C_1-V_0-2\mathcal H\ge169+4I+3\delta+2z-2w
 \ge85+4I+3\delta>0.
\]
Selecting none or all four outer peeled occurrences fills
$[\mathcal H,S-\mathcal H]$.

Complement symmetry at both layers, in the same strip
$\{1,\ldots,\mathcal H-1\}$, now gives
$z=2|\Gamma_E|$, $w=2|\Gamma_D|$ and
$\Gamma_D\subseteq\Gamma_E$. Translation of the absolute support by
$-S_o$ places exactly the reflected strip holes at the top of each target.
In particular $\eta_v:=S_v-|F_v|=w/2\le7$.

After replacing $w\le56$ by $w\le14$, the source rows give
$\mathcal H\le232-I-\delta$, $R\le368-I-\delta$. Moreover
\[
 S_v-\mathcal H-R_o
 \ge\frac58(2r_v-53)^2+\frac{275}{8}+2\delta
 \ge35+2\delta,\qquad T_{r_v}-R\ge40+\delta.
\]
If $1\le u\le S_v$ is a target hole, the $r_o\ge19$ values
$u-P_o$ are positive hull integers, and at most seven can be holes.
If $1-M_v\le u\le0$, then $u+M_v$ is in the complete lower target
interval, since $S_v-M_v\ge C_v>\mathcal H$. This proves activity.

The greedy first-gap rule says that entries below $m$ have total $m-1$
and generate $[0,m-1]$, and every later entry exceeds $m$.
If $R_v\le m$, the whole side lies in that prefix and
$C_v\le m-1<R$, contrary to $C_v\ge T_{r_v}>R$.
All peeled entries therefore exceed $m$, which persists as a hole in
$B_D$. It belongs to $\Gamma_D$, giving $w\ge2$.
\end{proof}

The outer centers have the common notation $O_v,H_v,U_v,V_v,C_{iv}$
($U_v=\mathrm{own2}_v$, $V_v=\mathrm{own3}_v$), with
\[
 d_v=v_1-v_2,\quad h_v=v_2-v_3,\quad g_v=v_3-v_4,\quad
 X_v=v_4-M_v,\quad X=X_s+X_t,\quad G=g_s+g_t.
\]
We use $\lambda_v=v_1-v_2-v_3-v_4$ for the leading shift.
Both leadings are active, since $K-\delta\ge19$. Thus $\lambda_v$ is
integral. A nonintegral lower coordinate $v_3$ or $v_4$ forces respectively
$\{U_v,C_{6o}\}$ or $\{V_v,C_{7o}\}$ to be missing singleton units.
Two such coordinates would require four holes. If exactly one were
nonintegral, these two holes would leave $H_v,\alpha_v$ active, contrary
to $H_v-\alpha_v=\lambda_v+v_3+v_4$. Hence all lower coordinates are
integral. A nonintegral $v_2$ then makes $O_v,H_v,\alpha_v$ three missing
singletons. All outer coordinates are consequently integral.

Define the cap predicate
\[
 \mathsf T_v:\quad O_v\text{ active, or both }\beta_v,H_v\text{ active}.
\]
Some $\mathsf T_v$ holds: if both $O$ units are missing, all $H$ units and
the relevant beta resources are retained. The raw-gap lemma gives
\begin{equation}
 \mathsf T_v\Longrightarrow d_v-M\le56+w/2\le63.
 \label{eq:p2-cap}
\end{equation}
For the first alternative use the gap from $H_v+M_v$ to $O_v-M_o$;
leading activity places its lower end above zero. For the second, with
$O_v$ empty, beta activity gives $h_v-d_v\ge X_v+1$, hence
$\alpha_v>U_v$; use the gap from $\alpha_v+M_v$ to $H_v-M_o$.
Only the 55 raw fixed labels and the $w/2$ holes of this one target pay
for either gap. No beta component capacity is used.

\subsection{Rank profiles and the first-gap reduction}
Let $k_v$ count entries below $m$, $k=k_s+k_t$, and $n_v=r_v-k_v$.
Their combined low mass is $m-1$. Distinctness on each sign implies
\begin{equation}
 I+\tau\ge70m-\frac54-(m+35)k+\frac{k^2}{4}
 +\left(r_s-\frac{69}{2}-\frac{k_s-k_t}{2}\right)^2.
 \label{eq:p2-lowmass}
\end{equation}
Since $2m<R\le368$, $m\le183$, and
$k(k+2)/4\le m-1\le182$ gives $k\le26$.
For $m\ge6$, the three ranges $k\le3$, $k=4$, $5\le k\le26$ give,
respectively,
\[
 I+\tau\ge\frac{1183}{4},\quad
 I+\tau\ge\frac{1283}{4},\quad
 I+\tau\ge\frac{-k^3+69k^2-4k+275}{4}\ge\frac{623}{2}.
\]
In the middle case $m\ge7$; in the last use
$69-k\ge43$, $k^2\ge25$, $-4k\ge-104$.
All contradict $I+\tau\le295$, so $m\le5$.

The following are coordinatewise minimal side profiles. Their excess
$\sigma$ is nonnegative; $B$ is the sum of their largest entries.
Opposite occurrences of the same magnitude remain distinct.
\begin{center}\small
\begin{tabular}{c|l|c|c}
$m$&low entries and side profiles&$I+\tau$&$B$\\\hline
1&both $(2,3,\ldots,r_v+1)$&$I+69+\sigma$&71\\
2&$(1,3,\ldots,k+1)$; $(3,4,\ldots,q+2)$&$102+(k-35)^2+\sigma$&72\\
3&both $(1,4,5,\ldots,r_v+2)$&$I+134+\sigma$&73\\
4, split&$(1,5,\ldots,r_s+3)$; $(2,5,\ldots,r_t+3)$&$202+I+\sigma$&75\\
4, together&$(1,2,5,\ldots,r+2)$; $(5,6,\ldots,q+4)$&$202+(r-35)(r-36)+\sigma$&75\\
5&$(1,2,6,\ldots,r+3)$; $(1,6,\ldots,q+4)$&$231+(r-35)^2+\sigma$&76
\end{tabular}
\end{center}
The low multisets are respectively empty, $\{1\}$, $\{1,1\}$,
$\{1,2\}$ and $\{1,1,2\}$: this follows from their total $m-1$, interval
completeness below $m$, and multiplicity at most two.

Delete the four largest core magnitudes on each side, leaving $Q$ of
size 62. With baseline maxima $B_v$, the maximum-deleted excess on side
$v$ is $\sigma_v-R_v+B_v\ge0$. The four ranks ending immediately before
the maximum therefore give
\[
 \operatorname{span}(Q)\le B-8+\frac{\sigma-R+B}{4},\qquad
 |Q\cap(Q+q)|\ge f-q,\qquad
 4f=524-5B-\sigma+R.                         \tag{P2-rank}
\]
This uses actual ranks, not realizability of the baseline support.

For a cap side $v$, choose one of the arms
$L_v-V_v-U_v$ and $L_o-C_{7o}-C_{6o}$ with at most one empty ordinary
vertex, and the first edge of the other arm. Such a choice exists since
$j=2$. The three edge gaps sum to $2X_v+g_v$.
The target-independent empty-vertex forest rule from Section~\ref{sec:separated-arms} charges these intersections to a group
of retained resources of deficiency $\Delta\le\delta$, while the cap-side
leading loss $\ell_v\le\Delta$. Thus
\begin{equation}
 d_v-M\ge K-M+3f-2\delta.                 \label{eq:p2-doubleforest}
\end{equation}
There is zero capacity on every empty vertex. For $m=4,5$, set $\chi$ to
the nonnegative square/product in the profile table. Substituting
\eqref{eq:p2-source} into \eqref{eq:p2-doubleforest} gives
\[
 \begin{array}{c|c|c}
 m&\delta\text{ cap}&d_v-M\text{ lower bound}\\\hline
 4&79+w-\chi-\sigma&97-2w+2\chi+5\sigma/4\ge69\\
 5&50+w-\chi-\sigma&153-2w+2\chi+5\sigma/4\ge125
 \end{array}
\]
Both contradict \eqref{eq:p2-cap}. Hence $m\in\{1,2,3\}$.

\subsection{Rooted forests and exhaustive activity reduction}
First, $\mathsf T_s,\mathsf T_t$ cannot both hold. The raw-gap caps would
imply $d_s+d_t\le2M+112+w$. The leading counts give
$d_s+d_t\ge2K+G+2X-\delta$.
The high-tail count $|F_v|\le d_v+g_v+341$ and
$S=2K+1332+w-\delta$ give
$d_s+d_t+G+\delta-2K\ge650$.
The four-edge star, allowing empty extreme singleton leaves with zero
capacity, gives $2X+4R+\delta\ge556$. At each target its positive and negative restricted edges occupy disjoint zones since $M>R$, so no point sees two empty leaf vertices. Setting
$Y=2M-2K+112+w+\delta$, these inequalities imply
\[
 2Y\ge1206-4R-\delta,
 \quad Y\le650+3w-2R-2I-2\tau-\delta,
 \quad4I+4\tau+\delta\le94+6w.
\]
But $\tau\ge63+2z$ and $w\le z$ would give $158+2z\le0$.
In particular at least one $O$ singleton is empty.

No $C_{7v}$ is empty. If it were, it and one $O$ would be the two holes,
so $C_{6v}$ is active. The empty-center dichotomy forces
$C_{7v}\ge S_v+1$, hence $\lambda_v>S_v+M_o$.
The leading contains $\lambda_v+(D_v\cup(M_v+D_v))$, at least $r_v+82$
points above the target; at least three of the higher singletons
$V_v,U_v,H_v,O_v$ are retained and each loses $r_v+5$ points.
Thus $\delta\ge4r_v+97$, and the source bound gives
$I+4r_v\le115+w\le129$, whereas
$I+4r_v=(r_v-32)(r_v-33)+134\ge134$.

If $\mathsf T_v$ and $V_v$ is retained, at least one of $U_v,C_{6o}$
is retained (one of the two holes is already an $O$). Choose the complete
arm and the other first edge. All vertices are now retained. Designate
$L_v$ as a root in the common forest lemma: its native loss is absorbed
in the \emph{same} budget. The three overlaps and the leading bound yield
\begin{equation}
 d_v-M\ge K-M+3f-\delta
 \ge R-269-w+I+\tau+3f.                   \label{eq:p2-root}
\end{equation}
The three profile rows give respectively
\[
 d_v-M\ge51-w+I+\sigma/4\ge37,\quad
 d_v-M\ge82-w+(k-35)^2+\sigma/4\ge68,\quad
 d_v-M\ge112-w+I+\sigma/4\ge98.
\]
Thus $m=2,3$ are impossible in this retained-$V_v$ case.

If $\beta_v$ and $O_v$ are both active, a stronger cap is available.
The two gaps from $\alpha_v+M_v$ to $H_v-M_o$ and from $H_v+M_v$ to
$O_v-M_o$ lie in the target hull, are disjoint, and contain no nonfixed
target points. Their union has $2(d_v-M-1)$ integers when nonempty.
Counting its target holes only once gives
\begin{equation}
 d_v-M\le(57+w/2)/2\le32.                 \label{eq:p2-two-gap}
\end{equation}
Consequently \eqref{eq:p2-root} also excludes $m=1$ with simultaneous
$\beta_v,O_v$ activity.

It remains to justify every missing-$V_v$ endpoint case. If $\mathsf T_v$
and $V_v$ is empty, the other hole must be $O_o$. If instead it were
$O_v$, then $O_o,H_v$ are active, so $\mathsf T_v$ would require
$\beta_v$, giving both cap predicates. Both $H$ units are retained.
Branch B would retain both beta singletons and again give two cap
predicates. In A, the retained beta pair is sole on $v$, because an
active $\beta_o$ would give $\mathsf T_o$. Therefore $\beta_v,O_v$ are
active and \eqref{eq:p2-two-gap} holds.
The empty center $V_v>-M_v$ is high, so the distinct retained $H_v,O_v$
each have native loss at least $r_v+5$. Use edges $L_v-U_v$ and
$L_o-C_{7o}$, root $L_v$, and those two unused native losses. This gives
\[
 d_v-M\ge R-269-w+I+\tau+2f+2(r_v+5).
\]
For $m=1$ this is at least $69-w\ge55$, using
$I+2r_v=(r_v-33)(r_v-34)+68$. For $m=2$ it is at least $86-w\ge72$,
using $(k-35)^2-I=35-k\ge-15$; for $m=3$ even discarding the two losses
gives $54-w\ge40$. All contradict 32.
We have proved, without an endpoint restriction,
\begin{equation}
 m=1,\qquad \text{no side has both }\beta_v\text{ and }O_v\text{ active}.
 \label{eq:p2-reduced}
\end{equation}

Here is a direct exhaustive count of the surviving units. If both $O$
are missing, branch B retains both $H$ and beta singletons, contrary to
the two-cap result. In A the beta pair is retained and actually sole,
because both $H$ are active. This is the form with holes $\{O_s,O_t\}$.
If exactly $O_v$ is missing and $W$ is the other hole, then $O_o$ is
active and \eqref{eq:p2-reduced} makes $\beta_o$ inactive.
In B this forces $W=\beta_o$, leaving $\beta_v,H_v$ active and giving two
cap sides, impossible. In A, if $W=H_v$, the retained beta pair is sole
on $v$. Otherwise $H_v$ is active, so the two-cap result makes
$\beta_v$ inactive too; thus $W$ is the whole beta pair. We obtain exactly
\[
 \{O_v,\mathrm{beta\ pair}\},\qquad
 \{O_s,O_t\},\qquad \{O_v,H_v\}.
                                                        \tag{P2-forms}
\]
All unlisted whole units are retained. This argument does not need a
choice of contributor for pair1, whose capacity is never separated below.

\subsection{The trim-six contradiction}
At $m=1$ write
$\sigma_v=\sum_i(x_{v,i}-i-1)$ and
$\xi_v=R_v-r_v-1\ge0$, with $\sigma=\sigma_s+\sigma_t$,
$\xi=\xi_s+\xi_t$. Thus
\begin{equation}
 I+\sigma+\delta\le212+w.                         \label{eq:p2-finalmass}
\end{equation}
The symbol $\xi$ is a radial surplus and is unrelated to the structural
mixed-sum excess $\kappa$.

In the first and third forms of (P2-forms), choose the side opposite the
missing $O_v$; it has $O,H$ active. In the middle form choose the actual
sole beta owner; it has $\beta,H$ active. On the chosen cap side $u$,
both $\alpha_o$ and $C_{5u}$ are retained singletons. Their activity and
the literal identity
\[
 \alpha_o+C_{5u}=d_u-2(v_4+o_4)-G
\]
imply $d_u\ge M+2X+G+2$. Combining with
\eqref{eq:p2-cap} gives
\begin{equation}
 2X+G\le54+w/2.                                  \label{eq:p2-finalcap}
\end{equation}

Delete six largest core magnitudes on each side, leaving $Q$ of size 58.
After excluding the maximum, the six ranks ending at $r_v-1$ dominate
six times the deviation at $r_v-6$, whence
\[
 \operatorname{span}(Q)\le59+(\sigma-\xi)/6,\qquad
 |Q\cap(Q+q)|\ge56-(\sigma-\xi)/6-q=:f-q.
\]
In every form all six singleton units $U_v,V_v,C_{7v}$ are retained.
Use, on each target, the tree $L_v-V_v-U_v$ with the additional edge
$L_v-C_{7v}$. The three gaps are $X_v,g_v,X_o$; the full leading is
counted once even though two of its restricted core copies certify edges.
The two trees give
\[
 \delta\ge6f-2X-G,
 \qquad 2X+G\ge336-\sigma+\xi-\delta.
\]
Combining this with \eqref{eq:p2-finalcap} and
\eqref{eq:p2-finalmass} forces
\[
 70+I+\xi\le\frac32w\le21,
\]
a contradiction. This excludes all $P_2$ configurations, including every
position of every endpoint; no normal-center premise was used.


\section{Three missing units with positive core holes: $P_3$}\label{sec:p3}
Here $j=3$ and $\delta=c-234$. We begin with the original positive-hole
range $1\le z\le14$ and $15\le r_s,r_t\le54$, not with a strengthened
hole profile. All signs in this section are physical signs. Put
$\mathcal H=h(P_s)+h(P_t)=\kappa+z$, and retain the source-offset notation
$U,C,J,L$ of Lemma~\ref{lem:common-envelope}; write $C_v=\sum P_v$, $C_0=C_s+C_t$. The symbols $U_0,V_0$ remain layer
\emph{masses}. The absolute values of a signed core or source are always
a labelled multiset: equal magnitudes on opposite signs are independently
selectable. Denote its subset support by $B_E$ or $B_D$.

\subsection{A profile-free baseline and the actual boundary holes}
\begin{proposition}\label{prop:p3-baseline}
Every original $P_3$ configuration satisfies
\[
22\le r_s,r_t\le47,\quad z\in\{2,4,6,8\},\quad
0\le w\le z,\quad w\equiv0\pmod2,\quad
\delta\le140-z-I\le138-I.
\]
\end{proposition}
\begin{proof}
The structural and layer rows specialize to
\begin{align}
\delta&\le207-I-\rho+2z,&
\rho+\kappa+\delta&\le207+3z-I,\label{eq:p3-initial}\\
U_0-C_0+\delta-2z&\le-442-2I,&
2V_0-S+\delta+2z&\le-570-2I-\mu-\rho.\nonumber
\end{align}
For $r=\min(r_s,r_t)\le20$, the general defect floors give
$\rho\ge50$ and $I\ge210$, hence $\delta\le235-210-50<0$.
For $r=21$, they give $\rho\ge59$, $I=182$ and
$\delta\le235-182-59<0$. Thus $r\ge22$, and the generic floors
$\rho\ge59$, $\mu\ge63$ apply. In particular
\[
\mathcal H\le148+4z-I-\delta\le204<T_{22}.
\]
The common-strip counting argument gives
$[\mathcal H,C_0-\mathcal H]\subseteq B_E$. The two layer rows give
\[
C_0-U_0-2\mathcal H\ge146+4I+3\delta-10z\ge6,
\quad
C_1-V_0-2\mathcal H\ge396+4I+3\delta-6z\ge312.
\]
Selecting none or all four occurrences of each peel therefore fills the
central intervals of both subsequent layers. Complement symmetry and
inclusion inside the \emph{same} strip $\{1,\ldots,\mathcal H-1\}$ give
\[
z=2|\Gamma_E|,\qquad w=2|\Gamma_D|,\qquad
\Gamma_D\subseteq\Gamma_E.
\]
This proves evenness and $0\le w\le z\le14$ without presuming persistence.

The source-only envelope, with $\mathcal A=I+\tau+\delta$, is
\begin{equation}\label{eq:p3-source}
\begin{gathered}
S=2K+1254+w-\delta,\qquad
2K=-29-w+\mathcal A+4R+U+2C+J,\\
J+U\le365+w-\mathcal A,\qquad
J\ge\mathcal A-47-w,\qquad L\ge\mathcal A-37-w,\\
\mathcal A\le203+w\le217,\qquad
K-M\ge R-191-w+\mathcal A,\qquad K\le M+R+21.
\end{gathered}
\end{equation}
The last bound counts the central interval of $M+R+1$ positions and at
most ten same-sign peeled-pair labels on each sign. The other rows are
Lemma~\ref{lem:common-envelope} and its two fringe comparisons.
Consequently $C_0\le1442$, $\tau\le217$, and
$R_v\le r_v+\tau_v\le264$.

At most four integers of $[1,17]$ are absent from both cores. Indeed, if
five missing values have sum at most $13+14+15+16+17=75$, avoiding them
forces on each side $\tau_v\ge5r_v+15-75$, whence $\tau\ge225$.
Choose thirteen distinct values present in $[1,17]$, using one actual
labelled occurrence for each. Their set $B$ has total $T\ge91$.
Lev's central-interval theorem, with $L=17$ and $|B|=13$, gives
$[9,T-9]\subseteq\FS(B)$, since $17\le3(13)/2-2$.
Adjoin every remaining labelled core occurrence in nondecreasing order.
If $n\ge33$ occurrences have been selected, their total is at least
$T_{16}+T_{17}=289$, so the next $x\le264\le T-17$.
If $n\le32$ and $x>T-17$, then $x\ge75$ and the remaining at least
37 occurrences have total at least $37\cdot75$; thus
$C_0\ge91+37\cdot75>1442$, impossible.
Hence always $x\le T-17$, and the intervals
$[9,T-9]$ and $[x+9,T+x-9]$ join. Induction proves
$[9,C_0-9]\subseteq B_E$.
The seed omits only four of the seventeen singleton positions, so there
are at most four low holes in $[1,8]$. Complement symmetry gives $z\le8$.
Finally $\tau=\mu+\kappa+z\ge63+2z$, so
\[
\delta\le203+w-I-\tau\le140+w-2z-I\le140-z-I.
\]
No first-gap conclusion or special direct-profile theorem was used.
\end{proof}

\begin{proposition}\label{prop:p3-boundary}
The first positive gap $m$ of $B_E$ persists in $B_D$ and satisfies
$m<\min(R_s,R_t)$. Moreover
\[
2\le w\le z\le8,\qquad \eta_s=\eta_t=w/2,\qquad
[1,S_v]\setminus F_v=\{S_v-h:h\in\Gamma_D\}.
\]
Every integral center in $[1-M_v,S_v]$ is genuinely active. An integral
empty center above a known active lower center is at least $S_v+1$.
\end{proposition}
\begin{proof}
Now $\mathcal H\le180-I-\delta$; the preceding complete strips remain
valid. Translating $B_D$ by $-S_o$ places its low holes strictly below zero
and its high holes precisely at $S_v-\Gamma_D$, because $S_o>\mathcal H$.
This proves the target-hole formula even if $w=0$.
The maximum-removal row gives $R+\delta\le280+4z-I\le312-I$.
For $r=r_v$, $S_v\ge T_r+4R_v+10$ and $R_v\ge r$, so
\[
S_v-\mathcal H-R_o
\ge\frac58(2r-53)^2+\frac{1139}{8}+2\delta
\ge143+2\delta>0.
\]
A target-hole center $u\in[1,S_v]$ has $r_o\ge22$ distinct positive
hull images $u-P_o$, whereas only $w/2\le4$ hull integers are holes.
An ordinary nonhole center uses the source point zero. If
$1-M_v\le u\le0$, the image $u+M_v$ is in the complete lower target
interval, since $S_v-M_v\ge C_v>\mathcal H$. This proves activity on
the claimed interval and the empty-center dichotomy. An active lower
center rules out its low arm.

For every integer $r$, writing $t=r-23$ yields
$T_r+I=408+t(3t+1)/2\ge408$. Hence $T_{r_v}-R\ge96+\delta$.
At the first failed greedy inequality for the sorted absolute-core
multiset, the preceding entries have total $m-1$, generate $[0,m-1]$,
and every remaining entry exceeds $m$. If $R_v\le m$, that entire
side would be in the prefix, giving $C_v\le m-1<R<T_{r_v}$.
Thus both core maxima exceed $m$; all peeled magnitudes exceed $m$ too,
so the gap persists. Central completeness puts it in $\Gamma_D$.
Therefore $w=2|\Gamma_D|\ge2$.
\end{proof}

\subsection{Activity, outer integrality, and the two exact forms}
Write $U_v=\mathrm{own2}_v$, $V_v=\mathrm{own3}_v$ and
$C_{iv}=\lambda_v-o_{i-3}$ for the common cross-center notation ($4\le i\le7$). The leading resource
is $\mathcal L_v=\lambda_v+(D_v+D_v)$. Put
$X_v=v_4-M_v$, $X=X_s+X_t$ and $G=g_s+g_t$.

Both leadings are active because $K\ge155>\delta$, so $\lambda_v$ is
integral. Each lower coordinate $v_3$ has the two singleton witnesses
$U_v,C_{6o}$, and $v_4$ has $V_v,C_{7o}$; all eight witnesses are distinct.
Two nonintegral lower coordinates would require four holes. If a sole
nonintegral lower coordinate $x$ remains, its two witnesses and one of
$H_v,\alpha_v$ are missing, since their difference is
$\lambda_v+v_3+v_4$. If $\alpha_v$ is missing, retained $O_v,H_v$ force
both top coordinates integral, contradicting the equation for $\lambda_v$.
If $H_v$ is missing, retained $O_v,\alpha_v$ and, in A, $C_{5o}$,
or in B, $\beta_v$, force $v_2$ integral and again force $x$ integral.
Thus the lower coordinates are integral. A nonintegral top pair would
make $O_v,H_v,\alpha_v$ and also $C_{5o}$ in A or $\beta_v$ in B missing.
All eight outer magnitudes are therefore integral.

Here are the branch-independent forest rows used in the classifier:
\begin{equation}\label{eq:p3-classifier-forests}
2X_v+X_o+g_v\ge256-\tau-\delta,\quad
2X_v+X_o+G\ge310-\tau-\delta,\quad
2X+4R+\delta\ge556.
\end{equation}
For completeness, deleting the largest $k$ core points on each side gives
$|Q_k|=70-2k$, span at most $69-2k+\tau/(k+1)$, and intersections at
integral gap $q$ of at least $70-2k-\tau/(k+1)-q$. For $k=0$ the full core has its actual span $R$, so the sharper intersection floor is $139-R-q$. Four full-core stars therefore give $556-4R-2X$, which proves the third forest row.
The two candidate arms for sign $v$ are
$\mathcal L_v-V_v-U_v$ and $\mathcal L_o-C_{7o}-C_{6o}$,
with gaps $X_v,g_v$. Four distinct ordinary vertices are involved, so
one arm has at most one empty vertex. For the first row choose it,
the unused $X_v$ edge, and either $X_o$ edge, using $k=3$.
For the second choose an arm for each sign and the unused $X_v$ edge,
using $k=4$. For the third use four full-core stars ($k=0$).
On one target the positive restricted supports lie above
$\lambda_v+M_v-R_o$ and negative ones below $\lambda_v-M_o+R_v$.
They are disjoint since $M>R$. At any point the restricted forest
therefore lies in one arm, with at most one empty vertex. Inside the
target it has no empty vertex, so its edge count is at most retained
multiplicity minus one; outside, it is at most retained multiplicity.
Summation charges it to the retained-resource union deficiency. This
proves all three rows with zero capacity on empty vertices and each
leading capacity counted once.

The $O_v$ raw cap needs no $H_v$ activity here: both leadings are active, so $H_v+M_v>\lambda_v+2M_v\ge1$, while $O_v$ activity gives $O_v-M_o\le S_v$. The first raw gap is therefore inside the actual target hull and has no nonfixed cover. This is the leading-activity alternative in Lemma~\ref{lem:common-raw-gaps}.

The following five activity assertions suffice:
\begin{enumerate}
\item $O_v,C_{6v}$ cannot both be active
\item $O_s,O_t$ cannot both be active
\item $\beta_v,H_v,C_{6v}$ cannot all be active
\item $C_{5v}$ active implies $C_{6v}$ active
\item $\beta_v,O_v,C_{7v}$ cannot all be active
\end{enumerate}
For assertions 1 and 3, the single raw-gap bound is
$d_v-M\le56+w/2$. When $O_v$ is empty the beta/$H$ version uses
$\alpha_v$ only as a numerical ceiling, not as a capacity-bearing vertex.
Actual $C_{6v}$ activity and its center identity give
$d_v-M\ge2X_v+X_o+G+1\ge311-\tau-\delta$.
Thus $52+I\le3w/2\le12$, impossible.
For assertion 5 use the double gap
$d_v-M\le57/2+w/4$ and
$d_v-M\ge2X_v+X_o+g_v+1\ge257-\tau-\delta$;
together they imply $102+4I\le5w\le40$.

For assertion 2, both raw-gap caps would give
$d_s+d_t\le2M+112+w$. The leading maxima and the high-tail count yield
\[
d_s+d_t\ge2K+G+2X-\delta,\qquad
 d_s+d_t+G+\delta-2K\ge572.
\]
For the tail count, subset sums above $d_v+g_v$, after adding the four
outer magnitudes, require at least three top occurrences in their
four-term representation; there are at most $4(78+6)+5=341$ values.
Set $Y=2M-2K+112+w+\delta$. The preceding rows and the four-star row give
$2Y\ge1128-4R-\delta$, whereas \eqref{eq:p3-source} gives
$Y\le494+3w-2R-2I-2\tau-\delta$. Hence
$140+4I+4\tau+\delta\le6w\le48$, a contradiction.
Finally, if $C_{5v}$ is active but $C_{6v}$ empty, the latter is HIGH.
Then $\lambda_v>S_v+M_o$, so the leading loses all 78 points of
$\lambda_v+D_v$. At least three of the five higher singletons
$C_{7v},V_v,U_v,H_v,O_v$ are retained, each losing at least $r_v+5$.
The original budget gives $\delta\ge78+3(r_v+5)\ge159>138$.
This proves assertion 4 without charging a $C_5$ component individually.

The owner grammar also follows directly from centers. For
$W=v_3+v_4+o_3+o_4>2M$,
\[
\begin{gathered}
C_{4s}+C_{4t}=-(s_2+s_3+s_4+t_2+t_3+t_4),\\
\beta_s+\beta_t+C_{5s}+C_{5t}=-2W,\\
C_{4v}+\beta_v+C_{5o}=-2v_3-2v_4-o_2-o_3-o_4,\\
C_{4v}+C_{5o}=-(\beta_s+\beta_t)-2W-d_o.
\end{gathered}
\]
Comparison with the actual lower activity endpoints $1-M_v$ proves:
retained pair1 has a sole actual owner; the displayed triple is
impossible; in B the two retained pairs have sole, matching owners.

We now exhaust the whole-unit holes. At least one $O$ is missing. If
both are missing, the third hole cannot be $C_{6v}$ in A, because
assertion 4 would also delete its singleton $C_{5v}$. Thus both $C_6$
are retained. An empty beta pair gives A0 below; otherwise a genuinely
active beta component $\beta_v$ forces $H_v$ missing by assertion 3.
The opposite beta is then inactive, and the owner grammar forces
pair1 owner $o$. In B, with both $O$ missing, assertion 3 needs a hole
in each disjoint triple $\{\beta_v,H_v,C_{6v}\}$, exceeding three holes.

If only $O_o$ is missing, $O_v$ is active, so $C_{6v}$ is missing and
$C_{5v}$ inactive. In A these are exactly the three singleton holes;
$H_o,C_{6o}$ force beta's sole owner to be $v$, and pair1's to be $o$.
Then $\beta_v,O_v,C_{7v}$ contradict assertion 5.
In B, write the empty $O$ as $O_v$: the holes are
$\{O_v,C_{6o},Z\}$, where assertion 3 forces
$Z\in\{\beta_v,H_v,C_{6v}\}$. Both pairs are retained with sole owner
$v$; assertion 4 excludes $Z=C_{6v}$. In either remaining case
$\beta_o,O_o,C_{7o}$ are active, again contradicting assertion 5.
Consequently exactly two forms remain:
\begin{equation}\label{eq:p3-forms}
\begin{array}{c|l|l}
\mathrm{A0}&\{\text{beta pair},O_s,O_t\}&\text{pair1 sole owner }s\text{ or }t\\
\mathrm{A1}&\{O_s,O_t,H_v\}&\text{beta sole owner }v,\ \text{pair1 sole owner }o
\end{array}
\end{equation}
Every unlisted whole unit is retained. No split or selected component
has acquired an individual capacity in this classification.

\subsection{First-gap profiles and uniform endpoint tools}
\begin{proposition}\label{prop:p3-profiles}
The first gap is one of $1,2,3,4$. The following side profiles are
coordinatewise lower bounds; $\sigma$ is excess above them and $B$ is
the sum of their maxima. Here $q=69-k$, and $k$ is the actual owner size of the unique
low 1 when $m=2$, or the side containing both low entries when $m=4$.
\begin{center}\small
\begin{tabular}{c|l|c|c}
$m$&side profiles&$I+\tau$&$B$\\\hline
1&both $(2,3,\ldots,r_v+1)$&$I+69+\sigma$&71\\
2&$(1,3,\ldots,k+1)$; $(3,4,\ldots,q+2)$&$102+(k-35)^2+\sigma$&72\\
3&both $(1,4,\ldots,r_v+2)$&$I+134+\sigma$&73\\
4, split&$(1,5,\ldots,r_s+3)$; $(2,5,\ldots,r_t+3)$&$202+I+\sigma$&75\\
4, together&$(1,2,5,\ldots,k+2)$; $(5,6,\ldots,q+4)$&$202+(k-35)(k-36)+\sigma$&75
\end{tabular}
\end{center}
In both $m=4$ placements, $z=w=2$ and $\delta\le3$.
\end{proposition}
\begin{proof}
Let $k_v$ count entries below $m$, with total low mass $m-1$.
Greedy completeness, sidewise distinctness, and $m<\min R_v$ give
\[
\tau_v\ge(m-k_v)(r_v-k_v),\qquad T_{k_v}\le m-1,
\qquad m\le155,\quad k_v\le17.
\]
For $m=5$, $k_v\le2$, $k_s+k_t\le3$, and one $k_v\le1$; thus
$\tau\ge3(69-3)+(22-1)=219$. For $m=6$ the same counts give
$\tau\ge4(69-3)=264$. For $m\ge7$, if $k_v\le3$ the side floor is
$4(r_v-3)$; otherwise $m-k_v\ge k_v+3$ and
\[
(k_v+3)(r_v-k_v)-4(r_v-3)
=(k_v-1)(r_v-k_v)-4k_v+12\ge k_v+7>0.
\]
Hence $\tau\ge252$. All contradict $\tau\le203+w\le211$.
For $m\le4$ the low multisets must be respectively empty, $\{1\}$,
$\{1,1\}$, $\{1,2\}$, which gives the table by summation.

At $m=4$, the table and $\mathcal A\le211$ give $\sigma+\delta\le9$.
A displaced tail rank propagates its positive integral deviation to all
later ranks, so every tail rank $i\le r_v-9$ equals its baseline.
Each side therefore contains $5,\ldots,12$. The absolute multiset
contains the seed $\{1,2,5,6,7,8,9,10,11,12\}$, of total 71.
The smaller seed $\{1,2,5,6,7\}$ has support
$[0,3]\cup[5,16]\cup[18,21]$. A support
$[0,3]\cup[5,T-5]\cup[T-3,T]$, on adjoining any $5\le x\le T-9$,
has the same form with total $T+x$: its two central intervals join,
and 4 and its reflection remain missing. Adjoin $8,\ldots,12$, then
all remaining labelled entries. Each remaining entry is between 5 and
$47+4+9=60\le71-9$, so the induction applies. Exactly the two core
holes $4,C_0-4$ remain. Thus $z=w=2$ and the mass cap improves to
$\delta\le3$.
\end{proof}

For all four profiles, deleting the source maxima gives the useful
penultimate estimate
$K\le R+C+p_s+p_t+176$: the central interval ends at
$R_v+c_v+p_v$, at most ten pairs per sign exceed it, and reinsertion
costs 155 labels. If $b$ is the profile's baseline $\tau$, subtracting
all but the final two core ranks gives
$p_s+p_t\le2B-2+\sigma-R$, hence $K\le C+2B+174+\sigma$.
Comparison with \eqref{eq:p3-source} yields
\begin{equation}\label{eq:p3-radial}
2R+I+\delta\le2B+209-b+w.
\end{equation}
The right sides for $m=1,2,3$ are respectively
$282+w,216+k+w,221+w$.
Delete the four largest core points on each side to form $Q$, of size 62.
The maximum-deleted profile excess is $\sigma_v-R_v+B_v\ge0$.
Four consecutive tail ranks ending before the maximum therefore give
\begin{equation}\label{eq:p3-rank}
\operatorname{span}Q\le B-8+\frac{\sigma-R+B}{4},\qquad
|Q\cap(Q+q)|\ge f-q,\qquad 4f=524-5B-\sigma+R.
\end{equation}
This rank argument asserts no interval completeness of $Q$.

\begin{lemma}[Compatible fixed credit]\label{lem:p3-fixed}
Both forms in \eqref{eq:p3-forms} have fixed credit $e=8$; A1 has $e=11$.
Every compatible grouping of retained whole resources has total
actual-target deficiency at most $\epsilon=\delta-e$.
\end{lemma}
\begin{proof}
Retained $C_{7v}$ below empty $O_v$ forces $O_v\ge S_v+1$.
The four fixed labels $C^v_{11},C^v_{12},C^v_{13},C^v_{14}$ are thus
HIGH on each sign. These eight labels are disjoint from the baseline
nonpositive-label witnesses of the original fixed-slot cover.
In A1, active beta lies below empty $H_v$, so $H_v\ge S_v+1$;
$C^v_{22},C^v_{23},C^v_{24}$ supply three further HIGH labels.
All are distinct labels, whether or not their values coincide.
Lemma~\ref{lem:common-fixed-credit}, or directly the 26 nonpositive
witnesses plus these HIGH labels, reduces the fixed allowance from 84
to $84-e$. Lemma~\ref{lem:common-loss} then gives the claimed one-budget
inequality, with each sole owner receiving its whole pair's capacity once.
\end{proof}
This immediately excludes $m=4$ in both forms: $\delta\ge8>3$.

For a retained singleton or genuinely sole pair owner $Z$ on target $v$,
write its native deficiency as
$\gamma_Z=78-|(Z+D_v)\cap F_v|=\ell_Z+h_Z+\nu_Z$,
where the three terms count low clipping, high clipping, and images at
interior target holes. There is no extra hole allowance.
\begin{lemma}[Uniform third-gap endpoints]\label{lem:p3-END2}
For $m=3$, $e\ge4$, and $\gamma_Z\le\epsilon=\delta-e$, one has
\[
Z\ge72-r_v-\ell_Z,\qquad Z\le S_v+r_o-71+h_Z.
\]
These rows include all centers outside the target hull.
\end{lemma}
\begin{proof}
The profile gives $\sigma+\epsilon\le77-e-I<74$.
For $Z\ge1$, distinct nonclipped negative images give the stronger row
$Z\ge74-r_v-\ell_Z$; for $Z\le S_v$, the stronger upper row is
$Z\le S_v+r_o-73+h_Z$.
If $Z\le0$, write $\ell_Z=74-r_v+j$. Since $\ell_Z<74$, the next
positive magnitude is a core rank $j+1$. If $j=0$ it is 1. Otherwise
its baseline is $j+3$; any displacement above that value would force
$\sigma_v\ge r_v-j=74-\ell_Z$, contradicting
$\sigma+\epsilon<74$. Thus $Z\ge1-(j+3)$, as required.
For $Z>S_v$ use the next opposite core magnitude with
$h_Z=74-r_o+j$; the identical rank argument proves the upper row.
\end{proof}

\subsection{Form A0: the beta pair is empty}
Let $s$ be the actual sole pair1 owner, put $P=C_{4s}$, and leave the
sizes unordered. The holes are exactly $\{\text{beta pair},O_s,O_t\}$.
Throughout take $e=8$ and $\epsilon=\delta-e$. Write
$\mathcal V=s_4+t_4=M+X$.

\paragraph{The first two gaps: one endpoint and forest argument.}
At $m=1$, $\tau=69+\sigma$, $R\ge71$, and
\[
\sigma+\delta+I\le134+w,\qquad L\ge32-w+I+\sigma+\delta.
\]
At $m=2$, let $k$ be the size of the \emph{physical} owner of the unique
1, independently of the pair1 owner $s$, and put $\chi=(k-35)^2$.
Then $R\ge72$ and
\[
\sigma+\delta+\chi\le101+w,\qquad
L\ge65-w+\chi+\sigma+\delta.
\]
By \eqref{eq:p3-rank}, the intersection floor $f$ is
$(169-\sigma+R)/4$ at $m=1$ and $(164-\sigma+R)/4$ at $m=2$.

Use the five distinct endpoint resources
\[
P,\quad\alpha_s,\quad\alpha_t,\quad C_{5t},\quad H_s.
\]
Their chosen clipping counts have respective weights
$1,2,1,1,1$, the first four lower and the last upper.
Let $N$ be the sum of their native deficiencies and $E_c$ this weighted
clipping sum; thus $E_c\le2N$. The weight two is an algebraic
coefficient, not a second copy of a resource.

Here is an explicit allowance for every endpoint position. Put $b=1$
for $m=1$ and $b=2$ for $m=2$. Normal endpoints satisfy
\[
Z\ge74-r_v-\ell_Z\quad(Z\ge1),\qquad
Z\le S_v+r_o-73+h_Z\quad(Z\le S_v).
\]
If a chosen endpoint is outside its hull, write its clipping count as
$c=74-r_v+j$ for a lower endpoint, or $c=74-r_o+j$ for an upper one.
Then $c\ge27$. The next magnitude supplies an error, relative to these
normal rows, at most $b+d$, where $d$ is its displacement above the true
profile baseline. If $c<74$, that magnitude is a core rank and
$(74-c)d\le\sigma_v$. In these two profiles
$\sigma+\epsilon\le134$ and $\sigma+\epsilon\le101$, respectively.
Thus $\sigma_v\le134-c$, and for $27\le c\le73$,
\[
c(75-c)-146=(c-2)(73-c)\ge0,
\qquad134-c\le c(74-c).
\]
The error is consequently at most $b+c$.
If $c\ge74$, retention gives $c\le77$, the next magnitude is peeled,
and its error is at most $b+\sigma_v+e_v$; here $e_v$ is the peeled
source offset on the physical side supplying that next magnitude.
For upper endpoints this is the opposite source side.
Because $N\le\epsilon\le134<148$, at most one selected clipping
count is at least 74. Call it heavy, with weight $q\in\{1,2\}$ if it
exists. Let $B_e$ be the weighted endpoint error sum. The exhaustive
bounds are
\begin{equation}\label{eq:p3-a0-errors}
\begin{array}{ll}
\text{no heavy count:}&B_e\le6b+2\epsilon,\\
\text{one heavy count }c:&B_e\le6b+2\epsilon-qc+q\sigma+qe_v,\quad
\delta\ge e+c.
\end{array}
\end{equation}
Normal endpoints have zero error and are included. This also covers a
next magnitude in a peeled tail, with no assumption about its owner.

The literal identity
\[
H_s=P+2\alpha_s+\alpha_t+C_{5t}+4\mathcal V+2G
\]
comes from $H_s=P+\alpha_s+\alpha_t+d_t+2\mathcal V+G$ and
$\alpha_s+C_{5t}=d_t-2\mathcal V-G$.
Apply the endpoint rows just proved. The lower constants total
$232-r_s$ and the upper constant is $S_s-r_s-4$, yielding
\[
S_s\ge236+4\mathcal V+2G-E_c-B_e.
\]
On each target take the full retained tree
$\mathcal L_v-V_v-U_v$ together with $\mathcal L_v-C_{7v}$.
The two trees have six edges, gaps totaling $2X+G$, and all their
vertices are distinct from the five endpoint resources.
If $D_F$ is their grouped deficiency, the forest lemma gives
$D_F\ge6f-(2X+G)$, and the simultaneous original grouping gives
$N+D_F\le\epsilon$. Thus $E_c+2D_F\le2\epsilon$.
Using the exact source mass and $K\le M+R+21$ proves
\begin{equation}\label{eq:p3-a0-main}
S_t+2M+12f\le2R+1060+w+\delta-2e+B_e.
\end{equation}
This single endpoint--forest budget already includes target holes.

For the final integer comparisons define $I_r=(r-34)(r-35)$. Direct
identities valid for every integer $r$ are
\begin{equation}\label{eq:p3-a0-quadratics}
\begin{aligned}
T_r+5r+3I_r&=670+\tfrac12(r-29)(7(r-29)+3)\ge670,\\
T_r+5r+4I_r&=695+\tfrac12(r-30)(9(r-30)-1)\ge695,\\
T_r+5r+5I_r&=711+\tfrac12(r-31)(11(r-31)+3)\ge711.
\end{aligned}
\end{equation}
Each product is nonnegative at every integer argument.

\paragraph{A0 at $m=1$.}
Put $r=r_t$. The three nonmaximal peeled values are at least
$R_t+1,R_t+2,R_t+3$, and hence
\[
S_t+2M\ge T_r+5r+\sigma_t+2R+2L+14.
\]
Substitute this, $12f=507-3\sigma+3R$ and $R\ge71$ into
\eqref{eq:p3-a0-main}:
\[
T_r+5r+\sigma_t+2L-3\sigma+2e
\le326+w+\delta+B_e.
\]
With no heavy count, \eqref{eq:p3-a0-errors}, the lower bound on $L$,
and $\sigma+\delta\le134+w-I$ force
\[
T_r+5r+3I+\sigma_t+4e\le402+4w\le434.
\]
Its left side is at least $670+32=702$.
With a heavy count, use $e_v\le L$, transfer $qe_v$ to the left,
and substitute the same lower bound with coefficient $2-q\ge0$.
Before the mass substitution the row is
\[
T_r+5r+(2-q)I+\sigma_t-(1+2q)\sigma-(1+q)\delta+4e+qc
\le268+32q+(3-q)w.
\]
Since the two negative terms are at least
$-(1+2q)(134+w-I)+q\delta$, this implies
\begin{equation}\label{eq:p3-a0-m1-heavy}
T_r+5r+(3+q)I+\sigma_t+q\delta+4e+qc
\le402+300q+(4+q)w.
\end{equation}
For $q=1$ the right side is at most 742, while the left is at least
$695+5\cdot8+2\cdot74=883$. For $q=2$ they are at most 1050
and at least $711+6\cdot8+4\cdot74=1055$. Both are strict contradictions.

\paragraph{A0 at $m=2$.}
Again $r=r_t$. If $t$ owns the unique 1 set $(u,a)=(1,1)$;
otherwise set $(u,a)=(2,0)$. Its baseline excess is $ur-a$.
Put $b_0=4u-a+10$, equal to 13 or 18. Then
\[
S_t+2M\ge T_r+(u+4)r+\sigma_t+2R+2L+b_0,
\]
and \eqref{eq:p3-a0-main} gives
\[
T_r+(u+4)r+\sigma_t+2L-3\sigma+2e
\le352-b_0+w+\delta+B_e.
\]
Here $\chi=I-k+35\ge I-12$, since the actual low-1 owner size
$k\le47$, regardless of which physical side it is.
For no heavy count the two substitutions give
\[
T_r+(u+4)r+3\chi+\sigma_t+4e
\le335-b_0+4w\le354.
\]
The left side is at least $670-36+32=666$.
For one heavy count, substituting the $L$ bound with coefficient $2-q$
and then $\sigma+\delta\le101+w-\chi$ gives
\begin{equation}\label{eq:p3-a0-m2-heavy}
T_r+(u+4)r+(3+q)\chi+\sigma_t+q\delta+4e+qc
\le335+267q-b_0+(4+q)w.
\end{equation}
The intermediate right side, before the last substitution, is
$234+65q-b_0+(3-q)w$. For $q=1$ the final right side is at most
629, but the left is at least $695-48+40+148=835$.
For $q=2$ they are at most 904 and at least
$711-60+48+296=995$. This excludes both physical low-1 placements,
including all heavy and light endpoint positions.

\paragraph{A0 at $m=3$.}
Here $\tau=134+\sigma$, $R\ge73$,
$\sigma+\delta+I\le69+w$, and
$L\ge93-w+I+\sigma+\delta$; the latter follows already from
$C\ge6$ and $J\ge\mathcal A-47-w$.
The rank floor is $f=(159-\sigma+R)/4$.
Apply Lemma~\ref{lem:p3-END2} to the lower endpoints $P,\alpha_s,\alpha_t$
and the upper endpoint $H_s$ in
\[
H_s=P+\alpha_s+\alpha_t+d_t+2\mathcal V+G.
\]
With their four selected clipping counts totaling $E_c$, this gives
$S_s\ge d_t+2\mathcal V+G+149-E_c$.
The opposite leading has native loss $\Lambda_t$ and maximum
$d_t-g_t-2X_t$, so $d_t\ge K+g_t+2X_t-\Lambda_t$.
These five resources are distinct, hence $E_c+\Lambda_t\le\epsilon$.
Subtracting from the total source mass and using $K\le M+R+21$ yields
\[
S_t+M+(2X+G)+(2X_t+g_t)\le R+1126+w-e.
\]
The two retained trees used above separately give
$2X+G\ge6f-\epsilon$.
The retained path $\mathcal L_t-V_t-U_t$ and star
$\mathcal L_s-C_{7s}$ give $2X_t+g_t\ge3f-\epsilon$.
These are two independently valid scalar inequalities; adding them
subtracts \emph{two} copies of $\epsilon$, not one. Thus
\[
S_t+M+\tfrac54R-\tfrac94\sigma-2\delta+3e
\le\tfrac{3073}{4}+w.
\]
The physical-side mass and the lower bound on $L$ give, for $r=r_t$,
\[
S_t+M\ge T_r+6r+89+\sigma_t+L
\ge T_r+6r+182-w+I+\sigma+\sigma_t+\delta.
\]
Using $R\ge73$ and then $\sigma+\delta\le69+w-I$ forces
\[
T_r+6r+\tfrac94I+\sigma_t+3e
\le\tfrac{2325}{4}+\tfrac{13}{4}w\le\tfrac{2429}{4}.
\]
But
\[
T_r+6r+\tfrac94(r-34)(r-35)
=666+\frac{(r-27)(11(r-27)-1)}4\ge666,
\]
already larger than $2429/4$. This excludes $m=3$. Together with the
fixed-credit exclusion of $m=4$, all A0 first gaps are impossible for
either actual pair1 owner.

% A1 proof. Requires the common foundations and the P3 profile lemmas.
% Structural rho and kappa are not reused here.
\subsection{A1: the opposite sole owners}\label{subsec:p3-a1}

\begin{proposition}[Exclusion of A1]\label{prop:p3-a1}
In the integral $P_3$ baseline, the following configuration is impossible:
\[
 \{\text{empty whole units}\}=\{O_s,O_t,H_s\},\qquad
 \text{sole beta owner}=s,\qquad
 \text{sole pair1 owner}=t.
\]
Here $s$ is the physical beta owner, with no sorting of the two side
sizes. Every other original whole unit is retained. The conclusion
includes every integral position of the active centers and each first
gap $m\in\{1,2,3,4\}$.
\end{proposition}

\begin{proof}
We keep the original deficit $\delta=c-234$ throughout. Write
\[
 r=r_s,\quad r_t=69-r,\quad T_r=\frac{r(r+1)}2,\quad
 I=(r-34)(r-35),\quad R=R_s+R_t,\quad L=e_s+e_t.
\]
Thus $22\le r_s,r_t\le47$, $0\le\delta\le138$, $2\le w\le z\le8$,
and the actual positive targets have $w/2$ holes each. The source has
$78$ points, $K=|D+D|\ge155$, and
\begin{equation}\label{eq:p3-a1-mass}
 S_s+S_t=2K+1254+w-\delta.
\end{equation}
The four increasing source magnitudes above $R_v$ have offsets
$a_v<b_v<c_v<e_v$, and $M_v=R_v+e_v$. In particular $e_v\ge4$.
All outer coordinates and centers used below are integral.

We abbreviate the ordinary centers by
\[
 U_v=\mathrm{own2}_v=\lambda_v+v_3,\qquad
 V_v=\mathrm{own3}_v=\lambda_v+v_4,\qquad P_t=C_{4t},
\]
and use the literal outer gaps
\[
 d_v=v_1-v_2,\quad g_v=v_3-v_4,\quad h_v=v_2-v_3,
 \quad X_v=v_4-M_v.
\]
The leading resource on target $v$ is denoted by
$\mathcal L_v=\lambda_v+(D_v+D_v)$; its capacity is $K$.
The letter $L$ always denotes the sum of the two fourth source offsets.

\paragraph{The fixed credit and licensed losses.}
Lemma~\ref{lem:p3-fixed} gives $e=11$ for this exact configuration.
Its eight generic labels are $C^v_{11},C^v_{12},C^v_{13},C^v_{14}$
on the two targets; the three additional labels are
$C^s_{22},C^s_{23},C^s_{24}$. They lie strictly above their respective
target maxima: the active $C_{7v}$ forces the empty $O_v$ to be on its
high arm, and the active $\beta_s$ similarly forces $H_s$ to be on
its high arm. The twenty-six nonpositive fixed-label witnesses remain
available, so these eleven labels reduce the fixed allowance to $73$.
Put
\begin{equation}\label{eq:p3-a1-budget}
 e=11,\qquad \varepsilon=\delta-e\ge0.
\end{equation}
Every compatible sum of deficiencies of distinct retained original
resources, or of disjoint groups of them, is at most $\varepsilon$.
This is Lemma~\ref{lem:common-loss} with this single fixed credit.

For an ordinary center $Z$ on target $v$, set
\[
 \Gamma_Z=78-| (Z+D_v)\cap F_v |
          =\ell_Z+h_Z+\iota_Z,
\]
where $\ell_Z$ and $h_Z$ count images below $1$ and above $S_v$, and
$\iota_Z$ counts images at the actual interior holes. The beta pair and
pair1 have actual sole owners $\beta_s$ and $P_t$, respectively; their
whole-unit deficiencies are therefore exactly these owners' native
$\Gamma$ values. Define
\begin{equation}\label{eq:p3-a1-reserved-losses}
 b_*=\Gamma_{\beta_s}+\Gamma_{P_t},\qquad
 h_*=\Gamma_{H_t},\qquad
 b_*+h_*\le\varepsilon.
\end{equation}
A grouping omitting both pairs consequently has available loss at most
$\varepsilon-b_*$; if it also omits $H_t$, the bound is
$\varepsilon-b_*-h_*$. These are remainders in the
same original budget, not new budgets.

\paragraph{The source rows used in this proof.}
We briefly specialize Proposition~\ref{prop:p3-profiles} and
\eqref{eq:p3-rank}. On each side $\sigma_v$ is excess above the
indicated coordinatewise baseline, and $\sigma=\sigma_s+\sigma_t$.
Let $Q$ be the $62$-point core obtained by removing its four largest
magnitudes on each sign and retaining zero; put $Q_s=Q$, $Q_t=-Q$.
For each nonnegative integral $q$,
\begin{equation}\label{eq:p3-a1-overlap}
 |Q\cap(Q+q)|\ge f-q,
\end{equation}
where the required profiles are
\begin{align}
 m=1:\quad &\tau_v=r_v+\sigma_v,\quad
 R_v\in[r_v+1,r_v+1+\sigma_v],\quad R\ge71,\notag\\
 &I+\sigma+\delta\le134+w,\qquad
 f=\frac{169-\sigma+R}{4};\label{eq:p3-a1-profile1}\\
 m=2:\quad &\tau=137-k+\sigma,\quad
 \sigma+\delta+\chi\le101+w,\quad \chi=(k-35)^2,\notag\\
 &f=41+\frac{R-\sigma}{4};\label{eq:p3-a1-profile2}\\
 m=3:\quad &\tau_v=2r_v-2+\sigma_v,\quad R_v\ge r_v+2,\notag\\
 &I+\sigma+\delta\le69+w,\quad
 f=\frac{159-\sigma+R}{4},\quad 2R+I+\delta\le221+w.
 \label{eq:p3-a1-profile3}
\end{align}
In $m=2$, $k$ is the size of the \emph{physical owner of the unique
core entry $1$}. It need not equal $r$. On this owner the baseline is
$(1,3,4,\ldots,k+1)$, giving $\tau_v=r_v-1+\sigma_v$ and
$r_v+1\le R_v\le r_v+1+\sigma_v$; on the other side the baseline is
$(3,4,\ldots,r_v+2)$, giving $\tau_v=2r_v+\sigma_v$ and
$r_v+2\le R_v\le r_v+2+\sigma_v$.
The $m=1$ baseline is $(2,3,\ldots,r_v+1)$, and the $m=3$ baseline
is $(1,4,5,\ldots,r_v+2)$. 

For use in all three profiles we also need
\begin{equation}\label{eq:p3-a1-L}
 L\ge I+\tau+\delta-37-w.
\end{equation}
Here is the source-only derivation, which fixes the constant.
Put $\mathcal A=I+\tau+\delta$, $U=a_s+b_s+a_t+b_t$,
$C=c_s+c_t$ and $J=L-C$. The exact mass row is
\[
 2K=-29-w+\mathcal A+4R+U+2C+J.
\]
Every pair sum is either in the central interval of cardinality
$2R+L+1$ or among the two peeled fringes of sizes $B_s,B_t$.
Lemma~\ref{lem:common-offsets} therefore gives
\[
 K\le2R+L+1+B_s+B_t,\qquad
 \mathcal A-31-w+U-J\le2(B_s+B_t)\le U+C+6.
\]
Cancel $U$ and rearrange to obtain \eqref{eq:p3-a1-L}. Its profile
specializations are
\begin{equation}\label{eq:p3-a1-L-profiles}
 \begin{array}{ll}
 m=1:&L\ge32-w+I+\sigma+\delta,\\
 m=2:&L\ge100-k-w+I+\sigma+\delta,\\
 m=3:&L\ge97-w+I+\sigma+\delta.
 \end{array}
\end{equation}
All intersections below are actual raw intersections of integral
sets. In particular, \eqref{eq:p3-a1-overlap} follows by counting in
the joint integer hull; it does not assert that $Q$ is an interval.

\paragraph{The ten-edge forest for $m=1,2$.}
For these two profiles take on target $s$ the arms
\[
 \mathcal L_s-V_s-U_s,\qquad
 \mathcal L_s-C_{7s}-C_{6s}-C_{5s},
\]
and on target $t$ the arms
\[
 \mathcal L_t-V_t-U_t-H_t,\qquad
 \mathcal L_t-C_{7t}-C_{6t}.
\]
Every ordinary vertex is a distinct retained singleton. For a positive
arm restrict its leading hub to $\lambda_v+M_v+Q_v$; for a negative
arm restrict it to $\lambda_v-M_o+Q_v$, where $o$ is the opposite
side. These are subsets of one and the same full leading vertex on
that target. Restrict ordinary vertices to their $Q_v$ translates.
The ten successive edge gaps have total
\[
 2X_s+2X_t+2g_s+2g_t+2h_t.
\]
For example $H_t-U_t=h_t$ and $C_{6s}-C_{5s}=h_t$.
Each edge contributes an intersection of size at least $f$ minus its
gap. Lemma~\ref{lem:common-forest} charges all ten to the deficiencies
of the displayed full resources. One leading capacity is used per
target, even though it contains two different restricted supports.
Both pairs are omitted, so
\begin{equation}\label{eq:p3-a1-ten-edges}
 X_s+X_t+g_s+g_t+h_t
 \ge5f-\frac{\varepsilon-b_*}{2}.
\end{equation}
Separately use only the four edges in
$\mathcal L_s-V_s-U_s$ and $\mathcal L_t-C_{7t}-C_{6t}$.
They omit both pairs and $H_t$, and their total gap is $2(X_s+g_s)$.
Thus
\begin{equation}\label{eq:p3-a1-four-edges}
 X_s+g_s\ge2f-\frac{\varepsilon-b_*-h_*}{2}.
\end{equation}
Finally restrict the two edges $\mathcal L_s-V_s$ and
$\mathcal L_t-C_{7t}$ to the entire $70$-point core $E$, of span $R$.
Each intersection is at least $139-R-X_s$. The same omissions give
\begin{equation}\label{eq:p3-a1-two-stars}
 X_s\ge139-R-\frac{\varepsilon-b_*-h_*}{2}.
\end{equation}
Adding \eqref{eq:p3-a1-ten-edges}, twice
\eqref{eq:p3-a1-four-edges}, and three times
\eqref{eq:p3-a1-two-stars} yields, for
$W=6X_s+X_t+3g_s+g_t+h_t$,
\begin{equation}\label{eq:p3-a1-W}
 W\ge9f+417-3R-3\varepsilon
             +3b_*+\frac52h_*.
\end{equation}
This is a positive linear combination of three separately valid scalar
inequalities, not a simultaneous charge of their repeated vertices.

\paragraph{Transport to $H_t$ with all endpoint errors exposed.}
For an ordinary center $Z$ on target $v$ choose nonnegative errors
$A_Z,D_Z$ such that
\begin{equation}\label{eq:p3-a1-endpoint-errors}
 Z\ge74-r_v-\ell_Z-A_Z,\qquad
 Z\le S_v+r_o-73+h_Z+D_Z.
\end{equation}
Only $A_{\beta_s},A_{P_t},D_{H_t}$ will be used. Set
\[
 E_{\rm err}=A_{\beta_s}+A_{P_t}+D_{H_t},\qquad
 E_{\rm clip}=\ell_{\beta_s}+\ell_{P_t}+h_{H_t}.
\]
Write $a=s_4$, $b=t_4$, and, only for this calculation,
$\mathcal V=a+b$, $G=g_s+g_t$. The literal center definitions give
\[
 \beta_s=h_s-a-d_s,\qquad
 P_t=d_t-2d_s-\beta_s-2\mathcal V-G,\qquad
 V_t=d_t-b-g_t.
\]
Eliminating $d_t$ and then using $H_t-V_t=g_t+h_t$ gives
\begin{equation}\label{eq:p3-a1-H-identity}
 H_t=2d_s+\beta_s+P_t+2a+b+g_s+g_t+h_t.
\end{equation}
In particular the definition of $\alpha_s$ is $h_s-a$, with no
additional $g_s$. Apply the two lower and one upper rows of
\eqref{eq:p3-a1-endpoint-errors}. Since $r_s+r_t=69$,
\begin{equation}\label{eq:p3-a1-H-endpoints}
 2d_s+2a+b+g_s+g_t+h_t
 \le S_t+r-152+E_{\rm clip}+E_{\rm err}.
\end{equation}
Let $\Lambda_s$ be the native deficiency of the distinct leading
resource $\mathcal L_s$. The three ordinary resources just used are
also distinct whole resources, so
$E_{\rm clip}+\Lambda_s\le\varepsilon$.
The maximum of $\mathcal L_s$ is
\[
 \lambda_s+2M_s=d_s-g_s-2X_s.
\]
Its $K-\Lambda_s>0$ target hits are positive integers, since
$\Lambda_s\le\delta<155\le K$. Consequently
\begin{equation}\label{eq:p3-a1-leading}
 d_s\ge K+g_s+2X_s-\Lambda_s.
\end{equation}
Insert this lower bound twice into \eqref{eq:p3-a1-H-endpoints}, use
\eqref{eq:p3-a1-mass}, and use
$E_{\rm clip}+2\Lambda_s\le2\varepsilon$. The resulting transport is
\begin{equation}\label{eq:p3-a1-H-transport}
 S_s+3g_s+4X_s+2a+b+g_t+h_t
 \le1102+w+r+\delta-2e+E_{\rm err}.
\end{equation}
The coefficient two on $\Lambda_s$ arises from using a scalar
inequality twice; the leading resource was counted once in the budget.
Interior target holes are already present in the native deficiencies.

\paragraph{Uniform endpoint errors for $m=1,2$.}
Call one of the three relevant endpoints \emph{light} if its clipping
count $q$ is less than $74$, and \emph{heavy} if $q\ge74$. Here
$q=\ell_Z$ for a lower endpoint and $q=h_Z$ for the upper endpoint.
The two profile caps and \eqref{eq:p3-a1-budget} imply
\begin{equation}\label{eq:p3-a1-light-caps}
 \begin{array}{ll}
 m=1:&\sigma+\varepsilon\le134+w-e-I\le131,\\
 m=2:&\sigma+\varepsilon\le101+w-e-\chi\le98.
 \end{array}
\end{equation}
Thus $\varepsilon<148$, and at most one of the three endpoints can be
heavy, by \eqref{eq:p3-a1-reserved-losses}.

If $Z\ge1$, its lower error can be zero: the $r_o+4$ negative source
values avoiding lower clipping occupy at most $Z-1$ integer positions,
which gives the claimed lower bound. If $Z\le S_v$, the analogous
positive-source count gives upper error zero.
For a light lower endpoint with $Z\le0$, all negative values and zero
are clipped. If $j$ positive values are additionally clipped, then
\[
 q=74-r_v+j<74,\qquad j<r_v,\qquad q\ge r_o+5\ge27.
\]
The next positive value $x_{j+1}$ is a core entry and satisfies
$Z\ge1-x_{j+1}$. Relative to the nominal lower endpoint $-j$, the
required error is at most $x_{j+1}-j-1$.
For $m=1$ this equals $1+d$ if its displacement above baseline is $d$.
For $m=2$ it is at most $2+d$. If the next rank is the forced $1$, the
assertion holds directly with $d=0$; otherwise all later baseline
ranks are consecutive, so monotonicity of their displacements gives
\[
 \sigma_v\ge d(r_v-j)=d(74-q).
\]
Because $\sigma_v+q\le\sigma+\varepsilon\le131$,
\begin{equation}\label{eq:p3-a1-displacement}
 d\le\frac{131-q}{74-q}\le q\qquad(27\le q\le73).
\end{equation}
For the last step use $q(75-q)\ge131$: this concave quadratic has
minimum on $[27,73]$ at an endpoint, whose values are $1296$ and
$146$. An upper endpoint with $Z>S_v$ is treated by reflecting the
source. If $j$ negative magnitudes are additionally upper-clipped,
then $q=74-r_o+j$ and $Z\le S_v+y_{j+1}$. Its error is bounded by
the same $1+d$ or $2+d$ calculation, now on the opposite core side.
Thus every light relevant endpoint admits
\[
 A_Z\text{ or }D_Z\ \le
 \begin{cases}1+q\le1+\Gamma_Z,&m=1,\\
               2+q\le2+\Gamma_Z,&m=2.
 \end{cases}
\]
If all three are light, this gives
\begin{equation}\label{eq:p3-a1-all-light}
 E_{\rm err}\le
 \begin{cases}3+b_*+h_*,&m=1,\\
               6+b_*+h_*,&m=2.
 \end{cases}
\end{equation}

For a heavy lower endpoint, activity alone gives $Z\ge1-M_v$;
for a heavy upper endpoint it gives $Z\le S_v+M_o$. In either case,
with $u=v$ in the lower case and $u=o$ in the upper case, a
nonnegative error no larger than
\[
 M_u-r_u+73-q\le M_u-r_u-1
\]
suffices. This quantity is nonnegative here (an active $78$-point
ordinary translate has clipping count at most $77$, and
$M_u\ge r_u+5$ in the relevant profiles). The baseline maximum bounds
therefore give heavy error at most $e_u+\sigma_u$ for $m=1$ and
at most $e_u+\sigma_u+1$ for $m=2$.
Let $G_{\rm light}$ be the sum of the native deficiencies of the other
two relevant centers. Combining their light bounds with the heavy one,
\begin{equation}\label{eq:p3-a1-one-heavy}
 E_{\rm err}\le
 \begin{cases}e_u+\sigma_u+2+G_{\rm light},&m=1,\\
               e_u+\sigma_u+5+G_{\rm light},&m=2.
 \end{cases}
\end{equation}
The heavy center has native deficiency at least $74$. Its coefficient
in $3b_*+\tfrac52h_*$ is at least $5/2$, while the
other two coefficients remain nonnegative after subtracting one copy.
Hence
\begin{equation}\label{eq:p3-a1-heavy-credit}
 3b_*+\frac52h_*-G_{\rm light}\ge185,
 \qquad L-e_u=e_o\ge4,\qquad
 \sigma_s-\sigma_u\ge-\sigma.
\end{equation}
The possibly large peeled offset has been retained, to be cancelled
against $L$, rather than assigned an unjustified uniform bound.

\paragraph{Source mass and the three $m=1,2$ scalar rows.}
The four positive peeled magnitudes on side $s$ have total at least
$M_s+3R_s+6$. With $a=M_s+X_s$, $b=M_t+X_t$ this gives
\begin{align}
 S_s+2a+b+4X_s
 &\ge T_r+\tau_s+5R_s+R+L+2e_s+6+6X_s+X_t.
 \label{eq:p3-a1-source-transport}
\end{align}
Use $e_s\ge4$ and the respective baseline rows. The right side is at
least, in the three cases,
\[
 \begin{array}{ll}
 m=1:&T_r+6r+19+\sigma_s+R+L+6X_s+X_t,\\
 m=2,\ k=r:&T_r+6r+18+\sigma_s+R+L+6X_s+X_t,\\
 m=2,\ k=69-r:&T_r+7r+24+\sigma_s+R+L+6X_s+X_t.
 \end{array}
\]
Add $3g_s+g_t+h_t$, use \eqref{eq:p3-a1-W}, and compare with
\eqref{eq:p3-a1-H-transport}. Substituting $\varepsilon=\delta-e$
yields exactly
\begin{align}
 m=1:\quad
 T_r+5r+\sigma_s-\frac94\sigma+\frac R4+L+5e
      +3b_*+\frac52h_*
 &\le\frac{1143}{4}+w+4\delta+E_{\rm err};
 \label{eq:p3-a1-scalar1}\\
 m=2,\ k=r:\quad
 T_r+5r+\sigma_s-\frac94\sigma+\frac R4+L+5e
      +3b_*+\frac52h_*
 &\le298+w+4\delta+E_{\rm err};
 \label{eq:p3-a1-scalar2s}\\
 m=2,\ k=69-r:\quad
 T_r+6r+\sigma_s-\frac94\sigma+\frac R4+L+5e
      +3b_*+\frac52h_*
 &\le292+w+4\delta+E_{\rm err}.
 \label{eq:p3-a1-scalar2t}
\end{align}
For example, the constant in $9f+417$ is $3189/4$ for $m=1$ and
$786$ for $m=2$. After adding $19,18,24$, respectively, subtraction
from $1102$ gives precisely $1143/4,298,292$. The net $R$ coefficient
is $1+9/4-3=1/4$. No radial upper bound is used for $m=1,2$.

\paragraph{Exclusion when $m=1$.}
If all three endpoints are light, insert \eqref{eq:p3-a1-all-light}
and the first row of \eqref{eq:p3-a1-L-profiles} into
\eqref{eq:p3-a1-scalar1}, and drop the nonnegative remaining native
credits. First this gives
\[
 T_r+5r+I+\sigma_s-\frac54\sigma+\frac R4+5e
 \le\frac{1027}{4}+2w+3\delta.
\]
Now $\delta\le134+w-I-\sigma$, so
\begin{equation}\label{eq:p3-a1-m1-light-final}
 T_r+5r+4I+\frac74\sigma+\sigma_s+\frac R4+5e
 \le\frac{2635}{4}+5w\le\frac{2795}{4}.
\end{equation}
For every integer $r$,
\begin{equation}\label{eq:p3-a1-q1}
 T_r+5r+4I=695+\frac{(r-30)(9r-271)}2\ge695.
\end{equation}
Indeed both factors are negative for $r\le29$, the first is zero at
$r=30$, and both are positive for $r\ge31$. Since $R\ge71$, even
these two terms give $695+71/4=2851/4>2795/4$, a strict contradiction.

If there is one heavy endpoint, insert \eqref{eq:p3-a1-one-heavy}
into \eqref{eq:p3-a1-scalar1} and apply
\eqref{eq:p3-a1-heavy-credit}. Cancelling $e_u$ against $L$ gives
\[
 T_r+5r-\frac{13}{4}\sigma+\frac R4+5e+189
 \le\frac{1151}{4}+w+4\delta.
\]
The same mass cap yields
\begin{equation}\label{eq:p3-a1-m1-heavy-final}
 T_r+5r+4I+\frac34\sigma+\frac R4+5e
 \le\frac{2539}{4}+5w\le\frac{2699}{4}<695,
\end{equation}
contradicting \eqref{eq:p3-a1-q1}. These alternatives cover all
integral endpoints, including clipping into the peeled tail.

\paragraph{Exclusion when $m=2$: both physical low-$1$ owners.}
Suppose first that all endpoints are light. Use
\eqref{eq:p3-a1-all-light}, the second row of
\eqref{eq:p3-a1-L-profiles}, and
$\delta\le101+w-(k-35)^2-\sigma$ in
\eqref{eq:p3-a1-scalar2s} or \eqref{eq:p3-a1-scalar2t}.
If $k=r$, the result is
\begin{equation}\label{eq:p3-a1-m2s-light-final}
 T_r+4r+I+3(r-35)^2+\frac74\sigma+\sigma_s+\frac R4+5e
 \le507+5w\le547.
\end{equation}
If $k=69-r$, it is
\begin{equation}\label{eq:p3-a1-m2t-light-final}
 T_r+7r+I+3(r-34)^2+\frac74\sigma+\sigma_s+\frac R4+5e
 \le570+5w\le610.
\end{equation}
The exact integral quadratic identities are
\begin{align}
 T_r+4r+I+3(r-35)^2
 &=680+\frac92(r-30)(r-31)\ge680,\label{eq:p3-a1-q2s-light}\\
 T_r+7r+I+3(r-34)^2
 &=743+\frac92(r-29)(r-30)\ge743.\label{eq:p3-a1-q2t-light}
\end{align}
The products of consecutive integers are nonnegative. Thus each row
contradicts its corresponding strict upper bound.

If there is one heavy endpoint, use \eqref{eq:p3-a1-one-heavy} and
\eqref{eq:p3-a1-heavy-credit} instead. The intermediate rows are
\[
 \begin{array}{ll}
 k=r:&T_r+5r-\frac{13}{4}\sigma+\frac R4+5e+189
       \le303+w+4\delta,\\
 k=69-r:&T_r+6r-\frac{13}{4}\sigma+\frac R4+5e+189
       \le297+w+4\delta.
 \end{array}
\]
After the corresponding mass cap these become
\begin{align}
 T_r+5r+4(r-35)^2+\frac34\sigma+\frac R4+5e
 &\le518+5w\le558,\label{eq:p3-a1-m2s-heavy-final}\\
 T_r+6r+4(r-34)^2+\frac34\sigma+\frac R4+5e
 &\le512+5w\le552.\label{eq:p3-a1-m2t-heavy-final}
\end{align}
But
\begin{align}
 T_r+5r+4(r-35)^2
 &=715+\frac92(r-30)(r-31)\ge715,\label{eq:p3-a1-q2s-heavy}\\
 T_r+6r+4(r-34)^2
 &=709+\frac92(r-29)(r-30)\ge709.\label{eq:p3-a1-q2t-heavy}
\end{align}
Again all other left terms are nonnegative. Both orientations are
impossible. No interchange of physical beta and low-$1$ ownership has
been used.

\paragraph{Exclusion when $m=3$: uniform endpoint transport.}
Use the third profile in \eqref{eq:p3-a1-profile3}. It gives
\[
 \sigma+\varepsilon\le69+w-e-I\le66<74.
\]
Therefore Lemma~\ref{lem:p3-END2} applies to every retained singleton
or actually sole paired owner used here. Its two endpoint bounds are
\begin{equation}\label{eq:p3-a1-m3-END2}
 Z\ge74-r_v-\ell_Z-2,\qquad
 Z\le S_v+r_o-73+h_Z+2.
\end{equation}
They include $Z\le0$ and $Z>S_v$. In either exceptional position a
positive displacement of the next relevant core rank above the
$m=3$ baseline would force $\sigma+\varepsilon\ge74$, which has
just been excluded. There is no separate interior-hole penalty.

This time use the literal identity for $V_t$ rather than $H_t$:
\begin{equation}\label{eq:p3-a1-m3-V-identity}
 V_t=2d_s+\beta_s+P_t+2a+b+g_s.
\end{equation}
Apply the lower rows of \eqref{eq:p3-a1-m3-END2} to $\beta_s,P_t$
and the upper row to $V_t$. Their six units of endpoint allowance
yield
\begin{equation}\label{eq:p3-a1-m3-endpoints}
 2d_s+2a+b+g_s\le S_t+r-146+E_0,\qquad
 E_0=\ell_{\beta_s}+\ell_{P_t}+h_{V_t}.
\end{equation}
The beta pair, pair1, and $V_t$ are three distinct whole resources.
Together with the distinct leading resource they give
$E_0+\Lambda_s\le\varepsilon$.
Insert \eqref{eq:p3-a1-leading} twice into
\eqref{eq:p3-a1-m3-endpoints}, then use \eqref{eq:p3-a1-mass}:
\begin{align}
 S_s+3g_s+4X_s+2a+b
 &\le1108+w+r-\delta+E_0+2\Lambda_s\notag\\
 &\le1108+w+r+\delta-2e.
 \label{eq:p3-a1-m3-transport}
\end{align}

For the ordinary gaps choose the two disjoint resource pairs
$\{U_s,V_s\}$ and $\{C_{6t},C_{7t}\}$, on their respective targets.
Both center gaps equal $g_s$. Each pair's deficiency dominates its
raw $Q$-overlap, so \eqref{eq:p3-a1-overlap} and the credited budget
imply
\begin{equation}\label{eq:p3-a1-m3-gap}
 2(f-g_s)\le\varepsilon,\qquad
 3g_s\ge3f-\frac32\varepsilon.
\end{equation}
These four vertices are retained singletons.

For the stars use, on each target $v$, the leading hub
$\mathcal L_v$ with leaves $V_v,C_{7v}$. Restrict the two hub supports
to $\lambda_v+M_v+E_v$ and $\lambda_v-M_o+E_v$, respectively.
The matching leaf intersections are at least $139-R-X_v$ and
$139-R-X_o$. The pointwise forest bound counts one full leading
vertex despite its two restricted supports. Using both stars gives
\[
 2(X_s+X_t)+4R+\varepsilon\ge556.
\]
Separately, using only the two edges whose gap is $X_s$ gives
\[
 2X_s+2R+\varepsilon\ge278.
\]
Half of the first scalar row plus five halves of the second gives
\begin{equation}\label{eq:p3-a1-m3-star}
 6X_s+X_t\ge973-7R-3\varepsilon.
\end{equation}
Neither the ordinary gap groups nor the stars use either pair. All
three rows are separately valid consequences of the same budget.

The $m=3$ mass baseline, $R_s\ge r+2$, and $e_s\ge4$ give from
\eqref{eq:p3-a1-source-transport}
\begin{equation}\label{eq:p3-a1-m3-source}
 S_s+2a+b+4X_s
 \ge T_r+7r+22+\sigma_s+R+L+6X_s+X_t.
\end{equation}
Insert the $m=3$ row of \eqref{eq:p3-a1-L-profiles},
\eqref{eq:p3-a1-m3-gap}, and \eqref{eq:p3-a1-m3-star} into this
lower bound, and compare with \eqref{eq:p3-a1-m3-transport}. With
$f=(159-\sigma+R)/4$, cancellation of the matching $\delta$ terms
gives the exact row
\begin{equation}\label{eq:p3-a1-m3-algebra1}
 T_r+6r+I+\frac14\sigma+\sigma_s+\frac{413}{4}-2w+2e
 \le\frac{21}{4}R+\frac92\varepsilon.
\end{equation}
For clarity, its constant is
$22+97+973+3(159/4)-1108=413/4$; its radial coefficient before
transposing is $1-7+3/4=-21/4$; and its excess coefficient is
$1-3/4=1/4$. This explicitly accounts for all terms of the
$m=3$ source transport.

Use $\varepsilon=\delta-e$ and the $m=3$ specialization of
\eqref{eq:p3-radial}, namely $R\le(221+w-I-\delta)/2$.
Equivalently, in radial-excess notation, $2\xi+I+\delta\le75+w$
with $\xi=R-73$. Then \eqref{eq:p3-a1-m3-algebra1} becomes
\begin{equation}\label{eq:p3-a1-m3-algebra2}
 T_r+6r+\frac{29}{8}I+\frac14\sigma+\sigma_s+\frac{13}{2}e
 \le\frac{3815}{8}+\frac{37}{8}w+\frac{15}{8}\delta.
\end{equation}
The coefficients are obtained by $1+21/8=29/8$,
$2+21/8=37/8$, $9/2-21/8=15/8$, $2+9/2=13/2$, and
$21\cdot221/8-413/4=3815/8$.
Finally substitute $\delta\le69+w-I-\sigma$ from the $m=3$ mass
profile. This gives
\begin{equation}\label{eq:p3-a1-m3-algebra3}
 T_r+6r+\frac{11}{2}I+\frac{17}{8}\sigma+\sigma_s+\frac{13}{2}e
 \le\frac{2425}{4}+\frac{13}{2}w\le\frac{2633}{4}.
\end{equation}
Indeed $29/8+15/8=11/2$, $1/4+15/8=17/8$,
$(3815+15\cdot69)/8=2425/4$, and $(37+15)/8=13/2$.
On the other hand,
\begin{equation}\label{eq:p3-a1-m3-quadratic}
 T_r+6r+\frac{11}{2}I
 =6r^2-373r+6545
 =748+(r-31)(6r-187)\ge748.
\end{equation}
For integral $r\le30$ both factors in the last product are negative;
at $r=31$ the product is zero; and for $r\ge32$ both are positive.
The other left terms in \eqref{eq:p3-a1-m3-algebra3} are
nonnegative, whereas $748>2633/4$. This excludes every $m=3$
endpoint and both physical size orientations.

\paragraph{Exclusion when $m=4$ and conclusion.}
Proposition~\ref{prop:p3-profiles} gives $z=w=2$ and $\delta\le3$
for both possible physical placements of the low multiset $\{1,2\}$.
But \eqref{eq:p3-a1-budget} gives $\delta\ge11>3$. No endpoint
restriction is needed in this last case. The four first-gap values
exhaust the stipulated $P_3$ baseline, proving the proposition.
\end{proof}


The activity classification exhausts the original $P_3$ stratum by A0
and A1, and the first-gap argument exhausts each by $m=1,2,3,4$.
Their exclusions therefore prove that $P_3$ is impossible.


\section{Four missing units with positive core holes: $P_4$}\label{sec:p4}
We prove the original $j=4$ stratum directly. The beta sum partitions it
into branch A, $\beta_s+\beta_t\le0$, and branch B,
$\beta_s+\beta_t>0$. In A, distinguish whether the whole beta pair is empty
(A0) or retained (A1). Initially a retained pair may have two contributing
components. The proof below establishes every sole-owner assertion before
using an individual native loss.

\subsection{The original two-hole source baseline}\label{subsec:p4-baseline}
Put $A=I+\tau+\delta$, reserving $C_0$ for total core mass.
The original structural budget, with $c=312+\delta$, gives
\[
 \delta\le157-I-n(r_s)-n(r_t).
\]
If a size is at most 20, then $I\ge210$ and the defect floors sum to 50;
if it is between 21 and 24, then $I\ge110$ and they sum to 59.
Both are impossible. Thus $25\le r_s,r_t\le44$, $\rho\ge59$, and
$\mu\ge63$. The stronger structural row gives
\[
 \rho+\kappa+\delta\le129+3z-I,\qquad
 \mathcal H:=h(P_s)+h(P_t)=\kappa+z\le70+4z-I-\delta\le126.
\]
Both side masses exceed $\mathcal H$. The common-strip lemma applies to
the core and then successively to the two peels, because the layer rows
imply
\[
 C_0-U_0-2\mathcal H\ge380+4I+3\delta-10z\ge240,
 \quad C_1-V_0-2\mathcal H\ge630+4I+3\delta-6z\ge546.
\]
Consequently $z,w$ are even and $0\le w\le z\le14$.

Write the four increasing peeled offsets above $R_v$ as
$a_v<b_v<c_v<e_v$, and set
\[
 U=a_s+b_s+a_t+b_t\ge6,\quad C=c_s+c_t,\quad
 J=(e_s-c_s)+(e_t-c_t),\quad L=C+J=M-R.
\]
The source mass and original cover identity give
\begin{equation}
 2K=49-w+A+4R+U+2C+J.                         \label{eq:p4-initialmass}
\end{equation}
Deleting both source maxima leaves 76 points. Its sums are covered by
an interval of cardinality $2R+C+1$ and at most twelve remaining
same-sign peeled-pair labels. Reinserting the maxima costs at most 155
labels, so $K\le2R+C+168$. The full-source central interval and the
four-offset lemma give $K\le2R+C+J+1+B_s+B_t$ and
$2(B_s+B_t)\le U+16$. Hence
\[
 J+U\le287+w-A,\qquad J\ge31-w+A,\qquad A\le125+w\le139.
                                                        \tag{P4-mass}
\]
In particular $C_0\le1364$, $\tau\le139$, and
$R_v\le r_v+\tau_v\le183$.

Let $m$ be the first missing positive subset sum of the absolute core.
The greedy rule gives total $m-1$ below $m$, with no entry equal to $m$.
If $k_v$ entries of side $v$ are below $m$, with mass $s_v$, then
\[
 \tau_v\ge s_v+(m-k_v)r_v+\frac{k_v(k_v-1)}2-mk_v
 \ge(m-k_v)(r_v-k_v).                         \tag{P4-low-rank}
\]
Since $m<183$, $T_{k_v}\le181$ gives $k_v\le18$, $r_v-k_v\ge7$.
At $m=5,6$, $k_v\le2$ and $k_s+k_t\le3$, giving respectively
$\tau\ge222$ and $\tau\ge264$. For $m\ge7$, a side with $k_v\le3$
has $\tau_v\ge4(r_v-3)$; with $k_v\ge4$, use
$m-k_v\ge T_{k_v-1}+1\ge k_v+3$ and
\[
 (k_v+3)(r_v-k_v)-4(r_v-3)
 \ge3k_v+5>0.
\]
Thus $\tau\ge252$. At $m=4$, the low multiset is $\{1,2\}$; split
placement gives $\tau\ge202$, while together on a size-$r$ side gives
$I+\tau\ge202+(r-35)(r-36)\ge202$.
All contradict $A\le139$. Thus $m=1,2,3$.

The following small supports finish the reduction. Here $T_0$ is the seed
mass, $H_0$ its lower hole set, $h=\max H_0$, and $L_9$ a lower bound for
the mass of any nine selected core occurrences.
\begin{center}\small
\begin{tabular}{l|l|r|l|r}
condition&seed&$T_0$&$H_0$&$L_9$\\\hline
$m=1$, 2 absent&$3,4,5,6$&18&$\{1,2\}$&43\\
$m=1$, 3 absent&$2,4,5,6$&17&$\{1,3\}$&41\\
$m=1$, 4 absent&$2,3,5,6$&16&$\{1,4\}$&39\\
$m=1$, $2,3,4,5$ present&$2,3,4,5$&14&$\{1\}$&34\\
$m=1$, $2,3,4,6$ present&$2,3,4,6$&15&$\{1\}$&34\\
$m=2$&$1,3,4,5$&13&$\{2\}$&37\\
$m=3$&$1,1,4,5,6$&17&$\{3\}$&39
\end{tabular}
\end{center}
Each seed has support $[0,T_0]\setminus(H_0\cup(T_0-H_0))$, as direct
addition up to $T_0/2$ and reflection verify. To see these rows exhaust
all cases, at $m=1$ two absent values $u,v\in\{2,\ldots,6\}$ would force
$\tau\ge207+12-2(1+u+v)\ge195$. Thus at most one is absent.
At $m=2$, a missing $q\in\{3,4,5\}$ would force
$\tau\ge211-r-2q\ge157$, where $r\le44$ owns the unique 1.
At $m=3$, the low entries are the two opposite-sign 1s; a missing
$q\in\{4,5,6\}$ forces $\tau\ge209-2q\ge197$.
The nine-entry bounds come from taking at most two copies of the smallest
allowed values, with the prescribed one or two 1s.

Adjoin the remaining actual occurrences in increasing order. If $n$
have been selected, of total $T$, joining the corrected central intervals
requires $x\le T-2h-1$ for the next entry. If this fails, every unselected
entry is at least $T-2h$. Any $N$ integers at least $x$, of multiplicity at
most two, have mass at least $Nx+\lfloor(N-1)^2/4\rfloor$.
For $n\le8$ use $N\ge61$; for $9\le n\le32$ use $N\ge37$ and
$T\ge L_9$. In the seven table rows the resulting core-mass lower bounds
are respectively
\[
\begin{array}{c|rrrrrrr}
n\le8&1772&1588&1404&1646&1708&1462&1588\\
9\le n\le32&1810&1660&1510&1542&1542&1582&1584
\end{array}
\]
all exceeding 1364. For $n\ge33$, $T\ge T_{16}+T_{17}=289$ and
$x\le183<289-9\le T-2h-1$. Every adjoining step is therefore valid.
The final support has at most $2|H_0|$ holes; a seed hole may fill.

In the first three rows $z\le4$, so (P4-mass) gives $A\le129$.
Yet absence of $1,f$, $f=2,3,4$, gives
$\tau\ge142-2f\in\{138,136,134\}$, impossible. The remaining $m=1,2$
rows have precisely two surviving holes, $m,C_0-m$. The $m=3$ row would
give $z\le2$, $A\le127$, contrary to $\tau\ge134$.
Since all peeled magnitudes exceed $m\le2$, the first gap survives in
$D$. Thus $w=z=2$, and translating the absolute support gives
\begin{equation}
 F_v=[1,S_v]\setminus\{S_v-m\},\qquad m\in\{1,2\},\qquad
 A\le127.                                           \label{eq:p4-target}
\end{equation}
At $m=1$, $\tau\ge69$, so $\delta\le58-I$.
At $m=2$, with unique-1 owner size $r$,
$I+\tau\ge102+(r-35)^2$, hence
$\delta\le25-(r-35)^2$.
Nonnegative $\delta$ forces $I\le58$, and therefore
\begin{equation}
 27\le r_v\le42,\quad \delta\le58-I,\quad
 R+\delta\le129-I.                                \label{eq:p4-baseline}
\end{equation}
For the last inequality use $\tau\ge R-2$ at $m=1$.
At $m=2$, if the other side has size $q$, the sharper
$\tau\ge R+q-4$ gives $R+\delta\le106-I$ already from $q\ge25$.
No stronger core-mass cap is assumed.

\subsection{Activity and outer integrality before classification}\label{subsec:p4-activity}
Write $U_v=\mathrm{own2}_v$, $V_v=\mathrm{own3}_v$ and retain the other
common center names. Leading loss $\ell_v\le\delta<K$ proves activity,
integral $\lambda_v$, and
\begin{equation}
 d_v\ge K+g_v+2X_v-\ell_v.                         \label{eq:p4-lead}
\end{equation}
The high-tail inequality and exact mass $S=2K+1178-\delta$ imply
\begin{equation}
 d_s+d_t+G+\delta-2K\ge494.                        \label{eq:p4-tail}
\end{equation}
Also \eqref{eq:p4-initialmass} with $w=2$ and the two-max estimate gives
$K-M\ge R-115+A\ge23+I+\delta$.

If $O_v$ is active, its lower source endpoint $O_v-M_o$ is integral and
at most $S_v$. The other nonfixed resources end at or below $H_v+M_v$,
which exceeds the positive leading maximum. Even if $H_v$ is not yet
integral, the raw55 gap has at least $d_v-M-1$ integer positions.
The one target hole gives $d_v\le M+57$.
If both $O$ were active, set $Y=2M-2K+114+\delta$. The leading and tail
rows give $G+2X\le Y$, $G\ge494-Y$, hence $2Y>494$; but $Y\le68$.
Thus they cannot both be active.

In fact neither is active. If $O_v$ is active, the preceding rows give
$2X_v+g_v\le172-R-I-\tau$.
For a lower coordinate $u=v_3$ or $v_4$, its witness singletons are
respectively $\{U_v,C_{6o}\}$ or $\{V_v,C_{7o}\}$.
Activity of either witness makes $u$ integral. The two actual 70-point
core overlaps at displacement $q=u-M_v$ each have size at least
$139-R-q$. Their distinct leading groups cost at most $\delta$,
including a possibly empty mate with zero capacity. Therefore
$q\ge139-R-\delta/2$. Since $q\le2X_v+g_v$, this forces
$I+\tau\le33+\delta/2\le62$, contradicting $\tau\ge69$.
All four witnesses would therefore be empty, exhausting $j=4$ and
forcing both $O$ active, already excluded. We conclude
\begin{equation}
 O_s,O_t\text{ are empty}.                         \label{eq:p4-O}
\end{equation}

Only two further unit holes remain. A nonintegral lower coordinate
forces two witness holes; two such coordinates are impossible. Exactly
one would additionally force $H_v$ or $\alpha_v$ empty through
$H_v-\alpha_v=\lambda_v+v_3+v_4$, again impossible.
All lower coordinates are integral, so $d_v$ is integral.
If $v_2$ were nonintegral, $H_v,\alpha_v$ would both be missing.
In A, the opposite $C_{5o}$ is a third missing \emph{singleton}; in B,
$\beta_v=v_2-d_v-v_3-v_4$ is a third missing singleton.
Both contradict the two remaining holes. Thus all eight outer
coordinates are integral in both branches, with the correct branch roles.

We record the precise empty-center rule. Source bounds give
$S_v\ge496$, $R_o\le129$ and first opposite peel
$x_{o,1}\le416-R_v-I-\delta\le389$.
The points $M_v,0,-R_o$ prove activity of every integral center in
$[1-M_v,S_v+R_o]$ except possibly $S_v+R_o-m$; at that point $-x_{o,1}$
produces a positive hit at least $132$, below $S_v-m$.
Hence every empty integral center satisfies
\begin{equation}
 Z\le-M_v\quad\text{or}\quad Z\ge S_v+R_o+1.       \label{eq:p4-empty}
\end{equation}
An empty own center cannot be low, by positive leading activity.
Every higher retained singleton then loses at least 74 native points,
contradicting $\delta\le58$. An empty cross center cannot be high,
since that would put the 78-point translate $\lambda_v+D_v$ above the
target. In particular actual cross activity is upward closed:
$C_{4v}\Rightarrow C_{5v}\Rightarrow C_{6v}\Rightarrow C_{7v}$.

We also need forests before full CORE8 retention. Delete $k$ largest
core magnitudes per sign. Monotone rank offsets give
\[
 |Q|=70-2k,\quad\operatorname{span}(Q)\le69-2k+\tau/(k+1),\quad
 |Q\cap(Q+q)|\ge70-2k-\tau/(k+1)-q.
\]
For each physical sign the two arms are $L_v-V_v-U_v$ and
$L_o-C_{7o}-C_{6o}$. At least one has at most one empty ordinary vertex,
since both $O$ are empty and only two other whole units can be missing.
Take its two edges, the other $X_v$ edge, and an $X_o$ edge for the
four-edge row. For the five-edge row take one two-edge arm on each sign
and the unused $X_v$ edge. Positive and negative restricted arms on any
target occupy disjoint zones since $M>R$, so the target-independent empty-vertex forest rule from Section~\ref{sec:separated-arms}
applies without duplicating the leading. The choices $k=3,4$
respectively give
\begin{align}
 2X_v+X_o+g_v&\ge256-\tau-\delta\ge129+I,\label{eq:p4-four}\\
 2X_v+X_o+G&\ge310-\tau-\delta\ge183+I.\label{eq:p4-five}
\end{align}
If $\beta_v,H_v$ are active, then
$h_v-d_v\ge X_v+1$, $\alpha_v>U_v$, and the alpha-to-$H$ raw gap gives
$d_v\le M+57$. An active $C_{6v}$, or an active $C_{7v}$, would give via
\eqref{eq:p4-five} or \eqref{eq:p4-four} respectively
$d_v\ge M+184+I$ or $d_v\ge M+130+I$.
Thus neither triple $\{\beta_v,H_v,C_{6v}\}$ nor
$\{\beta_v,H_v,C_{7v}\}$ can be wholly active.

\subsection{One quantitative toolkit for all remaining $P_4$ cases}\label{subsec:p4-tools}
Assume CORE8, the eight singletons $U_v,V_v,C_{6v},C_{7v}$, is retained.
Below this condition is verified separately in A0, A1 and B.
Let $e\ge0$ be a valid additional fixed credit, and put
$\varepsilon=\delta-e\ge0$; all compatible native or grouped losses
are at most $\varepsilon$, while mass identities retain $\delta$.
Set $\sigma_v=\tau_v-r_v$, $\sigma=\tau-69$, and
$\eta=0$ at $m=1$, $\eta=1$ at $m=2$.
In the latter case $\sigma_v\ge-1$, with only the unique-1 side possibly
negative; $\sigma\ge26$. At $m=1$ all $\sigma_v\ge0$.
The four-offset inequalities and exact mass give
\begin{align}
 2K&=116+I+\sigma+\delta+4R+U+2C+J,\label{eq:p4-mass116}\\
 K&\le M+R+21,& J&\ge98+I+\sigma+\delta,\nonumber\\
 L&\ge108+I+\sigma+\delta=39+I+\tau+\delta.\label{eq:p4-L}
\end{align}
For the last two rows, compare the central-fringe cover with the mass
identity, using respectively $2B_v\le a_v+b_v+8$ and
$2B_v-(a_v+b_v+c_v)\le3$.

Let $p_v$ be the penultimate core magnitude. Its smaller prefix has mass
at least $T_{r_v-1}-1$, except for one extra unit of deficit at the unique
1. Hence $p_s+p_t\le140+\sigma+\eta-R$.
After deleting just the two source maxima, the interval with endpoints
$-(R_t+c_t+p_t)$ and $R_s+c_s+p_s$ covers all sums except at most twenty
same-sign pair labels. Reinsertion costs 155, giving
$K\le R+C+p_s+p_t+176\le C+316+\sigma+\eta$.
Comparison with \eqref{eq:p4-mass116} gives
\begin{equation}
 2R+I+\delta\le206+\eta.                          \label{eq:p4-radial}
\end{equation}

Delete the four largest core magnitudes per sign, leaving $Q_{62}$.
The maximum-deleted prefix excess is
$\sigma_v-R_v+r_v+1$, with one additional unit on the unique-1 side.
Its last four deviations dominate the rank $r_v-4$ deviation. Thus
\begin{equation}
 \operatorname{span}(Q_{62})\le63+\frac{\sigma-R+71+\eta}{4},\qquad
 f_\eta=\frac{169-\eta-\sigma+R}{4},\qquad
 |Q_{62}\cap(Q_{62}+q)|\ge f_\eta-q.              \label{eq:p4-Q62}
\end{equation}
The four disjoint ordinary pairs are
$(U_s,V_s),(C_{6t},C_{7t})$ of gap $g_s$ and their reflections of gap
$g_t$. Their grouped deficiencies give every subset's sum of the bounds
$f_\eta-g_v$. The leading/CORE8 stars give, separately,
\begin{equation}
 2X_v+2R+\varepsilon\ge278,
 \qquad 2(X_s+X_t)+4R+\varepsilon\ge556.           \label{eq:p4-stars}
\end{equation}
Each follows from actual 70-point intersections $139-R-X_v$; a leading
has only one full-resource capacity even when two core copies are used.

For a genuinely charged center $Z$ let $l_Z,h_Z$ be its lower and upper
hull clipping losses. At $m=1$, $I+\sigma+\delta\le58$ licenses the
uniform endpoint lemma at every integral center. At $m=2$,
$\delta\le25$, while an outside-hull center clips at least
$\min(r_s,r_t)+5\ge32$ points, so every charged center is normal.
Writing $a=1$ at $m=1$ and $a=0$ at $m=2$, both conclusions are
\begin{equation}
 Z\ge74-r_v-l_Z-a,
 \qquad Z\le S_v+r_o-73+h_Z+a.                    \label{eq:p4-END}
\end{equation}
Only a retained singleton or an actual sole pair component supplies
such a native loss. In particular no split witness is charged.

\subsection{Two reusable owner obstructions}\label{subsec:p4-owner}
These arguments concern actual charged resources and do not use the sign
of $\beta_s+\beta_t$. A beta resource below may be a sole beta-pair
component in A or an ordinary beta singleton in B.

\paragraph{Opposite owners.}
Suppose $\beta_s$ and $P_t:=C_{4t}$ are native capacity-78 resources,
CORE8 is retained, and their actual losses share the original budget.
The center identity
\[
 V_t=2d_s+\beta_s+P_t+2s_4+t_4+g_s
\]
uses only the literal centers. Apply lower endpoints to $\beta_s,P_t$
and an upper endpoint to $V_t$; with the distinct leading loss $\ell_s$,
the three selected endpoint losses $E$ satisfy $E+\ell_s\le\varepsilon$.
Substitute twice \eqref{eq:p4-lead}. Since $E+2\ell_s\le2\varepsilon$,
\begin{equation}
 S_s+3g_s+4X_s+2s_4+t_4
 \le1026+3a+r_s+\delta-2e.                       \label{eq:p4-opp-transport}
\end{equation}
The two ordinary gap-$g_s$ pairs give
$3g_s\ge3f_\eta-3\varepsilon/2$, and the stars give
$6X_s+X_t\ge973-7R-3\varepsilon$.
The positive source mass is
$S_s\ge T_r+\tau_s+M_s+3R_s+6$, $r=r_s$.
Use $R_s\ge r+1$, $M_s\ge R_s+4$, and
$\tau_s\ge r$ at $m=1$, $\tau_s\ge r-1$ at $m=2$.
Combining these rows, \eqref{eq:p4-L}, \eqref{eq:p4-Q62}, and
\eqref{eq:p4-radial} in \eqref{eq:p4-opp-transport} gives the exact cards
\[
\begin{array}{c|c|c}
 m& T_r+5r+\frac{29}{8}I+\frac\sigma4-
 \frac{15}{8}\delta+\frac{13}{2}e\ \le&
 T_r+5r+\frac{11}{2}I+\frac\sigma4+\frac{13}{2}e\ \le\\\hline
 1&343&1837/4\\
 2&2755/8&3685/8
\end{array}                                                    \tag{P4-opp}
\]
The last column uses only the weaker consequence $\delta\le62-I$.
For example before the radial substitution the $m=1$ row is
$T_r+5r+I+\sigma/4+803/4+2e\le21R/4+9\varepsilon/2+3$;
the $m=2$ row replaces $803/4$ by $199$ and omits the $+3$.
These display every endpoint allowance. Since $r\ge27$,
$T_r+5r\ge513$, whereas $1837/4=459.25$ and $3685/8=460.625$.
Thus opposite actual owners are impossible, already with $e=0$.

\paragraph{Same owner with active alpha.}
Suppose $\alpha_s,P_s=C_{4s},\beta_s$ are three distinct charged
capacity-78 resources and CORE8 is retained. No activity of $\alpha_t$
is assumed. The identity and strict outer order give
\[
 \alpha_s-P_s-\beta_s
 =\alpha_t+d_t+2(s_4+t_4)+G
 >d_t+2s_4+t_4+G.
\]
An upper endpoint for $\alpha_s$ and lower endpoints for $P_s,\beta_s$
use three one-unit allowances at $m=1$ and none at $m=2$.
Together with the distinct leading loss $\ell_t$, they yield
\begin{equation}
 S_t+g_t+2X_t+2s_4+t_4+G
 \le K+1095+3a-r_t-e.                            \label{eq:p4-same-transport}
\end{equation}
Three disjoint CORE8 pairs give $g_s+2g_t\ge3f_\eta-\varepsilon$;
the stars give $2X_s+3X_t\ge695-5R-3\varepsilon/2$.
Use $S_t\ge T_r+\tau_t+M_t+3R_t+6$, $r=r_t$,
$K\le M+R+21$, $R_t\ge r+1$, and \eqref{eq:p4-L}.
The combined row before the radial substitution is
\[
 T_r+5r+I+\sigma/4+\delta+\sigma_t+e
 \le \begin{cases}721/4+17R/4+5\varepsilon/2,&m=1,\\
                   178+17R/4+5\varepsilon/2,&m=2.
       \end{cases}
\]
It follows that
\begin{equation}
 T_r+5r+\frac{25}{8}I+\frac\sigma4+\frac58\delta
       +\sigma_t+\frac72e
 \le\begin{cases}618,&m=1,\\4943/8,&m=2.\end{cases}
                                                        \label{eq:p4-same-final}
\end{equation}
For integer $r$,
\[
 T_r+5r+\frac{25}{8}(r-34)(r-35)
 =\frac{5390+(r-29)(29(r-29)+1)}8\ge\frac{5390}{8}.
\]
At $m=1$ the remaining terms are nonnegative. At $m=2$,
$\sigma_t\ge-1$, while $\sigma,\delta,e\ge0$.
Thus the left side is at least $673.75$ or $672.75$, respectively,
contradicting \eqref{eq:p4-same-final}. The unique-1 owner was never
identified with either pair owner.

\subsection{The remaining retained-beta case: empty pair1}\label{subsec:p4-T5}
Suppose the exact missing units are $O_s,H_s,O_t,\mathrm{pair1}$,
the beta pair is actually sole on $s$, and every other singleton,
including CORE8, is retained. This is the only additional A1 predicate
needed below. We prove it for both first gaps, using precisely the
caps 58 and 25 established above.

First there is fixed credit $e=13$. Equation \eqref{eq:p4-lead} and
\eqref{eq:p4-radial} imply $d_v-g_v>R$ (already the weaker
$\delta\le62-I$ suffices). Empty $O_s,H_s,O_t$ are high by
\eqref{eq:p4-empty}, so the eleven fixed labels
\[
 C^s_{11},C^s_{12},C^s_{13},C^s_{14},C^s_{22},C^s_{23},C^s_{24},
 C^t_{11},C^t_{12},C^t_{13},C^t_{14}
\]
are high and belong to none of the 26 baseline groups.
The empty $\beta_t$ is low, since its higher retained $\alpha_t$ would
otherwise lose at least 74 points. Hence $V^t_5,V^t_6,V^t_7,V^t_8$ are low,
as are $F^v_3,F^v_4,F^v_5$ on both signs, where
$F_3=g-d,F_4=g-h-d,F_5=g-2h-d$.
Both empty $C_4$ centers are low; each forces
$Q_v=\{C^v_{55},C^v_{56},C^v_{57},C^v_{58}\}$ low.
Finally genuine $\beta_s,C_{5t}$ imply $d_t>d_s$, while $C_{5s}$ activity
gives $t_2<d_s$. Thus $V^t_3=h_t+g_t-d_t$ and $V^t_4=h_t-d_t$ are low.

For an exact label check, name the fifteen rows of the baseline group
list $A,\ldots,O$ in displayed order, with subscripts for reflected rows.
The ten labels
$Z=\{F^s_3,F^s_4,F^s_5,F^t_3,F^t_4,F^t_5,V^t_5,V^t_6,V^t_7,V^t_8\}$
meet exactly
\[
 B_s,D_s,K_s,B_t,D_t,K_t,O_s,F,M_t,L_t.
\]
The eight labels $Q_s\cup Q_t$ add precisely the seven groups
$G_t,N_s,H_t,C,G_s,N_t,H_s$; the two labels $V^t_3,V^t_4$ add only $I_t$,
since $H_s$ was already met. There are 20 distinct low labels meeting
18 groups, disjoint from the eleven high labels. The fixed-label formula
therefore gives $e=11+20-18=13$, and $\varepsilon=\delta-13$.

The CORE8 pairs and stars give the separately valid rows
\[
 3g_s+g_t\ge4f_\eta-\frac32\varepsilon,
 \qquad 6X_s+2X_t\ge1112-8R-3\varepsilon.
\]
For the first, add half the all-four-pair row to the two gap-$g_s$ pair
row; for the second, add the all-star row to twice the side-$s$ star row.
Use upper endpoints of $\alpha_s,H_t$ and lower endpoints of
$\beta_s,\alpha_t$. Then reuse the lower endpoint of $\beta_s$ and use
that of $C_{5t}$ in
$\beta_s+C_{5t}=d_t-d_s-2(s_4+t_4)-G$.
Twice the leading-$s$ row yields
\begin{equation}
 2M+6X_s+2X_t+3g_s+g_t\le943+6a+\delta-2e.       \label{eq:p4-T5transport}
\end{equation}
Indeed the weighted loss is
$2\ell_s+h_{\alpha_s}+h_{H_t}+l_{\alpha_t}+2l_{\beta_s}+l_{C_{5t}}
\le2\varepsilon$: it involves six distinct resources, each with coefficient
at most two. The six endpoint weakenings are exactly $6a$.
Using \eqref{eq:p4-L}, \eqref{eq:p4-Q62}, and
\eqref{eq:p4-radial} in \eqref{eq:p4-T5transport} gives
\begin{equation}
 39-\frac72\eta-6a+\frac92I+\sigma-\delta+\frac{13}{2}e\le0.
                                                        \label{eq:p4-T5final}
\end{equation}
With $e=13$ and $\delta\le58-I$, its left side is at least
$119/2+11I/2+\sigma$ at $m=1$ and $62+11I/2+\sigma$ at $m=2$.
Both are positive. Endpoint validity used $I+\sigma+\delta\le58$ at
$m=1$ and native loss at most 25 at $m=2$; no cap 54 or 21 is needed.

\subsection{Exhausting A1 and the strict branch B}\label{subsec:p4-A1B}
\paragraph{A1.}
Suppose both beta components were active. The triple prohibition after
\eqref{eq:p4-five} forces a hole in each of
$\{H_s,C_{6s}\}$ and $\{H_t,C_{6t}\}$.
These and the two $O$ holes exhaust $j=4$, so both $C_5$ singletons are
active. But
\[
 \beta_s+\beta_t+C_{5s}+C_{5t}
 =-2(s_3+s_4+t_3+t_4)<-4M
\]
contradicts the activity lower bound $2-2M$ (or its weaker $-2M$).
Thus beta is actually sole; rename its owner $s$.
If $H_s$ were active, $C_{6s}$ would be empty, and upward-closed cross
activity forces $C_{5s}$ empty. These are the two additional holes,
so $C_{7s}$ is active, contradicting the other triple prohibition.
Hence $O_s,H_s,O_t$ are the three known holes, with one remaining unit $Z$.

No $C_6$ or $C_7$ can be $Z$, by the cross chain. A missing $V_v$ has
higher retained $U_v$; missing $U_t$ has higher retained $H_t$; missing
$U_s$ has higher retained $\alpha_s$, because beta activity gives
$\alpha_s>U_s$. Each violates the 74-point loss rule.
CORE8 is therefore retained and
\[
 Z\in\{H_t,\alpha_s,\alpha_t,C_{5s},C_{5t},\mathrm{pair1}\}.
\]
Pair1 has at most one actual contributor: its two centers sum to
$-(s_2+s_3+s_4+t_2+t_3+t_4)<-3M$, whereas two activities would give
$>-M$. If $Z$ is pair1, \S\ref{subsec:p4-T5} applies.
Otherwise pair1 is sole. Opposite owners are excluded by
\S\ref{subsec:p4-owner}. With both owners $s$, the identity
\[
 C_{4s}+\beta_s+C_{5t}=-2s_3-2s_4-t_2-t_3-t_4
\]
precludes $C_{5t}$ activity; its right side is less than $-4M_s-3M_t$,
while three activities would give more than $-2M_s-M_t$.
Thus $Z=C_{5t}$, leaving $\alpha_s$ retained. The same-owner obstruction
applies. This exhausts all of A1, including initially split beta.

\paragraph{B.}
Here $\beta_v$ is a singleton and $C_5$ is one whole pair.
If $V_v$ were empty, higher $U_v,H_v$ would also have to be empty,
requiring three extra holes. If $U_v$ were empty, so would $H_v$;
then $\beta_v,\alpha_v$ are retained and beta activity gives
$\alpha_v>U_v$, again contradicting the high-empty rule.
Thus $U,V$ are retained.
If $C_{7v}$ were empty, $C_{6v}$ is empty too; all three
$\beta_o,H_o,C_{6o}$ would then be retained. If $C_{6v}$ alone were
empty, retained $C_{7v}$ forces a second hole among $\beta_v,H_v$,
leaving the same forbidden opposite triple. Hence CORE8 is retained.
The two additional holes are now exactly one each from
$\{H_s,\beta_s\}$ and $\{H_t,\beta_t\}$.
Both alphas and both pairs are retained.

An empty beta is low, since its higher retained alpha prohibits the high
arm. Both betas cannot be empty, for that would contradict
$\beta_s+\beta_t>0$. Pair1 is sole as above. Writing
$W=s_3+s_4+t_3+t_4>2M$, the identities
\[
 C_{5s}+C_{5t}=-(\beta_s+\beta_t)-2W,
 \qquad C_{4v}+C_{5o}=-(\beta_s+\beta_t)-2W-d_o
\]
show that the $C_5$ pair is sole and has the same actual owner as pair1.
Every active beta on the opposite side is excluded by the opposite-owner
obstruction: its singleton has exactly the native capacity and loss used
there. Thus there is exactly one active beta, on the common owner $s$;
the exact missing units are $O_s,O_t,H_s,\beta_t$.
Its retained $\alpha_s$, sole pair1 resource and beta singleton satisfy
the same-owner obstruction. The unselected $C_5$ pair keeps only one
capacity, and the empty beta singleton is already among the four missing
whole units. This proves the strict-B exclusion with its original roles.

\subsection{The beta-empty branch A0: three complete cases}\label{subsec:p4-A0}
The holes already include $O_s,O_t$ and the whole beta pair. Exactly one
further unit $Z$ is missing. Cross-chain activity prohibits $C_6,C_7$ as
$Z$; a missing $V_v$ has higher retained $U_v$, and missing $U_v$ has
higher retained $H_v$. Therefore CORE8 is retained and
\[
 Z\in\{H_s,H_t,\alpha_s,\alpha_t,C_{5s},C_{5t},\mathrm{pair1}\}.
                                                        \tag{A0-cover}
\]
If pair1 is retained it is actually sole. Normalize its owner as $s$;
then $C_{5s}$ is retained by upward-closed cross activity, leaving
$Z\in\{H_s,H_t,\alpha_s,\alpha_t,C_{5t}\}$.
On each side at least one of $H_v,\alpha_v$ is retained and higher than
the empty beta. Thus both beta centers are low by the 74-point rule.

There is genuine fixed credit $e=9$. The eight labels
$C^v_{11},C^v_{12},C^v_{13},C^v_{14}$ on both signs are high because both
$O$ centers are high. None belongs to a baseline group.
The two labels $V^s_6=\beta_s-t_2$, $V^t_6=\beta_t-s_2$ are low and
form exactly one baseline group, disjoint from the high labels.
Hence $e=8+2-1=9$, $\varepsilon=\delta-9\ge0$.
No hole allowance is added to this fixed-label count.

Use throughout the harmless weakened endpoints
\begin{equation}
 Z\ge73-r_v-l_Z,\qquad Z\le S_v+r_o-72+h_Z,
                                                        \label{eq:p4-A0end}
\end{equation}
valid at every $m=1$ center by END1, and weaker than the normal $m=2$
endpoints. Use also
\[
 f=\frac{237+R-\tau}{4}\le f_\eta,
 \quad M\ge R+39+I+\tau+\delta,
 \quad 2R+I+\delta\le207,                         \tag{A0-source}
\]
with the sharper radial 206 at $m=1$.

One joint-loss detail is useful. For a group of $n$ retained resources,
let $\Gamma$ be its actual union deficiency and let $u$ be one selected
member's lower or upper clipping loss. Pointwise multiplicity counting
proves
\[
 \Gamma\ge u+n|Q_{62}|-\left|\bigcup_{i=1}^n(Z_i+Q_{62})\right|.
                                                        \tag{joint-loss}
\]
Inside the target, multiplicity minus one pays overlaps. Outside,
full multiplicity also pays the one clipping indicator. Thus clipping
and overlap do not receive independent budgets. In particular a pair of
gap $g$ has deficiency at least $f-g$, a triple of center span $g+h$
has deficiency at least $f+62-g-h$, and two disjoint pairs of the same
gap with total deficiency $\Omega$ give $2g\ge2f-\Omega$.
The four CORE8 pairs give $2G\ge4f-\varepsilon$.
For two retained same-target centers $Y,Z$, pair deficiency $\Gamma$ and
lower clipping loss $l$ of the smaller center, \eqref{eq:p4-A0end} and
(joint-loss) give
\begin{equation}
 Y+Z\ge146-2r_v+f-\Gamma-l
      \ge146-2r_v+f-2\Gamma.                     \label{eq:p4-pairsum}
\end{equation}

\paragraph{Case 1: pair1 is empty.}
All $H,\alpha,C_5$ and CORE8 singletons are retained.
Put $a=s_4$, $b=t_4$, $a+b=M+X$, where $X=X_s+X_t$.
The exact identities are
\[
 H_v=2(a+b)+G+\alpha_s+\alpha_t+C_{5v}.
\]
Let $\Gamma_v$ be the deficiency of $\{\alpha_v,C_{5v}\}$ and $u_v$
the upper clipping loss of $H_v$. These disjoint groups satisfy
$\Gamma_s+\Gamma_t+u_s+u_t\le\varepsilon$.
The two pair-sum rows give total at least
$154+2f-2(\Gamma_s+\Gamma_t)$; the two lower alpha endpoints give
$\alpha_s+\alpha_t\ge77-l_{\alpha_s}-l_{\alpha_t}$.
As each alpha clipping is included in its pair deficiency, the two upper
$H$ endpoints yield
\[
 S\ge4(M+X)+2G+306+2f-3\varepsilon.
\]
Since $S\le2M+2R+1220-\delta$ by the source central-fringe bound,
\[
 2M+4X+2G+2f\le2R+914-\delta+3\varepsilon.
\]
Insert (A0-source), $4X\ge1112-8R-2\varepsilon$ from the stars, and
$2G\ge4f-\varepsilon$. This gives
\[
 276+2I+2\tau+6f+3\delta\le8R+6\varepsilon,
 \qquad \frac{1371}{2}+2I+\frac\tau2\le\frac{13}{2}R+3\delta.
\]
The radial 207 row bounds the last right side by
$2691/4-13I/4-\delta/4$. Equivalently,
\begin{equation}
 51+21I+2\tau+\delta\le0,                         \label{eq:p4-A0case1}
\end{equation}
which is impossible. All positions of all charged endpoints are covered.

\paragraph{Case 2: pair1 is retained and the opposite alpha is active.}
Let its actual owner be $s$, $P=C_{4s}$, $r=r_s$, $k=r_t=69-r$;
assume $\alpha_t$ is retained. No condition on $\alpha_s,C_{5t},H_s,H_t$
is required. The exact identity
\[
 U_s=P+\alpha_t+C_{6t}+2a+4b+2G
\]
uses four distinct retained native resources. Three lower endpoints and
the upper endpoint of $U_s$ give
\[
 S_s\ge222-2k+2a+4b+2G-E,\qquad E\le\varepsilon.
                                                        \tag{A0-cross-end}
\]
The constant 222 includes all four unit allowances of
\eqref{eq:p4-A0end}. Use
$S_t\ge T_k+\tau_t+4R_t+10$, $S\le2M+2R+1220-\delta$,
$R_t\ge k+1$, $M_t\ge k+5$,
$2X_s+4X_t\ge834-6R-2\varepsilon$ and
$2G\ge4f-\varepsilon$. Straight substitution gives
\begin{equation}
 T_k+4k+133-\tau_s\le7R+3\delta.                 \label{eq:p4-A0cross}
\end{equation}
At $m=1$ use radial 206 and $\tau_s\le127-I-\delta-\tau_t$:
\[
 T_k+4k+\frac92I+\frac32\delta+\tau_t\le715.
\]
But $\tau_t\ge k$, $\delta\ge9$, and for every integer $k$,
\[
 T_k+5k+\frac92(k-34)(k-35)
 =705+5(k-30)(k-31)\ge705.
\]
Thus $718.5\le715$, impossible. This exact $7/2$ margin includes all
four endpoint allowances, with no normal-center shortcut.
At $m=2$ instead use $R+\delta\le106-I$ in
\eqref{eq:p4-A0cross}; it gives
\[
 T_k+4k+8I+5\delta+\tau_t\le736.
\]
Here $\tau_t\ge k-1$, even when $t$ owns the unique 1. With $x=k-32$,
\[
 T_k+5k+8(k-34)(k-35)-1
 =735+\frac{17x^2-5x}{2}\ge735
\]
for every integer $x$. The left side is at least $735+45=780$, again
impossible. The physical low-1 owner remains unrestricted.

\paragraph{Case 3: pair1 is retained and the opposite alpha is empty.}
The one-extra-hole cover makes $Z=\alpha_t$; all $H_s,H_t,\alpha_s,C_5$
and CORE8 units are retained. Give no endpoint to the missing alpha.
With $r=r_s$, the center identities
\[
\begin{aligned}
 P&=d_s-d_t-2a-b-G-h_t,& C_{5t}&=d_t-2b-a-G-h_s,\\
 \alpha_s&=h_s-a,&V_s&=d_s-a-g_s,\qquad H_t=d_t+h_t-b
\end{aligned}
\]
allow upper endpoints at $V_s,H_t$, a lower endpoint at $P$, and twice
each the lower endpoints at $C_{5t},\alpha_s$. The exact result is
\[
 S\ge5a+4b+2g_s+3g_t+2h_t+302-r-E,
 \quad E=h_{V_s}+h_{H_t}+l_P+2l_{C_{5t}}+2l_{\alpha_s}
 \le2\varepsilon.                                \tag{A0-last-end}
\]
For the constant, the two upper endpoints first contribute 75; substitution
of $P$ and twice $C_{5t}$ gives $156+r$, and twice the alpha endpoint adds
$146-2r$. All seven weighted unit allowances are thus included.
The source upper bound implies
\begin{equation}
 2M+M_s+5X_s+4X_t+2g_s+3g_t+2h_t
 \le2R+918+r-\delta+2\varepsilon.                  \label{eq:p4-A0last}
\end{equation}
Use the following disjoint ordinary groups: the two gap-$g_s$ pairs,
combined deficiency $\Omega$; the triple $\{H_t,U_t,V_t\}$,
deficiency $\Psi$; and the pair $\{C_{6s},C_{7s}\}$,
deficiency $\Lambda$. They satisfy $\Omega+\Psi+\Lambda\le\varepsilon$,
so (joint-loss) gives
\[
 2g_s\ge2f-\Omega,\quad g_t+h_t\ge f+62-\Psi,
 \quad g_t\ge f-\Lambda,
\]
and therefore $2g_s+3g_t+2h_t\ge5f+124-2\varepsilon$.
Also
\[
 2M+M_s\ge2R+82+r+2I+2\tau+2\delta,
 \qquad5X_s+4X_t\ge1251-9R-\tfrac52\varepsilon.
\]
Substitute these in \eqref{eq:p4-A0last} and then use
$\varepsilon=\delta-9$, $4f=237+R-\tau$:
\[
 539+2I+2\tau+5f+3\delta\le9R+\tfrac{13}{2}\varepsilon,
 \qquad \frac{3575}{4}+2I+\frac34\tau
 \le\frac{31}{4}R+\frac72\delta.
\]
Radial 207 bounds the last right side by
$6417/8-31I/8-3\delta/8$. Equivalently,
\begin{equation}
 733+47I+6\tau+3\delta\le0,                        \label{eq:p4-A0case3}
\end{equation}
a contradiction.

Cases 1--3 are exhaustive: pair1 is empty, or it has an actual sole
owner and its opposite alpha is active or empty. Together with the A1
and B arguments, this proves that no original $P_4$ configuration exists.


\section{Compact source guide}\label{app:sources}
The main text is the consolidated argument. The following companion paths locate the retained foundations and the executable scope map, without inserting research chronology into the proof.
\begin{description}[style=nextline,leftmargin=0pt]
\raggedright
\item[The consolidated proof]
\nolinkurl{proof/report.tex} is the clean, standalone source of this PDF.
\item[The structural theorem]
\nolinkurl{foundations/proof/STRUCTURAL86.md} gives the full retained structural derivation and its input boundary. Its exact finite components are mapped in \nolinkurl{foundations/COVERAGE.json}.
\item[General additive and non-AP inputs]
\nolinkurl{foundations/source/structural/foundations/general_lemmas_source.tex} records the general statements and their references; \nolinkurl{nonap_defect.md} in the same folder gives the finite-to-infinite non-AP reduction. The published results include subset-sum stability and small-increment inverse results, finite-group sumset estimates with Kneser's theorem, simultaneous inverse and short-diameter results, and a central-interval theorem. They are invoked only at their stated hypotheses.
\item[High-hole foundations]
\nolinkurl{foundations/proof/dense_holes_direct_defects.tex} contains the bounded-density induction, profile inference and direct-defect reduction used by Theorem~\ref{thm:dense-input}. The correction in \nolinkurl{foundations/proof/endpoint_bridge_correction.txt} is incorporated in the common bridge proof here.
\item[Executable scope and source identity]
\texttt{README.md}, \nolinkurl{COVERAGE_MAP.json}, and \nolinkurl{SOURCE_MAP.json} explain the top-level checks. The corresponding files under \nolinkurl{foundations/} identify the exact finite foundation domains and source transformations. Sample full-run results are in \nolinkurl{sample_results/full/}.
\end{description}
Some retained foundation sources record an earlier global status. They are included for the specific lemmas identified above. The current theorem and the complete case assembly are those in this report.

\subsection{Published inputs}
The following references identify the published or publicly posted theorems used through the retained general sources. The applications and their exact hypotheses are spelled out in the bundled general-lemma supplement; the bibliography does not replace those applicability checks.
\begin{enumerate}
\item R. Carpenter, C. Defant and N. Kravitz, \emph{Sets with few subset sums}, arXiv:2605.05498, version 2 (24 August 2026). \href{https://arxiv.org/html/2605.05498v2}{Version 2 full text}. The stability input says that if $X$ is a positive $r$-element real set, $r\ge4$, $0\le M\le r-4$, and $\card{\FS(X)}\le1+T_r+M$, then $X=cB$ for a positive integer $r$-element set $B$ with total at most $T_r+M$.
\item E. Balandraud, B. Girard, S. Griffiths and Y. Ould Hamidoune, \emph{Subset sums in abelian groups}, European Journal of Combinatorics \textbf{34} (2013), 1269--1286. \href{https://doi.org/10.1016/j.ejc.2013.05.009}{Publisher record}. The retained finite-group inputs are Lemma 3.6 and Theorem 1.5, together with Kneser's theorem in the stated aperiodic applications.
\item Y. Stanchescu, \emph{On addition of two distinct sets of integers}, Acta Arithmetica \textbf{75} (1996), 191--194. \href{https://doi.org/10.4064/aa-75-2-191-194}{Primary journal record}. The supplement distinguishes its simultaneous inverse statement from the short-diameter Theorem 4 and Corollary 2. No extra gcd premise is inserted into the latter.
\item V. F. Lev, \emph{On consecutive subset sums}, Discrete Mathematics \textbf{187} (1998), 151--160. \href{https://doi.org/10.1016/S0012-365X(98)80006-X}{Publisher record}. The retained central-interval form is: if $B\subseteq[1,L]_{\Z}$, $\card B=s$ and $L\le3s/2-2$, then
\[
[2L-2s+1,\ \sum B-(2L-2s+1)]_{\Z}\subseteq\FS(B).
\]
\end{enumerate}

\end{document}
